% ============ 05. Documentación por archivo de código ============

\section{Documentación y crítica por archivo de código}

Esta sección documenta, para cada archivo de código del repositorio:
qué hace (descripción), cómo está construido (lógica), una crítica de
calidad (fortalezas, debilidades, mejoras), su ejecutabilidad real y
si produce los resultados que los investigadores buscaban.

\noindent Archivos documentados: \textbf{51}.

\subsection{Ingesta de datos}
\noindent Archivos: 3.

\subsubsection{src/ingesta/\_\_init\_\_.py}
\textbf{Rol:} ingesta · \textbf{Líneas:} 53

\paragraph{Descripción funcional}
Punto de entrada de carga de datos del paquete ingesta. Expone cargar\_consolidado() y cargar\_produccion(), con prioridad de rutas (CSV crudo Socrata -> XLSX local para investigadores; CSV crudo obligatorio para producción), normaliza nombres de columnas a MAYÚSCULAS y reporta conteos cargados.

\paragraph{Análisis de la lógica}
Define rutas absolutas ROOT/datos/raw y ROOT/datos/tarea\_join. cargar\_consolidado() prefiere datos/raw/investigadores\_consolidado.csv (Socrata) y cae a datos/tarea\_join/investigadores\_consolidado.xlsx; cargar\_produccion() exige datos/raw/produccion\_grupos.csv y lanza FileNotFoundError si falta, sugiriendo python -m src.ingesta.produccion. Es el contrato de entrada que usan todos los scripts (from ingesta import cargar\_consolidado, cargar\_produccion).

\paragraph{Fortalezas}
\begin{itemize}
  \item Mecanismo de fallback raw -> tarea\_join permite ejecutar análisis de investigadores sin datos crudos.
  \item Uppercase centralizado de columnas: garantiza el contrato MAYÚSCULAS que exigen analisis/*.py y Transformacion.py.
  \item Mensajes de error accionables que indican el comando de ingesta a ejecutar.
\end{itemize}

\paragraph{Debilidades}
\begin{itemize}
  \item cargar\_consolidado() no valida el shape esperado (77.237 x 30) ni los tipos: con low\_memory=True, id\_persona\_pr con ceros a la izquierda puede leerse de forma inconsistente (int vs str).
  \item cargar\_produccion() solo acepta el CSV crudo: no hay fallback ni lectura desde el modelo DuckDB; bloquea todo el análisis de producción en el clon.
  \item Depende de openpyxl para read\_excel, declarado solo en requirements.txt (no en pyproject.toml): una instalación vía poetry pura fallaría.
  \item El comando sugerido python -m src.ingesta.minciencias no pagina (hace una sola llamada con limit=100.000) y puede truncar silenciosamente los 77.237 registros si Socrata limita la respuesta por petición (contraste con ingesta/produccion.py que sí pagina).
\end{itemize}

\paragraph{Mejoras propuestas}
\begin{itemize}
  \item Añadir validación de shape y conteo contra datos/catalogo.yaml (tamanio\_registros).
  \item Especificar dtypes en read\_csv (id\_persona\_pr como str) para preservar ceros a la izquierda y garantizar el cruce con producción.
  \item Declarar openpyxl y pyyaml en pyproject.toml.
  \item Añadir fallback para cargar\_produccion desde parquet o desde el modelo DuckDB si el CSV crudo no está.
\end{itemize}

\paragraph{Ejecutabilidad}
cargar\_consolidado(): SÍ en el clon --- cae al XLSX datos/tarea\_join/investigadores\_consolidado.xlsx (existe, 8 MB) porque datos/raw/investigadores\_consolidado.csv NO existe (solo .gitkeep). cargar\_produccion(): NO --- datos/raw/produccion\_grupos.csv no existe (gitignored; requiere descarga de \textasciitilde{}1,2 GB). El módulo en sí es importable sin datos.

\paragraph{¿Produce lo esperado?}
No produce archivos; es capa de lectura. Sus salidas (DataFrames) alimentan transformar() y analisis.*. Cumple el contrato documentado en README (ingesta/ cargar\_consolidado + cargar\_produccion).

\begin{tcolorbox}[riesgo]
\textbf{Bugs / problemas detectados:}
\begin{itemize}
  \item src/ingesta/\_\_init\_\_.py:44-48: cargar\_produccion() lanza FileNotFoundError sin alternativa; en el clon (sin datos/raw) todos los scripts sprint5/sprint6 de producción mueren en el arranque.
\begin{quote}\footnotesize\ttfamily
\noindent\emph{Evidencia en src/ingesta/\_\_init\_\_.py (líneas 44-48):}
\end{quote}
\begin{lstlisting}
    if not _PROD.exists():
        raise FileNotFoundError(
            "Dataset de producción no encontrado. "
            "Ejecute: python -m src.ingesta.produccion"
        )
\end{lstlisting}

  \item src/ingesta/\_\_init\_\_.py:26: read\_csv(low\_memory=False) + read\_excel sin dtype explícito -> ID\_PERSONA\_PR puede inferirse como int en un lado y str en otro, rompiendo el cruce id\_persona\_pd <-> id\_persona\_pr.
\begin{quote}\footnotesize\ttfamily
\noindent\emph{Evidencia en src/ingesta/\_\_init\_\_.py (línea 26):}
\end{quote}
\begin{lstlisting}
        df = pd.read_csv(_RAW, low_memory=False)
\end{lstlisting}

  \item src/ingesta/\_\_init\_\_.py:18-36: no verifica duplicados por (ID\_PERSONA\_PR, ID\_CONVOCATORIA) pese a que el README define esa granularidad para el dataset.
\begin{quote}\footnotesize\ttfamily
\noindent\emph{Evidencia en src/ingesta/\_\_init\_\_.py (líneas 18-36):}
\end{quote}
\begin{lstlisting}
def cargar_consolidado() -> pd.DataFrame:
    """Carga el dataset consolidado de investigadores (77 237 registros, 30 cols).

    Prioridad: datos/raw/investigadores_consolidado.csv (Socrata)
               datos/tarea_join/investigadores_consolidado.xlsx (Excel local)
    """
[...]
\end{lstlisting}

\end{itemize}
\end{tcolorbox}

\paragraph{Recomendación de auditoría}
\textbf{mejorar --- diseño sólido de fallback, pero requiere validación de tipos/duplicados y dependencias declaradas (openpyxl/pyyaml) para que la instalación reproducible funcione.}

\medskip
\hrule

\subsubsection{src/ingesta/minciencias.py}
\textbf{Rol:} ingesta · \textbf{Líneas:} 70

\paragraph{Descripción funcional}
Descarga el dataset consolidado de investigadores reconocidos de MinCiencias desde datos.gov.co (API Socrata, sodapy; recurso bqtm-4y2h) y lo guarda en datos/raw/investigadores\_consolidado.csv. Actualiza el campo ultimo\_pull del catálogo YAML. CLI con --limit (muestra de pruebas) y --token (app token Socrata).

\paragraph{Análisis de la lógica}
descargar() crea el cliente Socrata, trae hasta limit (por defecto 100.000) registros y construye un DataFrame. guardar() persiste a CSV UTF-8 en datos/raw. actualizar\_catalogo() reescribe el YAML con la fecha de hoy. main() orquesta y, si no hay límite, actualiza el catálogo. Logging con formato estándar.

\paragraph{Fortalezas}
\begin{itemize}
  \item CLI limpia y documentada (--limit para pruebas evita descargas completas).
  \item Separación descargar/guardar/actualizar\_catalogo permite reutilización y tests.
  \item Maneja app token Socrata para evitar throttling.
\end{itemize}

\paragraph{Debilidades}
\begin{itemize}
  \item Sin paginación para datasets grandes (el límite por defecto 100.000 cubre el padrón de 77k, pero es frágil si el dataset crece).
  \item Escribe en datos/raw (gitignored): el artefacto no se versiona ni se valida.
  \item actualizar\_catalogo reescribe el YAML completo con yaml.dump, perdiendo comentarios y formato original del catálogo.
  \item Sin control de calidad post-descarga (conteo de filas, codificación UTF-8) --- el mojibake del CSV canónico no se detecta aquí.
\end{itemize}

\paragraph{Mejoras propuestas}
\begin{itemize}
  \item Añadir paginación genérica (offset) para datasets >100k.
  \item Validar el resultado (filas esperadas, codificación) tras la descarga.
  \item Actualizar el catálogo conservando comentarios o con un campo separado.
\end{itemize}

\paragraph{Ejecutabilidad}
Parcial: requiere red y API Socrata (no ejecutable offline); en la auditoría no se ejecutó porque los datos ya están en el clon (el XLSX versionado proviene de esta fuente). datos/raw está gitignored, por lo que el resultado no es visible en el repositorio.

\paragraph{¿Produce lo esperado?}
Sí, en su momento: los 77.237 registros (según README/catálogo) se descargaron de esta fuente; el XLSX consolidado versionado es el derivado. No verificable en clon por gitignore.

\begin{tcolorbox}[riesgo]
\textbf{Bugs / problemas detectados:}
\begin{itemize}
  \item yaml.dump reescribe el catálogo y puede alterar su estructura/orden (riesgo de formato).
\end{itemize}
\end{tcolorbox}

\paragraph{Recomendación de auditoría}
\textbf{mantener --- simple y funcional; mejorar la validación de codificación y el manejo del catálogo.}

\medskip
\hrule

\subsubsection{src/ingesta/produccion.py}
\textbf{Rol:} ingesta · \textbf{Líneas:} 106

\paragraph{Descripción funcional}
Descarga paginada del dataset de Producción de Grupos (Socrata 33dq-ab5a, 3.166.629 filas) a datos/raw/produccion\_grupos.csv, con modo --limit para pruebas y actualización de metadatos en datos/catalogo.yaml. Es la única ingesta del proyecto con paginación real.

\paragraph{Análisis de la lógica}
Usa sodapy.Socrata con timeout=300; descargar() pagina con \$offset/\$limit de 50.000 y order=':id' hasta que un chunk venga vacío o con menos de PAGE\_SIZE; con --limit hace una sola llamada. guardar() escribe CSV UTF-8; actualizar\_catalogo() actualiza ultimo\_pull y tamanio\_registros de la fuente 33dq-ab5a en catalogo.yaml; main() orquesta y omite el catálogo en modo prueba.

\paragraph{Fortalezas}
\begin{itemize}
  \item Paginación robusta con orden estable (:id) y condición de corte doble (chunk vacío o < PAGE\_SIZE).
  \item CLI argparse con modo limit para pruebas sin descargar 1,2 GB.
  \item Cierra el cliente Socrata explícitamente y loguea progreso por página.
  \item Actualiza el catálogo de metadatos automáticamente tras la descarga completa.
\end{itemize}

\paragraph{Debilidades}
\begin{itemize}
  \item Importa yaml en el módulo pero PyYAML NO está declarado en pyproject.toml ni requirements.txt (sodapy no lo trae) -> ImportError en un entorno limpio.
  \item Descarga todas las columnas sin proyección (select), incluyendo títulos largos (nme\_producto\_pd) -> payload y tiempo innecesarios en 3,2 M de filas.
  \item Sin reintentos/backoff ante errores de red ni validación del conteo final contra el catálogo (3.166.629).
  \item yaml.dump reescribe catalogo.yaml reformateándolo (indentación y comillas), generando diffs ruidosos en un archivo versionado con formato manual.
  \item if limit: con limit=0 ejecutaría la descarga completa (valor falsy no contemplado).
\end{itemize}

\paragraph{Mejoras propuestas}
\begin{itemize}
  \item Añadir select de columnas analíticas (id\_persona\_pd, id\_convocatoria, ano\_convo, nme\_tipo\_medicion\_pd, cod\_grupo\_gr, etc.).
  \item Declarar pyyaml en pyproject.toml y requirements.txt.
  \item Verificar len(df) contra catalogo['tamanio\_registros'] y reintentar páginas fallidas con backoff exponencial.
  \item Guardar en parquet con tipos explícitos (id\_persona\_pd como str de 10 caracteres).
\end{itemize}

\paragraph{Ejecutabilidad}
Ejecutable con red: python -m src.ingesta.produccion (completa) o --limit N. Requiere sodapy, pandas y PyYAML (este último no declarado en el proyecto). Escribe en datos/raw/ (gitignored; en el clon el directorio existe con .gitkeep y es escribible localmente, pero la salida NO está versionada ni presente en el clon).

\paragraph{¿Produce lo esperado?}
Produce datos/raw/produccion\_grupos.csv (gitignored, \textasciitilde{}1,2 GB), el input que exige cargar\_produccion(). No genera evidencias/ ni artifacts/ directamente; coincide con lo documentado en README (ingesta por API Socrata paginada, Sprint 5).

\begin{tcolorbox}[riesgo]
\textbf{Bugs / problemas detectados:}
\begin{itemize}
  \item src/ingesta/produccion.py:24: import yaml no cubierto por pyproject.toml -> falla en instalación reproducible (probablemente ejecutado ad-hoc con PyYAML instalado manualmente).
\begin{quote}\footnotesize\ttfamily
\noindent\emph{Evidencia en src/ingesta/produccion.py (línea 24):}
\end{quote}
\begin{lstlisting}
import yaml
\end{lstlisting}

  \item src/ingesta/produccion.py:47-53: en modo limit main() igualmente guarda el DataFrame como CSV 'completo' en SALIDA: si alguien corre --limit sin querer, sobrescribe el archivo crudo con una muestra.
\begin{quote}\footnotesize\ttfamily
\noindent\emph{Evidencia en src/ingesta/produccion.py (líneas 47-53):}
\end{quote}
\begin{lstlisting}
    if limit:
        registros = client.get(SOCRATA_ID, limit=limit)
        client.close()
        df = pd.DataFrame.from_records(registros)
        log.info("Descargados (modo limit=%d): %d filas, %d columnas",
                 limit, len(df), len(df.columns))
[...]
\end{lstlisting}

  \item src/ingesta/produccion.py:59-65: si la última página devuelve exactamente PAGE\_SIZE y la siguiente viene vacía, el loop hace una petición extra vacía (ineficiencia menor, no error).
\begin{quote}\footnotesize\ttfamily
\noindent\emph{Evidencia en src/ingesta/produccion.py (líneas 59-65):}
\end{quote}
\begin{lstlisting}
        chunk = client.get(SOCRATA_ID, limit=PAGE_SIZE, offset=offset, order=":id")
        if not chunk:
            break
        paginas.append(pd.DataFrame.from_records(chunk))
        offset += len(chunk)
        if len(chunk) < PAGE_SIZE:
[...]
\end{lstlisting}

\end{itemize}
\end{tcolorbox}

\paragraph{Recomendación de auditoría}
\textbf{mantener/mejorar --- única ingesta con paginación correcta del proyecto; corregir la dependencia PyYAML declarada y el sobrescrito en modo limit para producción segura.}

\medskip
\hrule

\subsection{Transformación y limpieza}
\noindent Archivos: 1.

\subsubsection{src/Transformacion.py}
\textbf{Rol:} transformacion · \textbf{Líneas:} 222

\paragraph{Descripción funcional}
Núcleo de limpieza y normalización del dataset consolidado de investigadores (77.237 registros). Aplica limpieza de texto (strip + mayúsculas), parseo robusto del año de convocatoria (fecha, serial de Excel o regex de 4 dígitos), estandarización de género y tratamiento de faltantes (mediana para numéricas, 'NO REPORTADO' para categóricas). Entrada: DataFrame crudo de ingesta.cargar\_consolidado(); salida: DataFrame transformado listo para análisis.

\paragraph{Análisis de la lógica}
Flujo en transformar(): limpiar\_texto -> parsear\_ano\_convocatoria -> estandarizar\_genero -> tratar\_faltantes. Usa constantes COLS\_NUMERICAS y COLS\_CATEGORICAS; funciones puras que copian el df. En \_\_main\_\_ inserta ROOT/src en sys.path y llama cargar\_consolidado() + transformar(). Se conecta con el resto del pipeline porque todos los scripts sprint* ejecutan df = transformar(cargar\_consolidado()) antes de llamar a analisis.*.

\paragraph{Fortalezas}
\begin{itemize}
  \item Funciones puras con docstrings claros y operaciones idempotentes (strip + upper).
  \item Parseo de año multiestrategia (fecha, serial Excel, regex) --- robusto frente a formatos heterogéneos de ANO\_CONVO.
  \item Centraliza la lógica de limpieza en un solo módulo consumido por todos los scripts.
\end{itemize}

\paragraph{Debilidades}
\begin{itemize}
  \item Imputa con mediana global sin crear columna indicadora de imputación: para un observatorio de calidad de datos, destruye la señal de 'missingness' (el propio proyecto reporta faltantes como hallazgo).
  \item No corrige la anomalía documentada de la convocatoria 833 (ANO\_CONVO=2019-12-06, año nominal 2018): ANO\_CONVO\_INT quedará en 2019 para esa convocatoria.
  \item Imputar mediana en columnas ordinales (ORDEN\_CLAS\_PR, NRO\_ORDEN\_FORM\_PR) puede producir valores no enteros que contaminan agrupaciones posteriores.
  \item Docstring dice 'Grupo 7' y 'generado por ingesta.py' mientras el repo es Grupo 8 y el módulo de carga es ingesta/\_\_init\_\_.py (desactualización cosmética).
\end{itemize}

\paragraph{Mejoras propuestas}
\begin{itemize}
  \item Añadir columna flag (ej. *\_imputado) por variable imputada y comparar análisis con/sin imputación.
  \item Aplicar corrección de año nominal para la conv. 833 según datos/catalogo.yaml (campo ano\_convo\_campo).
  \item Validar shape/nulos contra catálogo (77.237 x 30) y loguear filas descartadas.
  \item Escribir tests para parsear\_ano\_convocatoria con casos: '2013', serial 41739, '2019-12-06', NaN.
\end{itemize}

\paragraph{Ejecutabilidad}
Sí, ejecutable tal cual en el clon para el dataset de investigadores: \_\_main\_\_ carga vía ingesta.cargar\_consolidado(), que cae al fallback datos/tarea\_join/investigadores\_consolidado.xlsx (existe, 8 MB) porque datos/raw/ solo tiene .gitkeep (sin CSV crudo). Requiere pandas, numpy (declarados) y openpyxl (solo en requirements.txt, no en pyproject.toml). NO depende de datos/raw ni de .duckdb.

\paragraph{¿Produce lo esperado?}
No escribe archivos (solo devuelve DataFrame y print de shape); el volcado a CSV lo hacen los scripts sprint*. Su output alimenta todos los análisis del README; consistente con lo documentado.

\begin{tcolorbox}[riesgo]
\textbf{Bugs / problemas detectados:}
\begin{itemize}
  \item Transformacion.py:107-109: la regex r'\textbackslash\{\}b(20[0-9]\{2\})\textbackslash\{\}b' solo captura años 2000-2099; correcto para el dominio, pero un año fuera de rango quedaría como NaT sin aviso.
\begin{quote}\footnotesize\ttfamily
\noindent\emph{Evidencia en Transformacion.py (líneas 107-109):}
\end{quote}
\begin{lstlisting}
    # Intento 3: año de 4 dígitos embebido en el texto
    anio_regex = s.str.extract(r"\b(20[0-9]{2})\b")[0]
    fecha_regex = pd.to_datetime(anio_regex, format="%Y", errors="coerce")
\end{lstlisting}

  \item Transformacion.py:99-100: pd.to\_datetime(s, dayfirst=True) sobre un texto numérico tipo '640' (id de convocatoria) produciría NaT; depende del contenido real de ANO\_CONVO, que el catálogo documenta como año o timestamp.
\begin{quote}\footnotesize\ttfamily
\noindent\emph{Evidencia en Transformacion.py (líneas 99-100):}
\end{quote}
\begin{lstlisting}
    fecha = pd.to_datetime(s, dayfirst=True, errors="coerce")

\end{lstlisting}

  \item Transformacion.py:140-141: si una columna numérica fuera 100\% NaN, mediana=NaN y fillna(NaN) no corrige nada (caso límite no manejado).
\begin{quote}\footnotesize\ttfamily
\noindent\emph{Evidencia en Transformacion.py (líneas 140-141):}
\end{quote}
\begin{lstlisting}
            mediana = df[col].median()
            df[col] = df[col].fillna(mediana)
\end{lstlisting}

\end{itemize}
\end{tcolorbox}

\paragraph{Recomendación de auditoría}
\textbf{mejorar --- núcleo correcto y bien estructurado; debe mejorarse la trazabilidad de imputaciones y la corrección del año nominal 2018 antes de seguir usándolo en un informe de calidad de datos.}

\medskip
\hrule

\subsection{Modelo}
\noindent Archivos: 2.

\subsubsection{src/modelo/\_\_init\_\_.py}
\textbf{Rol:} modelado · \textbf{Líneas:} 2

\paragraph{Descripción funcional}
Marcador de paquete Python para src/modelo/ (módulo del modelo dimensional).

\paragraph{Análisis de la lógica}
Sin lógica: archivo de marcador.

\paragraph{Fortalezas}
\begin{itemize}
  \item Correcto por convención.
\end{itemize}

\paragraph{Ejecutabilidad}
N/A.

\paragraph{¿Produce lo esperado?}
N/A.

\paragraph{Recomendación de auditoría}
\textbf{mantener --- estándar de Python.}

\medskip
\hrule

\subsubsection{src/modelo/dimensional.py}
\textbf{Rol:} modelado · \textbf{Líneas:} 286

\paragraph{Descripción funcional}
Construye el modelo dimensional (Issue \#14) en DuckDB: 7 dimensiones (convocatoria, investigador, categoría, área, territorio, formación, institución) + fact\_clasificacion (investigador x convocatoria) + tabla puente N:N bridge\_hecho\_institucion. Entrada: DataFrame transformado; salida: datos/processed/observatorio.duckdb (gitignored).

\paragraph{Análisis de la lógica}
Construye cada dimensión con drop\_duplicates sobre su clave natural y SK solo para territorio/institucion; construir\_fact\_clasificacion() une la SK de territorio por columnas geográficas e inserta sk\_hecho; construir\_bridge\_institucion() descompone inst\_filia con un dict sk y un loop de filas. cargar\_en\_duckdb() crea el archivo (mkdir parents) y hace DROP + CREATE TABLE AS SELECT por tabla; resumen\_modelo() lista tablas y conteos. Importa normalizar\_categoria/ORDEN\_CATEGORIAS desde analisis.longitudinal (requiere ROOT/src en sys.path).

\paragraph{Fortalezas}
\begin{itemize}
  \item Esquema claro y documentado (asociación dimensión -> hecho, más bridge para la relación N:N de instituciones).
  \item SK para territorio e institución y ordenamiento determinista antes de drop\_duplicates en varias dimensiones.
  \item DB\_PATH centralizado y carga idempotente (DROP + CREATE).
  \item La tabla puente resuelve correctamente la granularidad N:N que una dimensión plana no podría representar.
\end{itemize}

\paragraph{Debilidades}
\begin{itemize}
  \item LA 'PRIMERA APARICIÓN' NO ES CRONOLÓGICA: construir\_dim\_investigador() ordena solo por ID\_PERSONA\_PR (no por ANO\_CONVO\_INT) y luego keep='first'; el snapshot de atributos no corresponde realmente a la primera convocatoria del investigador.
  \item El fact referencia claves naturales (id\_convocatoria, id\_persona\_pr, id\_clas\_pr, id\_area\_con\_pr, id\_niv\_formacion\_pr) sin SK: las dimensiones (salvo territorio/institucion) carecen de clave sustituta y no hay FKs/constraints --- es más un conjunto de tablas de referencia que un modelo estrella estricto.
  \item \_drop\_and\_create() es código muerto (definido y nunca llamado), sugiriendo una intención de DDL con constraints que no se implementó.
  \item construir\_bridge\_institucion() itera las \textasciitilde{}77.237 filas con iterrows: lento y sin vectorización.
  \item El archivo .duckdb es gitignored y NO existe en el clon: el modelo no es reproducible sin volver a ejecutar (requiere duckdb >= 1.0 y el XLSX).
\end{itemize}

\paragraph{Mejoras propuestas}
\begin{itemize}
  \item Ordenar por [ANO\_CONVO\_INT, ID\_PERSONA\_PR] en construir\_dim\_investigador() para que 'primera aparición' sea real.
  \item Añadir SK (surrogate) a todas las dimensiones y referenciarlas en el fact; definir PK/FK con ALTER TABLE o CREATE TABLE con constraints.
  \item Vectorizar el bridge (explode + map) en lugar de iterrows.
  \item Documentar el modelo como 'tablas de referencia + hecho' si no se van a imponer FKs, para no sobrevender el término 'estrella'.
\end{itemize}

\paragraph{Ejecutabilidad}
Sí para la parte de investigadores en el clon: requiere duckdb, pandas y el XLSX de tarea\_join (presente); escribe datos/processed/observatorio.duckdb (gitignored; en el clon solo hay .gitkeep, el archivo NO está). No requiere datos/raw. NOTA: para incorporar producción (sprint5\_duckdb.py) sí haría falta el CSV crudo.

\paragraph{¿Produce lo esperado?}
Sí (con datos): genera las tablas documentadas y evidencias/duckdb\_resumen\_modelo.txt existe como evidencia de ejecución del equipo; el README reporta el modelo estrella como completado (Issue \#14). El .duckdb final no está versionado por diseño (.gitignore: *.duckdb).

\begin{tcolorbox}[riesgo]
\textbf{Bugs / problemas detectados:}
\begin{itemize}
  \item src/modelo/dimensional.py:69-74: sort solo por ID\_PERSONA\_PR -> 'keep first' no es la primera aparición temporal; los atributos del snapshot (género, región) pueden pertenecer a una convocatoria posterior.
\begin{quote}\footnotesize\ttfamily
\noindent\emph{Evidencia en src/modelo/dimensional.py (líneas 69-74):}
\end{quote}
\begin{lstlisting}
    dim = (
        df[cols]
        .sort_values("ID_PERSONA_PR")
        .drop_duplicates(subset=["ID_PERSONA_PR"], keep="first")
        .reset_index(drop=True)
    )
\end{lstlisting}

  \item src/modelo/dimensional.py:41-44: \_drop\_and\_create nunca se usa (código muerto).
\begin{quote}\footnotesize\ttfamily
\noindent\emph{Evidencia en src/modelo/dimensional.py (líneas 41-44):}
\end{quote}
\begin{lstlisting}
def _drop_and_create(con: duckdb.DuckDBPyConnection, ddl: str) -> None:
    table = ddl.split("(")[0].split()[-1]
    con.execute(f"DROP TABLE IF EXISTS {table}")
    con.execute(ddl)
\end{lstlisting}

  \item src/modelo/dimensional.py:192-197: el merge por columnas geográficas con NaN (código DANE ausente) no empareja en pandas (NaN != NaN) -> fact con sk\_territorio NULL para filas sin geografía, sin contarlas ni advertirlo.
\begin{quote}\footnotesize\ttfamily
\noindent\emph{Evidencia en src/modelo/dimensional.py (líneas 192-197):}
\end{quote}
\begin{lstlisting}
    fact = fact.merge(
        dim_territorio[["sk_territorio", *geo_cols]],
        on=geo_cols,
        how="left",
    )
    fact = fact.drop(columns=geo_cols)
\end{lstlisting}

  \item src/modelo/dimensional.py:52-56: construir\_dim\_convocatoria conserva ANO\_CONVO\_INT = 2019 para la convocatoria id 20 (año nominal 2018), propagando la anomalía de datos documentada al modelo.
\begin{quote}\footnotesize\ttfamily
\noindent\emph{Evidencia en src/modelo/dimensional.py (líneas 52-56):}
\end{quote}
\begin{lstlisting}
    cols = ["ID_CONVOCATORIA", "NME_CONVOCATORIA", "ANO_CONVO_INT", "ANO_CONVO_FECHA"]
    dim = df[cols].drop_duplicates(subset=["ID_CONVOCATORIA"]).copy()
    dim = dim.sort_values("ANO_CONVO_INT").reset_index(drop=True)
    dim.columns = dim.columns.str.lower()
    return dim
\end{lstlisting}

\end{itemize}
\end{tcolorbox}

\paragraph{Recomendación de auditoría}
\textbf{mejorar --- modelo útil y correcto en granularidad, pero necesita SK/constraints para ser un 'modelo estrella' real y corregir el snapshot cronológico del investigador antes de citarlo como modelo dimensional canónico.}

\medskip
\hrule

\subsection{Análisis estadístico}
\noindent Archivos: 8.

\subsubsection{src/analisis/\_\_init\_\_.py}
\textbf{Rol:} orquestacion · \textbf{Líneas:} 80

\paragraph{Descripción funcional}
Fachada del paquete analisis: re-exporta las funciones públicas de genero, territorial, redes, diversidad, longitudinal y produccion, y declara \_\_all\_\_. Define la API estable que consumen scripts/sprint*\_*.py y streamlit\_app.py.

\paragraph{Análisis de la lógica}
Importa explícitamente cada función/constante de los 6 submódulos y las lista en \_\_all\_\_. No contiene lógica de negocio; su función es el contrato de importación (from analisis import pct\_femenino\_por\_area, hhi, construir\_grafo, comparar\_dane, ...). Cualquier script que importe analisis dispara la importación de networkx (vía redes) y de longitudinal en el arranque.

\paragraph{Fortalezas}
\begin{itemize}
  \item API pública explícita y bien documentada para todos los análisis del proyecto.
  \item \_\_all\_\_ completo y ordenado facilita autocompletado e inspección.
  \item Permite a scripts y dashboard importar sin conocer la estructura interna del paquete.
\end{itemize}

\paragraph{Debilidades}
\begin{itemize}
  \item Importación eager de los 6 submódulos: importar 'analisis' carga networkx y longitudinal (peso y acoplamiento); si una dependencia opcional falta, todo el paquete falla al importar.
  \item Duplicación de símbolos entre el bloque de imports y \_\_all\_\_ (mantenimiento manual: añadir una función exige tocar dos lugares).
  \item No hay validación de esquema en la frontera: las funciones asumen columnas MAYÚSCULAS y ANO\_CONVO\_INT sin verificar (contrato implícito con Transformacion.py).
\end{itemize}

\paragraph{Mejoras propuestas}
\begin{itemize}
  \item Convertir redes (y pyvis) en import perezoso o submódulo opcional para no penalizar al dashboard.
  \item Generar \_\_all\_\_ a partir de una sola lista de nombres para evitar duplicación.
  \item Añadir una función validate\_esquema(df) que verifique columnas requeridas antes de los análisis.
\end{itemize}

\paragraph{Ejecutabilidad}
Sí, importable en el clon si las dependencias están instaladas (pandas, numpy, networkx; pyvis solo dentro de generar\_html\_pyvis). No lee archivos: la carga la hace ingesta. El fallo de cargar\_produccion en el clon no afecta el import, pero sí a las funciones de .produccion al ejecutarlas.

\paragraph{¿Produce lo esperado?}
No produce outputs; habilita el resto del pipeline. Coincide con la estructura documentada en el README (src/analisis como 'lógica reutilizable').

\begin{tcolorbox}[riesgo]
\textbf{Bugs / problemas detectados:}
\begin{itemize}
  \item src/analisis/\_\_init\_\_.py:4-36: importa funciones que no usa directamente; si analisis/longitudinal.py fallara al importar (dependencias), derriba también genero/territorial, que son independientes (acoplamiento innecesario).
\begin{quote}\footnotesize\ttfamily
\noindent\emph{Evidencia en src/analisis/\_\_init\_\_.py (líneas 4-36):}
\end{quote}
\begin{lstlisting}
from .genero import pct_femenino_por_area, tabla_pivot_pct_femenino, brecha_genero, evolucion_pct_femenino
from .territorial import hhi, hhi_por_convocatoria, top_territorios, tabla_cuotas
from .redes import (
    construir_pares, construir_grafo, metricas_grafo,
    tabla_nodos, tabla_aristas, metricas_por_convocatoria,
    subgrafo_grado_minimo, generar_html_pyvis, normalizar_institucion,
[...]
\end{lstlisting}

\end{itemize}
\end{tcolorbox}

\paragraph{Recomendación de auditoría}
\textbf{mantener --- fachada correcta; refactorizar a imports perezosos o dividir el paquete para reducir el acoplamiento del dashboard.}

\medskip
\hrule

\subsubsection{src/analisis/calidad.py}
\textbf{Rol:} analisis · \textbf{Líneas:} 105

\paragraph{Descripción funcional}
Validación documentada de la calidad de los datos del padrón, con foco en atípicos de edad (EDAD\_ANOS\_PR). Documenta explícitamente la decisión metodológica (eliminación frente a imputación por mediana del grupo) para que cualquier auditor pueda reproducirla y discutirla. Expone detectar\_atipicos\_edad, resumen\_tratamiento, filtrar (método elegido) e imputar\_mediana (alternativa comparativa).

\paragraph{Análisis de la lógica}
Umbral UMBRAL\_EDAD\_MAX=100. detectar\_atipicos\_edad devuelve los registros con edad>100 (mínimas columnas de identificación). resumen\_tratamiento produce el reporte antes/después con justificación. filtrar aplica eliminación (decisión por defecto). imputar\_mediana reemplaza los atípicos por la mediana del grupo (gran área, clasificación, género) como alternativa reproducible. El docstring documenta el razonamiento con la cifra de su momento ('diez registros' sobre el total entonces disponible, 0.013\%); el conteo vigente en el README y en las evidencias versionadas (git HEAD) es de 22 casos sobre 77.237 (0.028\%), sesgo despreciable, eliminación preferida por trazabilidad.

\paragraph{Fortalezas}
\begin{itemize}
  \item Decisión metodológica documentada y justificada en el propio código (transparencia auditable).
  \item Ofrece comparación explícita entre eliminación e imputación (análisis de sensibilidad).
  \item Funciones puras y pequeñas, sin efectos laterales.
  \item El filtro se aplica solo donde el análisis depende de la edad.
\end{itemize}

\paragraph{Debilidades}
\begin{itemize}
  \item Docstring dice 'diez registros' (10) mientras el conteo real documentado en README y evidencias es 22 --- inconsistencia documental verificada.
  \item Umbral fijo (100) sin parametrización.
  \item imputar\_mediana usa apply fila a fila (lento con 77k filas).
  \item El resto de anomalías (NO DISPONIBLE 5.2\%, NO REPORTADO 2.7\%, conv. 833 etiquetada 2019) se documenta en el docstring pero no se instrumenta como funciones.
\end{itemize}

\paragraph{Mejoras propuestas}
\begin{itemize}
  \item Corregir el docstring a 22 casos y parametrizar el umbral.
  \item Vectorizar imputar\_mediana (groupby transform en vez de apply).
  \item Instrumentar también las otras dimensiones de calidad (completitud, consistencia, unicidad) como funciones.
\end{itemize}

\paragraph{Ejecutabilidad}
Sí, verificada empíricamente: sprint6\_validacion\_calidad.py (que lo usa) ejecutado en clon fresco con EXIT=0 en 61.6 s. Requiere el XLSX consolidado (presente).

\paragraph{¿Produce lo esperado?}
Sí: evidencias/calidad\_atipicos\_edad.csv, calidad\_comparacion.csv y calidad\_resumen.json existen. Matiz verificable: la versión versionada en git (HEAD) reporta 22 atípicos sobre 77.237 (0.0285\%); el working tree regenerado por la propia ejecución de la auditoría sobre el XLSX parcial (3 convocatorias, 50.891 filas) reporta 19 (0.0373\%). El tratamiento documentado (eliminación) es el mismo en ambos casos.

\begin{tcolorbox}[riesgo]
\textbf{Bugs / problemas detectados:}
\begin{itemize}
  \item Docstring (línea \textasciitilde{}10): 'diez registros' vs 22 reales --- inconsistencia que puede confundir al lector del informe.
\end{itemize}
\end{tcolorbox}

\paragraph{Recomendación de auditoría}
\textbf{mejorar --- lógica correcta y bien documentada; corregir la inconsistencia documental y ampliar la instrumentación de calidad.}

\medskip
\hrule

\subsubsection{src/analisis/diversidad.py}
\textbf{Rol:} analisis · \textbf{Líneas:} 162

\paragraph{Descripción funcional}
Análisis de variables de diversidad (Issue \#20): cobertura temporal de ID\_VICTIMA\_CONFLICTO, TXT\_GRUPO\_ETNICO y TXT\_POBLACION\_DISCA, distribución por categoría, comparación de subrepresentación frente a DANE 2018/RUV 2021, cruce interseccional género x etnia y categorías por minoría. Entrada: DataFrame transformado (mayúsculas aplicadas por src/Transformacion.py); salida: DataFrames exportados a evidencias/diversidad\_*.csv.

\paragraph{Análisis de la lógica}
Define DANE\_REFERENCIA (constantes poblacionales) y VALORES\_NULOS por variable. cobertura\_por\_convocatoria() calcula \% NO REGISTRA/NO DISPONIBLE y cobertura por año. comparar\_dane() (diseñada para df\_2021) mapea nombres MinCiencias -> DANE con regex y calcula la razón de subrepresentación. interseccional\_genero\_etnia() y categoria\_por\_minoria() cruzan con género y categoría académica. El script sprint4\_diversidad.py filtra df\_2021 antes de llamar a comparar\_dane.

\paragraph{Fortalezas}
\begin{itemize}
  \item Hallazgo central documentado en el propio módulo: las variables solo se capturaron desde 2021 --- honestidad metodológica.
  \item Razón de subrepresentación (DANE/MinCiencias) correcta y con manejo de división por cero (None).
  \item Constantes de referencia (DANE 2018 / RUV 2021) centralizadas y comentadas.
\end{itemize}

\paragraph{Debilidades}
\begin{itemize}
  \item Contrato frágil: comparar\_dane, interseccional\_genero\_etnia y categoria\_por\_minoria asumen que el caller ya filtró la convocatoria 2021; si se pasa el df completo, los resultados mezclan años sin advertencia.
  \item ACOPLAMIENTO ROTO CON Transformacion.py: éste convierte NaN categórico en 'NO REPORTADO', pero diversidad espera 'NO REGISTRA'/'NO DISPONIBLE'; los faltantes reales etiquetados 'NO REPORTADO' no se cuentan como nulos -> cobertura inflada.
  \item Bug metodológico en comparar\_dane(): el comentario (línea 78) dice excluir 'NINGUN GRUPO ETNICO' del denominador, pero el código (línea 79) solo excluye 'NO DISPONIBLE'; 'NINGUN GRUPO ETNICO' queda en el denominador y diluye la subrepresentación de las minorías (verificado: comentario vs código).
  \item Comparación 'SÍ' con tilde (línea 121): funciona con los datos actuales porque transformar() pone en mayúsculas 'Sí' -> 'SÍ', pero es frágil ante variantes sin tilde ('Si'/'SI') --- el catálogo documenta 'Si'.
  \item interseccional\_genero\_etnia() puede lanzar KeyError 'FEMENINO' solo si NINGÚN grupo étnico tiene mujeres (con unstack(fill\_value=0) la columna existe si al menos un grupo tiene mujeres); caso límite, crash en runtime.
\end{itemize}

\paragraph{Mejoras propuestas}
\begin{itemize}
  \item Que las funciones filtren internamente por ANO\_CONVO\_INT == 2021 (parametrizado) para hacer la API autónoma.
  \item Alinear los valores nulos entre Transformacion y diversidad (constante compartida) o contar también 'NO REPORTADO' como faltante.
  \item Excluir de verdad 'NINGUN GRUPO ETNICO' del denominador (alinear comentario y código) y documentar el impacto en las razones de subrepresentación.
  \item Normalizar 'SÍ'/'SI'/'Si' con strip + upper antes de comparar (o eliminar acentos).
  \item Usar .reindex(columns=['FEMENINO', 'MASCULINO'], fill\_value=0) tras el unstack en interseccional\_genero\_etnia.
\end{itemize}

\paragraph{Ejecutabilidad}
Sí en el clon: solo DataFrame transformado (XLSX presente); las funciones de 2021 requieren que el caller filtre (sprint4\_diversidad.py lo hace). Sin datos/raw ni .duckdb.

\paragraph{¿Produce lo esperado?}
Sí: evidencias/diversidad\_cobertura\_por\_convocatoria.csv, diversidad\_comparacion\_dane.csv, diversidad\_distribucion\_*.csv, diversidad\_interseccional\_genero\_etnia.csv y diversidad\_categorias\_por\_etnia.csv existen, junto con artifacts/sprint4\_diversidad/fig01-03.png --- consistente con el README (Issue \#20, hallazgos 4-6).

\begin{tcolorbox}[riesgo]
\textbf{Bugs / problemas detectados:}
\begin{itemize}
  \item src/analisis/diversidad.py:78-79: contradicción comentario-código sobre 'NINGUN GRUPO ETNICO' en el denominador de comparar\_dane() --- diluye la subrepresentación de minorías (verificado por lectura directa).
\begin{quote}\footnotesize\ttfamily
\noindent\emph{Evidencia en src/analisis/diversidad.py (líneas 78-79):}
\end{quote}
\begin{lstlisting}
    # Etnia: excluir "NINGUN GRUPO ETNICO" del denominador para ver minorias
    etnia = df_2021[df_2021["TXT_GRUPO_ETNICO"] != "NO DISPONIBLE"]
\end{lstlisting}

  \item src/analisis/diversidad.py:145: d['FEMENINO'] puede lanzar KeyError solo si NINGÚN grupo étnico tiene mujeres (con unstack(fill\_value=0) la columna no se crea en ese caso) --- crash en caso límite (reproducido).
\begin{quote}\footnotesize\ttfamily
\noindent\emph{Evidencia en src/analisis/diversidad.py (línea 145):}
\end{quote}
\begin{lstlisting}
        .assign(pct_femenino=lambda d: (d["FEMENINO"] / d["total"] * 100).round(1))
\end{lstlisting}

  \item src/analisis/diversidad.py:44-46: los NaN convertidos a 'NO REPORTADO' por Transformacion.py no se cuentan como nulos -> pct\_no\_registra subestimado y cobertura sobreestimada para registros con faltante real.
\begin{quote}\footnotesize\ttfamily
\noindent\emph{Evidencia en src/analisis/diversidad.py (líneas 44-46):}
\end{quote}
\begin{lstlisting}
            pct_nulo = (grupo[col] == val_nulo).mean() * 100
            fila[f"{col}_pct_no_registra"] = round(pct_nulo, 2)
            fila[f"{col}_pct_cobertura"] = round(100 - pct_nulo, 2)
\end{lstlisting}

  \item src/analisis/diversidad.py:121: comparación 'SÍ' sensible a acentos; no es un bug activo (los datos traen 'Sí' y transformar() lo pone en mayúsculas), pero fallaría si la fuente usara 'Si' sin tilde.
\begin{quote}\footnotesize\ttfamily
\noindent\emph{Evidencia en src/analisis/diversidad.py (línea 121):}
\end{quote}
\begin{lstlisting}
    n_vict = (vict["ID_VICTIMA_CONFLICTO"] == "SÍ").sum()
\end{lstlisting}

\end{itemize}
\end{tcolorbox}

\paragraph{Recomendación de auditoría}
\textbf{mejorar --- hallazgos valiosos y bien enmarcados; prioridad alta: corregir el denominador de comparar\_dane(), alinear valores nulos con Transformacion.py y robustecer comparaciones de cadenas.}

\medskip
\hrule

\subsubsection{src/analisis/genero.py}
\textbf{Rol:} analisis · \textbf{Líneas:} 85

\paragraph{Descripción funcional}
Análisis de género por área OCDE (Issue \#22): porcentaje femenino, tablas pivote y brecha de género por convocatoria a 3 niveles de área. Entrada: DataFrame transformado con NME\_GENERO\_PR estandarizado y ANO\_CONVO\_INT; salida: DataFrames tidy/pivot listos para exportar a evidencias/genero\_*.csv y graficar.

\paragraph{Análisis de la lógica}
pct\_femenino\_por\_area() filtra género en \{FEMENINO, MASCULINO\} (excluye NO REPORTADO del denominador, decisión documentada), agrupa por (año, área), hace value\_counts().unstack() y calcula n\_femenino, n\_masculino, n\_total y pct. tabla\_pivot\_pct\_femenino() pivotea área x año con promedio y ordena; brecha\_genero() deriva pct\_masculino = 1 - pct\_femenino y la brecha; evolucion\_pct\_femenino() reordena. El script sprint2\_genero\_ocde.py llama a estas funciones y exporta los CSVs.

\paragraph{Fortalezas}
\begin{itemize}
  \item Decisión metodológica explícita y correcta: excluir NO REPORTADO del denominador evita sesgos de imputación.
  \item Funciones puras, parametrizables (col\_area, col\_year) y reutilizables en scripts y dashboard.
  \item Documenta el contrato de entrada (género estandarizado, ANO\_CONVO\_INT disponible).
\end{itemize}

\paragraph{Debilidades}
\begin{itemize}
  \item evolucion\_pct\_femenino() es redundante (mismo DataFrame ordenado); el propio docstring lo admite.
  \item No aplica dropna sobre col\_area: un área con NaN se pierde silenciosamente del groupby (aceptable pero implícito).
  \item pct\_femenino = n\_fem/n\_total puede dividir por cero si un grupo tiene 0 registros reportados -> NaN propagado a pivote y brecha sin aviso.
  \item brecha usa pct\_masculino = 1 - pct\_femenino, válido solo dentro del universo 'reportado'; la columna n\_masculino calculada no se aprovecha.
\end{itemize}

\paragraph{Mejoras propuestas}
\begin{itemize}
  \item Eliminar evolucion\_pct\_femenino o fusionarla con un parámetro sort.
  \item Añadir manejo explícito de grupos vacíos (n\_total=0 -> pct=NaN con advertencia).
  \item Escribir tests con df pequeño (área 100\% hombres, área con solo NO REPORTADO).
\end{itemize}

\paragraph{Ejecutabilidad}
Sí en el clon: solo requiere el DataFrame transformado, que se obtiene de cargar\_consolidado() con fallback al XLSX presente. Sin dependencias de datos/raw ni .duckdb.

\paragraph{¿Produce lo esperado?}
Sí: evidencias/genero\_pct\_femenino\_por\_area.csv, genero\_pct\_femenino\_por\_gran\_area.csv, genero\_brecha\_por\_gran\_area.csv y genero\_tabla\_pivot\_gran\_area.csv existen en el clon, junto con artifacts/sprint2\_genero\_ocde/fig01-04.png --- consistente con el README (Issue \#22).

\begin{tcolorbox}[riesgo]
\textbf{Bugs / problemas detectados:}
\begin{itemize}
  \item src/analisis/genero.py:32-33: conteos.get('FEMENINO', 0) funciona porque DataFrame.get acepta default escalar, pero si un grupo queda sin filas reportadas, n\_total=0 y pct\_femenino=NaN (división por cero) sin advertencia --- caso límite real si una convocatoria tuviera 100\% NO REPORTADO.
\begin{quote}\footnotesize\ttfamily
\noindent\emph{Evidencia en src/analisis/genero.py (líneas 32-33):}
\end{quote}
\begin{lstlisting}
    resultado["n_femenino"] = conteos.get("FEMENINO", 0)
    resultado["n_masculino"] = conteos.get("MASCULINO", 0)
\end{lstlisting}

  \item src/analisis/genero.py:67: pct\_masculino = 1 - pct\_femenino asume universo binario; correcto tras el filtro, pero la columna n\_masculino ya existe y no se usa (inconsistencia menor).
\begin{quote}\footnotesize\ttfamily
\noindent\emph{Evidencia en src/analisis/genero.py (línea 67):}
\end{quote}
\begin{lstlisting}
    base["pct_masculino"] = 1 - base["pct_femenino"]
\end{lstlisting}

\end{itemize}
\end{tcolorbox}

\paragraph{Recomendación de auditoría}
\textbf{mantener --- análisis correcto y bien documentado; mejoras menores de robustez (grupos vacíos) y limpieza de redundancia.}

\medskip
\hrule

\subsubsection{src/analisis/longitudinal.py}
\textbf{Rol:} analisis · \textbf{Líneas:} 207

\paragraph{Descripción funcional}
Núcleo del análisis longitudinal (Issue \#11): construye el panel investigador$\times$convocatoria a partir de ID\_PERSONA\_PR, compara trayectorias entre periodos consecutivos (sube/baja/se mantiene/desaparece), calcula matrices de transición de conteos y probabilidades, y tasas de retención con entrada de nuevos investigadores. Cubre las 6 convocatorias históricas (2013-2021).

\paragraph{Análisis de la lógica}
ORDEN\_CATEGORIAS define la escala ordinal (Junior=1 \dots{} Emérito=4) y normalizar\_categoria mapea el texto libre a las 4 etiquetas canónicas. construir\_panel deduplica (ID, año) conservando la categoría más alta. comparar\_periodo hace left join desde t0 hacia t1 (los que desaparecen quedan representados) y clasifica el resultado. matriz\_transicion calcula conteos (crosstab) y probabilidades normalizadas por fila, con opción de incluir 'Desaparece' como estado final. tasa\_retencion\_por\_periodo computa retención (presentes en t0 que reaparecen) y nuevos en t1.

\paragraph{Fortalezas}
\begin{itemize}
  \item Implementación estadística correcta: crosstab normalizado por fila = matriz de Markov empírica; left join conserva la atrición.
  \item Escala ordinal explícita para clasificar sube/baja/se mantiene.
  \item Estado 'Desaparece' tratado como opción configurable.
  \item Docstrings con interpretación metodológica (caveat eméritos vitalicios documentado en README/sprint6).
\end{itemize}

\paragraph{Debilidades}
\begin{itemize}
  \item No se testea la propiedad markoviana ni la homogeneidad temporal; con intervalos irregulares (2013, 2014, 2015, 2017, 2019, 2021) las tasas no son directamente comparables entre periodos.
  \item comparar\_periodo usa apply con closure fila a fila (lento con 77k filas).
  \item La 'desaparición' incluye a los eméritos vitalicios (diseño del sistema), que sesgan la interpretación si no se aísla (documentado en sprint6, no en el código).
  \item No hay análisis de supervivencia (Kaplan-Meier/Cox) para el tiempo hasta la salida.
\end{itemize}

\paragraph{Mejoras propuestas}
\begin{itemize}
  \item Añadir test de markovianidad y matrices condicionadas a historia.
  \item Vectorizar la clasificación del tracking (np.select en vez de apply).
  \item Complementar con modelos de supervivencia y reportar IC bootstrap de las tasas.
\end{itemize}

\paragraph{Ejecutabilidad}
Sí, verificada empíricamente: usado por sprint2\_panel\_longitudinal.py (EXIT=0, 108 s) y sprint2\_matrices\_transicion.py (EXIT=0, 74 s) en clon fresco. Requiere XLSX consolidado (presente).

\paragraph{¿Produce lo esperado?}
Sí: matrices de transición (evidencias/matriz\_*.csv), paneles y tasas versionados coinciden con la lógica; los regenerados con el clon (3 convocatorias) cubren solo 2 periodos --- limitación de datos, no de código.

\begin{tcolorbox}[riesgo]
\textbf{Bugs / problemas detectados:}
\begin{itemize}
  \item Ninguno crítico; riesgo metodológico: interpretar las transiciones como cadena de Markov sin validar la propiedad (caveat recomendado).
\end{itemize}
\end{tcolorbox}

\paragraph{Recomendación de auditoría}
\textbf{mantener con mejoras --- lógica sólida para el alcance descriptivo; la ampliación inferencial (markovianidad, supervivencia) es la recomendación de mayor valor.}

\medskip
\hrule

\subsubsection{src/analisis/produccion.py}
\textbf{Rol:} analisis · \textbf{Líneas:} 340

\paragraph{Descripción funcional}
Análisis cruzado producción <-> investigadores (Sprint 5): cobertura de reconocimiento, productividad por categoría/género/área/territorio, mix de tipologías y reconocidos sin producción. Entrada: prod (3.166.629 filas, Socrata 33dq-ab5a) e inv (consolidado transformado), cruzados por ID\_PERSONA\_PD <-> ID\_PERSONA\_PR (+ ID\_CONVOCATORIA); salida: DataFrames exportados a evidencias/produccion\_*.csv y consumidos por sprint5\_produccion.py.

\paragraph{Análisis de la lógica}
normalizar\_produccion() pone columnas en MAYÚSCULAS, convierte IDs a Int64 y extrae ANO\_CONVO\_INT. cobertura\_reconocidos() marca es\_reconocido por isin global y agrega por año; cobertura\_autores\_unicos() devuelve un dict de solapamiento. productividad\_por\_categoria()/por\_genero\_area() agregan n\_productos por (persona, convocatoria) y hacen left merge con inv. productividad\_territorial() suma por departamento con per cápita y ratio de concentración. mix\_tipologias() y tipologias\_por\_categoria() agregan por NME\_TIPO\_MEDICION\_PD. reconocidos\_sin\_produccion() cruza conjuntos de autores por ID\_CONVOCATORIA.

\paragraph{Fortalezas}
\begin{itemize}
  \item Responde de forma completa las 6 preguntas críticas del Sprint 5 con funciones bien nombradas y docstrings que explican la interpretación de cada métrica.
  \item Manejo correcto de tipos para el cruce (to\_numeric -> Int64 en ambos lados).
  \item Coberturas y brechas con denominadores explícitos (pct\_con\_al\_menos\_un\_producto, razon\_f\_m, ratio\_concentracion).
  \item Evidencias producidas y versionadas: 7 CSVs + resumen JSON en evidencias/.
\end{itemize}

\paragraph{Debilidades}
\begin{itemize}
  \item Requiere datos/raw/produccion\_grupos.csv (\textasciitilde{}1,2 GB, gitignored): en el clon ninguna función de este módulo es ejecutable hasta descargar el archivo.
  \item normalizar\_produccion() solo usa pd.to\_datetime para ANO\_CONVO: si el CSV exporta años como enteros (2013), to\_datetime los interpreta como nanosegundos -> año 1970; no replica la estrategia multi-formato de Transformacion.py (riesgo latente según el formato real del export).
  \item reconocidos\_sin\_produccion() hace groupby(...).apply(set) sobre 3,2 M filas: lento y con alto uso de memoria.
  \item Los merges asumen granularidad única (persona, convocatoria) en inv; duplicados inflarían n\_productos (sin verificación previa).
  \item brecha\_productividad\_genero() lanza KeyError si un área no tiene mujeres (columna ausente tras el pivot).
\end{itemize}

\paragraph{Mejoras propuestas}
\begin{itemize}
  \item Replicar el parseo multiestrategia de año de Transformacion.py en normalizar\_produccion() o unificarlo en un helper compartido.
  \item Vectorizar reconocidos\_sin\_produccion() con merge sobre autores únicos por convocatoria.
  \item Verificar unicidad de (ID\_PERSONA\_PR, ID\_CONVOCATORIA) en inv antes de mergear.
  \item Reindexar columnas en brecha\_productividad\_genero para áreas sin un género.
\end{itemize}

\paragraph{Ejecutabilidad}
NO en el clon: cargar\_produccion() falla porque datos/raw/produccion\_grupos.csv no existe (gitignored; requiere python -m src.ingesta.produccion, \textasciitilde{}1,2 GB y red). Las funciones son puras y ejecutables si se les pasan DataFrames, pero el flujo documentado exige el CSV crudo. No depende del .duckdb.

\paragraph{¿Produce lo esperado?}
Sí (con datos): evidencias/produccion\_cobertura\_por\_convocatoria.csv, produccion\_productividad\_por\_categoria.csv, produccion\_brecha\_genero.csv, produccion\_concentracion\_territorial.csv, produccion\_mix\_tipologias.csv, produccion\_tipologias\_por\_categoria.csv, produccion\_reconocidos\_sin\_produccion.csv + produccion\_resumen\_cobertura.json + 6 figuras en artifacts/sprint5\_produccion/ --- todos existen en el clon y coinciden con el README (Sprint 5, hallazgos 8-10).

\begin{tcolorbox}[riesgo]
\textbf{Bugs / problemas detectados:}
\begin{itemize}
  \item src/analisis/produccion.py:43: pd.to\_datetime(ANO\_CONVO) sin manejo de ints/seriales -> año 1970 si el campo es numérico (inconsistente con Transformacion.py:99-116).
\begin{quote}\footnotesize\ttfamily
\noindent\emph{Evidencia en src/analisis/produccion.py (línea 43):}
\end{quote}
\begin{lstlisting}
        df["ANO_CONVO_INT"] = pd.to_datetime(df["ANO_CONVO"], errors="coerce").dt.year.astype("Int64")
\end{lstlisting}

\begin{quote}\footnotesize\ttfamily
\noindent\emph{Evidencia en Transformacion.py (líneas 99-116):}
\end{quote}
\begin{lstlisting}
    fecha = pd.to_datetime(s, dayfirst=True, errors="coerce")

    # Intento 2: serial numérico de Excel
    s_num = pd.to_numeric(s, errors="coerce")
    fecha_serial = pd.to_datetime(
        s_num, unit="D", origin="1899-12-30", errors="coerce"
[...]
\end{lstlisting}

  \item src/analisis/produccion.py:199: pivot.rename no falla si falta 'FEMENINO', pero pivot['prom\_femenino'] sí lanza KeyError en áreas sin mujeres.
\begin{quote}\footnotesize\ttfamily
\noindent\emph{Evidencia en src/analisis/produccion.py (línea 199):}
\end{quote}
\begin{lstlisting}
    pivot["razon_f_m"] = (pivot["prom_femenino"] / pivot["prom_masculino"]).round(3)
\end{lstlisting}

  \item src/analisis/produccion.py:89: división por len(autores\_unicos) -> ZeroDivisionError si prod está vacío (caso límite).
\begin{quote}\footnotesize\ttfamily
\noindent\emph{Evidencia en src/analisis/produccion.py (línea 89):}
\end{quote}
\begin{lstlisting}
        "pct_autores_reconocidos": round(len(overlap) / len(autores_unicos) * 100, 2),
\end{lstlisting}

  \item src/analisis/produccion.py:63-65: isin usa IDs normalizados a Int64; si id\_persona\_pr llega con ceros a la izquierda desde el XLSX (leído como str) y prod como int, el cruce pierde matches (depende de la inferencia de tipos).
\begin{quote}\footnotesize\ttfamily
\noindent\emph{Evidencia en src/analisis/produccion.py (líneas 63-65):}
\end{quote}
\begin{lstlisting}
    ids_reconocidos = set(inv["ID_PERSONA_PR"].dropna().astype("Int64"))
    p = prod.copy()
    p["es_reconocido"] = p["ID_PERSONA_PD"].isin(ids_reconocidos)
\end{lstlisting}

\end{itemize}
\end{tcolorbox}

\paragraph{Recomendación de auditoría}
\textbf{mantener/mejorar --- módulo metodológicamente sólido y bien documentado; priorizar la robustez de tipos/años y la reproducibilidad del input crudo (o versionar un parquet de producción) para que el cruce sea reproducible sin 1,2 GB.}

\medskip
\hrule

\subsubsection{src/analisis/redes.py}
\textbf{Rol:} analisis · \textbf{Líneas:} 225

\paragraph{Descripción funcional}
Análisis de redes de co-filiación institucional (Issues \#15/\#16): normaliza nombres de instituciones, extrae pares de instituciones de INST\_FILIA (separador '|'), construye grafos NetworkX, calcula métricas (grado, betweenness, densidad, componentes) y genera HTML interactivo con Pyvis. Entrada: DataFrame transformado con INST\_FILIA; salida: DataFrames de pares/nodos/aristas (evidencias/redes\_*.csv) y HTML en hallazgos/.

\paragraph{Análisis de la lógica}
normalizar\_institucion() aplica heurística de paréntesis (si el paréntesis contiene una institución madre, lo usa; si no, lo elimina), elimina sufijos SEDE/SECCIONAL y normaliza mayúsculas. construir\_pares() filtra filas con '|', divide en 2 partes (n=1), normaliza, elimina self-loops y ordena pares para que (A,B)==(B,A). construir\_grafo() agrega pesos por grupo (a,b). metricas\_grafo(), tabla\_nodos(), tabla\_aristas(), metricas\_por\_convocatoria() y subgrafo\_grado\_minimo() componen el resto. generar\_html\_pyvis() pinta nodos por grado ponderado y betweenness.

\paragraph{Fortalezas}
\begin{itemize}
  \item Normalización institucional con reglas documentadas y heurística de paréntesis --- valor real frente a las 2.762 cadenas crudas de inst\_filia.
  \item Eliminación de self-loops tras normalizar y orden canónico de pares: evita aristas duplicadas (A,B)==(B,A).
  \item Módulo completo (construcción, métricas, subgrafos, visualización) con lazy import de pyvis.
  \item Artefactos versionados en evidencias/redes\_*.csv y artifacts/sprint3\_redes/.
\end{itemize}

\paragraph{Debilidades}
\begin{itemize}
  \item CRÍTICO: construir\_pares() divide con n=1, capturando solo el PRIMER par de investigadores con 3+ instituciones; se pierden las combinaciones restantes (A,C), (B,C) -> el grafo subestima la co-filiación.
  \item El grafo global (anio=None) pondera por filas investigador x convocatoria: un investigador con el mismo par en 2 convocatorias cuenta doble, inflando el peso de aristas 'investigadores compartidos'.
  \item import networkx en mitad del archivo (línea 49) y un import networkx as nx\_local sin uso dentro de generar\_html\_pyvis --- estilo e higiene.
  \item Si INST\_FILIA termina en '|' (cadena vacía tras split), se crea un nodo con nombre '' en el grafo (no se filtra la cadena vacía, a diferencia de modelo/dimensional.py).
  \item El bloque 'if G.has\_edge' en construir\_grafo es código muerto (cada grupo (a,b) es único en el groupby).
\end{itemize}

\paragraph{Mejoras propuestas}
\begin{itemize}
  \item Usar itertools.combinations sobre las instituciones separadas para generar todos los pares de investigadores multi-filiados.
  \item Agregar peso por investigadores únicos (drop\_duplicates por ID\_PERSONA\_PR antes de contar) en el grafo global.
  \item Filtrar cadenas vacías y NaN después del split (strip + dropna + != '').
  \item Mover los imports al inicio del módulo; eliminar nx\_local no usado.
\end{itemize}

\paragraph{Ejecutabilidad}
Sí en el clon (solo DataFrame transformado, XLSX presente); requiere networkx y, para generar\_html\_pyvis, pyvis (ambos declarados en pyproject.toml). Sin datos/raw ni .duckdb.

\paragraph{¿Produce lo esperado?}
Sí: evidencias/redes\_pares\_cofiliacion.csv, redes\_nodos\_global.csv, redes\_aristas\_global.csv y redes\_metricas\_por\_convocatoria.csv existen, junto con artifacts/sprint3\_redes/fig01-03.png; los HTML Pyvis los genera sprint3\_grafo\_interactivo.py. Coincide con el README (Issues \#15/\#16).

\begin{tcolorbox}[riesgo]
\textbf{Bugs / problemas detectados:}
\begin{itemize}
  \item src/analisis/redes.py:62-64: str.split('|', n=1) descarta afiliaciones 3+; red incompleta (subestimación de aristas y nodos).
\begin{quote}\footnotesize\ttfamily
\noindent\emph{Evidencia en src/analisis/redes.py (líneas 62-64):}
\end{quote}
\begin{lstlisting}
    multi[["inst_a", "inst_b"]] = (
        multi["INST_FILIA"]
        .str.split("|", n=1, expand=True)
\end{lstlisting}

  \item src/analisis/redes.py:74-79: sorted() sobre pares que contienen '' (de 'A | ') crea un nodo vacío en el grafo; con tipos no comparables podría lanzar TypeError.
\begin{quote}\footnotesize\ttfamily
\noindent\emph{Evidencia en src/analisis/redes.py (líneas 74-79):}
\end{quote}
\begin{lstlisting}
    multi = multi[multi["inst_a"] != multi["inst_b"]].copy()
    # Normalizar orden para que (A,B) == (B,A)
    multi[["inst_a", "inst_b"]] = pd.DataFrame(
        [sorted([a, b]) for a, b in zip(multi["inst_a"], multi["inst_b"])],
        index=multi.index,
    )
\end{lstlisting}

  \item src/analisis/redes.py:93-98: el peso del grafo global cuenta filas investigador x convocatoria, no investigadores únicos: 'investigadores compartidos' queda inflado entre años.
\begin{quote}\footnotesize\ttfamily
\noindent\emph{Evidencia en src/analisis/redes.py (líneas 93-98):}
\end{quote}
\begin{lstlisting}
    for (a, b), grp in pares.groupby(["inst_a", "inst_b"]):
        peso = len(grp)
        if G.has_edge(a, b):
            G[a][b]["weight"] += peso
        else:
            G.add_edge(a, b, weight=peso)
\end{lstlisting}

  \item src/analisis/redes.py:16-48: normalizar\_institucion devuelve el valor original (posiblemente sin estandarizar) en el caso base.
\begin{quote}\footnotesize\ttfamily
\noindent\emph{Evidencia en src/analisis/redes.py (líneas 16-48):}
\end{quote}
\begin{lstlisting}
def normalizar_institucion(nombre) -> str:
    """
    Colapsa variantes (sedes, parentesis, mayusculas) a una sola entidad.

    Reglas:
    - Si hay parentesis y contiene una institucion (Universidad/Instituto/etc),
[...]
\end{lstlisting}

\end{itemize}
\end{tcolorbox}

\paragraph{Recomendación de auditoría}
\textbf{refactorizar --- análisis con buen valor analítico, pero la pérdida de pares multi-filiación y el conteo duplicado de pesos comprometen las métricas publicadas (n\_aristas, peso de aristas); corregir antes de citar resultados de red.}

\medskip
\hrule

\subsubsection{src/analisis/territorial.py}
\textbf{Rol:} analisis · \textbf{Líneas:} 107

\paragraph{Descripción funcional}
Análisis de concentración territorial (Issue \#13): índice Herfindahl-Hirschman (HHI) por convocatoria, top territorios y tabla de cuotas departamento x año. Entrada: DataFrame transformado; salida: DataFrames exportados a evidencias/territorial\_*.csv.

\paragraph{Análisis de la lógica}
hhi() normaliza conteos a cuotas y suma sus cuadrados (devuelve NaN si la suma es 0). hhi\_por\_convocatoria() filtra 'NO REPORTADO' y dropna, agrupa por año y reporta hhi, n\_total y n\_territorios. top\_territorios() calcula conteos por (año, territorio) con pct sobre el total anual y filtra los n territorios más frecuentes del total acumulado. tabla\_cuotas() pivotea departamento x año con las cuotas. sprint2\_territorial.py los consume (HHI por departamento y por región).

\paragraph{Fortalezas}
\begin{itemize}
  \item HHI implementado correctamente (cuotas normalizadas, rango [0,1]) con docstring de interpretación.
  \item Filtro explícito de 'NO REPORTADO' y dropna antes de calcular.
  \item Métricas auxiliares (n\_total, n\_territorios) dan contexto al índice.
\end{itemize}

\paragraph{Debilidades}
\begin{itemize}
  \item top\_territorios mezcla dos alcances: el 'top n' se calcula sobre el total acumulado del período, pero pct se calcula por año -> un territorio top en un año específico puede quedar excluido de la tabla (pérdida de información por año).
  \item hhi() con una sola unidad devuelve 1.0 (concentración total), correcto matemáticamente pero debe interpretarse junto con n\_territorios.
  \item El filtro de 'NO REPORTADO' depende de que Transformacion.py haya etiquetado así los faltantes: acoplamiento implícito a un valor mágico.
  \item Sin normalización adicional de nombres de departamentos (p.ej. 'Bogotá D.C.' vs 'Bogota') más allá del upper del pipeline.
\end{itemize}

\paragraph{Mejoras propuestas}
\begin{itemize}
  \item Separar el alcance del top-n (por año o global) en parámetros explícitos.
  \item Mover el filtro de 'NO REPORTADO' a una constante compartida (p.ej. en analisis/\_\_init\_\_) para evitar valores mágicos.
  \item Añadir HHI normalizado (HHI*) para comparar entre años con distinto número de territorios.
  \item Tests: serie de conteos, suma cero, una sola unidad.
\end{itemize}

\paragraph{Ejecutabilidad}
Sí en el clon: solo requiere DataFrame transformado (XLSX presente). Sin datos/raw ni .duckdb.

\paragraph{¿Produce lo esperado?}
Sí: evidencias/territorial\_hhi\_por\_convocatoria.csv, territorial\_hhi\_region\_por\_convocatoria.csv y territorial\_cuotas\_departamento.csv existen, junto con artifacts/sprint2\_territorial/fig01-03.png --- consistente con el README (Issue \#13).

\begin{tcolorbox}[riesgo]
\textbf{Bugs / problemas detectados:}
\begin{itemize}
  \item src/analisis/territorial.py:81-87: el top-n global combinado con pct anual puede producir tablas donde ningún año muestre el territorio \#1 de ese año; no es error de cálculo, pero distorsiona el concepto de 'top territorios por convocatoria'.
\begin{quote}\footnotesize\ttfamily
\noindent\emph{Evidencia en src/analisis/territorial.py (líneas 81-87):}
\end{quote}
\begin{lstlisting}
    top = (
        sub[col_geo]
        .value_counts()
        .head(n)
        .index.tolist()
    )
[...]
\end{lstlisting}

  \item src/analisis/territorial.py:44 y 71: si col\_geo contuviera valores no-cadena (el NaN ya se dropea, pero p.ej. códigos numéricos), .str.strip() lanzaría AttributeError --- depende del contrato de ingesta.
\begin{quote}\footnotesize\ttfamily
\noindent\emph{Evidencia en src/analisis/territorial.py (línea 44):}
\end{quote}
\begin{lstlisting}
    sub = sub[sub[col_geo].str.strip().str.upper() != "NO REPORTADO"]
\end{lstlisting}

\end{itemize}
\end{tcolorbox}

\paragraph{Recomendación de auditoría}
\textbf{mantener --- análisis correcto y exportable; mejorar la semántica del top-n y centralizar el filtro de faltantes.}

\medskip
\hrule

\subsection{Orquestación de análisis (scripts por sprint)}
\noindent Archivos: 17.

\subsubsection{scripts/generar\_manual.py}
\textbf{Rol:} documentacion · \textbf{Líneas:} 610

\paragraph{Descripción funcional}
Genera docs/manual.ipynb, un notebook ejecutable de 49 celdas (23 markdown + 26 código) que reproduce los doce hallazgos del observatorio de investigadores MinCiencias. Construye el notebook programáticamente con nbformat: va apilando celdas markdown y de código mediante los helpers md()/code() y las escribe en docs/ al final. Cada sección (1-14) combina texto explicativo con cifras del observatorio y celdas que reutilizan el pipeline moderno (src/ingesta, src/Transformacion.py, src/analisis/*) en lugar de duplicar lógica.

\paragraph{Análisis de la lógica}
El script no ejecuta análisis: solo ensambla el notebook. Declara ROOT como el padre del directorio del script y crea docs/ con mkdir(exist\_ok=True). Luego encadena 14 secciones: portada, setup (celda que inserta src/ en sys.path), carga de datos (from ingesta import cargar\_consolidado, cargar\_produccion; transformar; normalizar\_produccion), y secciones temáticas que importan funciones de src/analisis (calidad, longitudinal, territorial, genero, diversidad, produccion). Incrusta figuras mediante IFrame con rutas relativas a artifacts/sprint6\_sankey/*.html y CSV ya generados en evidencias/ (tabla\_maestra\_ies.csv, geografia\_*.csv, ocde\_*.csv). Cuatro celdas lanzan subprocess.run(["python3", "scripts/sprint6\_*.py"], cwd=ROOT, check=True) para regenerar Sankeys y CSVs. Al final, main() crea nbformat.v4.new\_notebook() con metadata kernelspec python3 y escribe el JSON en docs/manual.ipynb con codificación utf-8.

\paragraph{Fortalezas}
\begin{itemize}
  \item Estructura clara y reproducible: 49 celdas en 14 secciones numeradas, con markdown explicativo de cada hallazgo, incluyendo citas textuales del director y caveats metodológicos (Emérito vitalicio, diversidad solo capturada desde 2021, comparación DANE vs población con educación superior).
  \item Reutiliza el pipeline moderno (ingesta, Transformacion, analisis.*) en las celdas en lugar de duplicar lógica de análisis, lo que mantiene coherencia con scripts/sprint*.py.
  \item Documenta explícitamente la decisión metodológica de tratamiento de atípicos (eliminación sobre imputación) con justificación cuantitativa, y ofrece la alternativa imputar\_mediana aunque no se aplique por defecto.
  \item El notebook generado es versionable y el script es idempotente (siempre escribe el mismo manual.ipynb desde cero).
  \item El notebook existente en el clon coincide con lo que el script genera: 49 celdas (23 markdown + 26 código), mismo título UTF-8 correcto y mismas celdas clave (carga, subprocess, matrices) --- verificado byte a byte en la sección relevante.
\end{itemize}

\paragraph{Debilidades}
\begin{itemize}
  \item Cifras hardcodeadas en markdown (22 atípicos de edad, 77.237 filas, 51\% Bogotá+Antioquia, 242 investigadores Antioquia$\rightarrow$Bogotá, retención del 92\%) que nunca se recalculan; ya hay discrepancia real: el docstring de src/analisis/calidad.py dice 'diez registros' y '0.013\%', mientras el manual dice 22 y 0,03\%.
  \item El script asume el dataset completo de 6 convocatorias (77.237 filas), pero el clon solo dispone del XLSX de respaldo datos/tarea\_join/investigadores\_consolidado.xlsx con 3 convocatorias (2017, 2019, 2021) y 50.891 filas, por lo que las cifras regeneradas no coincidirían con el texto (p. ej. 'Cinco diagramas' de Sankey solo podrían ser 2 pares).
  \item No declara dependencias del script ni del notebook: nbformat (necesario para ejecutar el script), seaborn e IPython (importados por las celdas) no están en requirements.txt.
  \item Depende de subprocess con 'python3' hardcodeado en 4 celdas, lo que lo hace no portable a Windows.
  \item Rutas relativas en IFrame ("../artifacts/...") dependen del cwd del kernel de Jupyter; frágiles si el notebook se abre desde otro directorio.
  \item No valida que los CSV leídos existan ni que los subprocess tengan éxito con mensajes claros (check=True lanza excepción genérica).
\end{itemize}

\paragraph{Mejoras propuestas}
\begin{itemize}
  \item Corregir la celda de matrices: matrices\_todos\_periodos() retorna dict de tuplas (conteos, probs); usar m[0] o cambiar la función para retornar dict con claves 'conteos'/'probabilidades'.
  \item Reemplazar "python3" por sys.executable (o "python") en los subprocess.run para portabilidad a Windows.
  \item Agregar nbformat, seaborn e IPython a requirements.txt.
  \item Recalcular cifras clave en celdas (n\_atipicos, top departamentos, flujos) y compararlas con asserts contra los valores del texto, o generar el markdown de hallazgos dinámicamente.
  \item Documentar que el XLSX de respaldo cubre solo 3 convocatorias y que las secciones 11-13 requieren datos/raw/produccion\_grupos.csv; añadir try/except en la carga de producción con mensaje de acción clara.
  \item Reemplazar subprocess por llamadas directas a funciones de src/analisis dentro del notebook (más robusto y auditable), o usar rutas basadas en ROOT en lugar de rutas relativas en IFrame.
\end{itemize}

\paragraph{Ejecutabilidad}
El script en sí (python scripts/generar\_manual.py) es ejecutable tal cual en el clon siempre que nbformat esté instalado (no está en requirements.txt, aunque probablemente sí en .venv por Jupyter); escribe docs/manual.ipynb y no requiere datos. El notebook generado NO es ejecutable de punta a punta en el clon: (1) la celda de carga llama cargar\_produccion(), que lanza FileNotFoundError porque datos/raw/produccion\_grupos.csv está ausente (gitignored), rompiendo las secciones 11-13; (2) las 4 celdas con subprocess.run(["python3", ...]) fallan en Windows (python3 no es ejecutable estándar; sería python o py); (3) la celda de matrices accede m["conteos"] sobre tuplas y lanza TypeError; (4) los IFrame usan rutas relativas a artifacts/ que dependen del cwd. Los CSV de evidencias/ y los HTML de artifacts/sprint6\_sankey/ existen (generados previamente con el dataset completo), por lo que las celdas de solo-lectura (IFrame, read\_csv) funcionarían si las anteriores no fallaran.

\paragraph{¿Produce lo esperado?}
Sí genera docs/manual.ipynb con 49 celdas (23 markdown + 26 código) y metadata kernelspec python3 / language\_info 3.11. El notebook existente en el clon coincide con lo que el script generaría: misma estructura (49 celdas), mismo título (em-dash UTF-8 correcto, sin BOM) y las mismas celdas clave (carga de datasets, subprocess con python3, IFrame y las celdas bugueadas con m["conteos"] y con columna= en HHI); el notebook no tiene outputs ni execution\_count, es decir, nunca fue ejecutado. El notebook resultante NO es totalmente ejecutable: además del TypeError de m["conteos"] y del FileNotFoundError de cargar\_produccion(), la celda HHI lanza TypeError por el kwarg 'columna=' (el parámetro real es col\_geo), usa import seaborn e IPython.display sin declarar dependencias y rutas relativas frágiles; solo la sección 1 y parcialmente la 2 serían viables de forma inmediata con el XLSX de respaldo.

\begin{tcolorbox}[riesgo]
\textbf{Bugs / problemas detectados:}
\begin{itemize}
  \item Línea 197 (celda 'Trayectoria longitudinal'): matrices\_todos\_periodos() (src/analisis/longitudinal.py:164-179) retorna \{periodo: (conteos\_df, probs\_df)\}, pero la celda hace m["conteos"] $\rightarrow$ TypeError: tuple indices must be integers or slices, not str en runtime; el notebook existente contiene el mismo bug.
\begin{quote}\footnotesize\ttfamily
\noindent\emph{Evidencia en src/analisis/longitudinal.py (líneas 164-179):}
\end{quote}
\begin{lstlisting}
def matrices_todos_periodos(
    panel: pd.DataFrame, incluir_desaparece: bool = True
) -> dict:
    """
    Calcula matrices de transición para todos los pares consecutivos.

[...]
\end{lstlisting}

  \item Línea 251 (celda HHI de la sección 5): hhi\_por\_convocatoria(inv\_clean, columna="NME\_DEPARTAMENTO\_RES\_PR") pasa el kwarg 'columna', pero el parámetro real de la función es 'col\_geo' (src/analisis/territorial.py:32) $\rightarrow$ TypeError en runtime; la sección 5 no es viable aunque no requiera producción.
\begin{quote}\footnotesize\ttfamily
\noindent\emph{Evidencia en src/analisis/territorial.py (línea 32):}
\end{quote}
\begin{lstlisting}
def hhi_por_convocatoria(
\end{lstlisting}

  \item Líneas 213, 318, 357 y 509: subprocess.run(["python3", "scripts/sprint6\_*.py"], ...) --- 'python3' no existe en Windows (el clon vive en Windows); falla con FileNotFoundError. Debería ser sys.executable o 'python'.
  \item Celda de carga (líneas 107-117): prod = normalizar\_produccion(cargar\_produccion()) lanza FileNotFoundError en el clon porque datos/raw/produccion\_grupos.csv no existe; rompe secciones 11, 12 y 13.
  \item Línea 136-140 vs src/analisis/calidad.py:1-26: el manual afirma '22 registros con edad superior a 100 años' y '0,03\%', pero el módulo de calidad documenta 'diez registros' y '0.013\%' --- cifras contradictorias sin recalcular.
\begin{quote}\footnotesize\ttfamily
\noindent\emph{Evidencia en src/analisis/calidad.py (líneas 1-26):}
\end{quote}
\begin{lstlisting}
"""
analisis/calidad.py

Validación documentada de calidad de los datos del padrón de investigadores
reconocidos. La decisión metodológica adoptada se documenta de forma explícita
para que cualquier auditor pueda reproducirla y discutirla.
[...]
\end{lstlisting}

  \item Líneas 44-47 de la portada: dice 'los doce hallazgos' pero enumera 14 secciones numeradas (1-14) --- incoherencia menor de numeración.
  \item requirements.txt no incluye nbformat (necesario para ejecutar el script), ni seaborn ni IPython (importados por celdas del notebook).
  \item Celda setup (líneas 89-94): ROOT = pathlib.Path.cwd() depende del directorio desde donde se lance el kernel; el fallback a ROOT.parent solo cubre un nivel, frágil en Jupyter Lab/VS Code con cwd distinto.
\end{itemize}
\end{tcolorbox}

\paragraph{Recomendación de auditoría}
\textbf{mejorar --- El script cumple bien su rol documental (estructura clara, reuso de src/, notebook reproducible y coincidente con el del clon), pero no es ejecutable de punta a punta en el estado actual del repositorio: hay un bug de API real (m["conteos"]), dependencia de un CSV de producción ausente, subprocess no portables a Windows y cifras hardcodeadas que ya divergen del código de calidad. Se recomienda una mejora enfocada: corregir la celda de matrices, usar sys.executable, declarar dependencias y recargar/validar las cifras --- sin necesidad de refactorizar la arquitectura del script.}

\medskip
\hrule

\subsubsection{scripts/sprint2\_duckdb.py}
\textbf{Rol:} orquestacion · \textbf{Líneas:} 127

\paragraph{Descripción funcional}
Construye el modelo dimensional estrella (7 dimensiones + fact\_clasificacion + bridge\_hecho\_institucion) en datos/processed/observatorio.duckdb a partir del dataset consolidado transformado, y lo valida con 6 consultas de negocio (investigadores por convocatoria, distribución de categorías, top-10 instituciones, gran área, concentración Bogotá/Medellín/Cali, edad promedio por categoría y año). Usa src/ingesta.cargar\_consolidado, src/Transformacion.transformar y src/modelo/dimensional (cargar\_en\_duckdb, resumen\_modelo, DB\_PATH). Salidas: la base DuckDB y evidencias/duckdb\_resumen\_modelo.txt.

\paragraph{Análisis de la lógica}
main() encadena [1/3] transformar(cargar\_consolidado()), [2/3] cargar\_en\_duckdb(df, path=DB\_PATH, verbose=True) + resumen\_modelo(con) y [3/3] ejecución de CONSULTAS (dict nombre$\rightarrow$SQL) con con.execute(sql).fetchdf(), imprimiendo cada resultado y acumulando lineas\_reporte que se escriben al final en evidencias/duckdb\_resumen\_modelo.txt (encoding utf-8). La construcción del esquema y el DROP/CREATE idempotente viven en src/modelo/dimensional.py, no en el script; este solo orquesta carga, validación y persistencia del reporte. Cierra la conexión al final.

\paragraph{Fortalezas}
\begin{itemize}
  \item Modelo estrella explícito y bien documentado en el docstring (dims, fact, bridge N:N con dim\_institucion), coherente con src/modelo/dimensional.py.
  \item Las consultas de validación son preguntas de negocio reales (concentración territorial, edad por categoría, distribución OCDE), no solo conteos triviales; funcionan como smoke test del modelo.
  \item Idempotente gracias al DROP TABLE IF EXISTS + CREATE en cargar\_en\_duckdb: permite re-ejecución sin estado corrupto acumulado.
  \item El reporte de texto en evidencias es verificable y de hecho existe en el clon con el formato exacto que genera este script.
\end{itemize}

\paragraph{Debilidades}
\begin{itemize}
  \item Usa duckdb fetchdf(), API deprecada desde duckdb 0.9.0 en favor de df(); con el pin duckdb \textasciicircum{}1.0 del pyproject, una instalación reciente puede emitir DeprecationWarning o directamente fallar (riesgo de ejecutabilidad futura).
  \item resumen\_modelo(con) solo imprime en stdout: el resumen del modelo estrella (conteos por tabla) NO queda en el reporte de evidencias, que solo captura las 6 consultas.
  \item Sin try/finally: si una consulta o la construcción falla, la conexión duckdb queda abierta y no hay mensaje de error contextual.
  \item No valida el total de registros (esperado 77.237) contra fact\_clasificacion ni reporta duplicados de id\_persona\_pr por convocatoria.
  \item DB\_PATH está fijado en src/modelo/dimensional.py y no hay parámetro CLI; el script no permite apuntar a otra base ni modo --solo-consultas.
\end{itemize}

\paragraph{Mejoras propuestas}
\begin{itemize}
  \item Reemplazar fetchdf() por df() (línea 110) y también en src/modelo/dimensional.py:283.
  \item Incluir el output de resumen\_modelo() en lineas\_reporte para que la evidencia contenga el esquema completo.
  \item Envolver el flujo en try/finally con con.close() y capturar excepción con mensaje accionable.
  \item Añadir una consulta de sanidad: COUNT(*) de fact\_clasificacion $\approx$ 77.237 y verificar que ninguna dim tenga claves NULL usadas por los JOINs.
  \item Parametrizar DB\_PATH vía argparse/entorno para reutilizar el script en entornos CI.
\end{itemize}

\paragraph{Ejecutabilidad}
Ejecutable en el clon solo tras instalar dependencias: pyproject declara duckdb \textasciicircum{}1.0, pero el .venv del clon está vacío (solo pip/setuptools) y requirements.txt no incluye duckdb. El input sí existe: cargar\_consolidado() prioriza datos/raw/investigadores\_consolidado.csv (NO existe, gitignored) y cae al fallback datos/tarea\_join/investigadores\_consolidado.xlsx (SÍ existe). La base datos/processed/observatorio.duckdb no existe en el clon (solo .gitkeep) y el script la crea. Riesgo adicional: fetchdf() deprecado según la versión de duckdb instalada.

\paragraph{¿Produce lo esperado?}
SÍ, con matiz: evidencias/duckdb\_resumen\_modelo.txt existe en el clon y su contenido (encabezado 'Modelo estrella --- observatorio.duckdb' + las 6 consultas con sus to\_string) coincide exactamente con lo que genera el script, lo que demuestra que se ejecutó al menos una vez con datos. El archivo observatorio.duckdb NO está en el clon (gitignored, datos/processed solo contiene .gitkeep), por lo que esa salida no es reproducible sin re-ejecución. Coincidencia declarado-vs-generado: OK para el reporte.

\begin{tcolorbox}[riesgo]
\textbf{Bugs / problemas detectados:}
\begin{itemize}
  \item scripts/sprint2\_duckdb.py:110 --- con.execute(sql).fetchdf(): fetchdf() está deprecado desde duckdb 0.9.0 (reemplazado por .df()); con duckdb \textasciicircum{}1.0 en pyproject.toml una versión reciente puede eliminarlo, rompiendo el script en runtime.
\begin{quote}\footnotesize\ttfamily
\noindent\emph{Evidencia en scripts/sprint2\_duckdb.py (línea 110):}
\end{quote}
\begin{lstlisting}
        resultado = con.execute(sql).fetchdf()
\end{lstlisting}

  \item scripts/sprint2\_duckdb.py:104 --- resumen\_modelo(con) imprime a stdout pero su salida no se agrega a lineas\_reporte: la evidencia guardada (duckdb\_resumen\_modelo.txt) omite el resumen del modelo estrella que el propio docstring describe como parte de la validación.
\begin{quote}\footnotesize\ttfamily
\noindent\emph{Evidencia en scripts/sprint2\_duckdb.py (línea 104):}
\end{quote}
\begin{lstlisting}
    resumen_modelo(con)
\end{lstlisting}

  \item scripts/sprint2\_duckdb.py:97-123 --- sin try/finally: cualquier excepción entre [2/3] y [3/3] deja la conexión duckdb sin cerrar y sin reporte parcial.
\begin{quote}\footnotesize\ttfamily
\noindent\emph{Evidencia en scripts/sprint2\_duckdb.py (líneas 97-123):}
\end{quote}
\begin{lstlisting}
def main() -> None:
    print("[1/3] Cargando y transformando datos...")
    df = transformar(cargar_consolidado())
    print(f"      Shape: {df.shape}")

    print("\n[2/3] Construyendo modelo estrella en DuckDB...")
[...]
\end{lstlisting}

\end{itemize}
\end{tcolorbox}

\paragraph{Recomendación de auditoría}
\textbf{mantener --- es la base del modelo dimensional del observatorio y sus consultas de validación son valiosas; solo requiere la migración fetchdf()$\rightarrow$df(), incluir el resumen en el reporte y robustecer el manejo de conexión.}

\medskip
\hrule

\subsubsection{scripts/sprint2\_genero\_ocde.py}
\textbf{Rol:} orquestacion · \textbf{Líneas:} 182

\paragraph{Descripción funcional}
Analiza la representación femenina y la brecha de género por área OCDE (gran área, área) en las 6 convocatorias 2013-2021. Usa src/analisis/genero (pct\_femenino\_por\_area, tabla\_pivot\_pct\_femenino, brecha\_genero, evolucion\_pct\_femenino) sobre transformar(cargar\_consolidado()). Salidas: 4 figuras PNG en artifacts/sprint2\_genero\_ocde/ (barras, heatmap gran área$\times$convocatoria, líneas de evolución, brecha primera vs última convocatoria) y 4 CSVs en evidencias/.

\paragraph{Análisis de la lógica}
main() imprime el \% femenino global (filtrando solo FEMENINO/MASCULINO), genera las 4 figuras con funciones dedicadas (fig\_barras\_pct\_femenino, fig\_heatmap\_pct\_femenino, fig\_lineas\_evolucion, fig\_brecha\_genero), y exporta las tablas: pct por gran área, pivot con promedio, brecha (M\%-F\%) y pct por área. Cada función de figura llama a la función de analisis/genero correspondiente y guarda con plt.savefig(dpi=150) tras plt.tight\_layout(); usa matplotlib.use('Agg') para entornos headless.

\paragraph{Fortalezas}
\begin{itemize}
  \item Código limpio y directo: separación clara carga$\rightarrow$figuras$\rightarrow$exportación, funciones de figura autocontenidas y reutilizables.
  \item Uso correcto de la capa analitica (src/analisis/genero): el script no recalcula métricas, solo orquesta y visualiza.
  \item Backend Agg antes de importar pyplot, correcto para ejecución por consola/CI.
  \item Salidas verificadas: las 4 figuras y los 4 CSVs existen en el clon, en las rutas declaradas en el docstring.
\end{itemize}

\paragraph{Debilidades}
\begin{itemize}
  \item Escalas fijas que saturan la información: heatmap con vmin=0, vmax=60 (fig02) y líneas con ylim(0,70) (fig03); si las áreas STEM superan 60\% de presencia masculina, el color y la curva quedan aplastados y no se distingue entre áreas.
  \item fig\_brecha\_genero alinea b0 con reindex sobre las áreas de la última convocatoria: si una gran área existe en a0 pero no en a1 (o viceversa), se produce NaN silencioso y la comparación pierde áreas sin ninguna advertencia.
  \item El resumen global (pct femenino total) solo se imprime en consola; no se persiste como evidencia.
  \item Sin validación del mínimo de registros por área/convocatoria: celdas con n muy bajo (posibles identificadores) se grafican igual que áreas con miles de registros.
\end{itemize}

\paragraph{Mejoras propuestas}
\begin{itemize}
  \item Autoeescalar el heatmap (vmin/vmax del dato o por cuantiles) o documentar por qué 60\% es el tope; igual para ylim de fig03.
  \item En fig\_brecha\_genero usar union de áreas (b0 y b1) y rellenar NaN con 0 o marcar 'sin dato', e imprimir advertencia.
  \item Exportar el resumen global a un CSV de evidencias.
  \item Anotar n por celda (o tamaño de punto) en el heatmap para visibilizar áreas con pocos investigadores.
\end{itemize}

\paragraph{Ejecutabilidad}
Ejecutable en el clon solo tras instalar dependencias (matplotlib, seaborn, pandas; .venv del clon vacío, requirements.txt no las lista aunque pyproject sí). El input existe: cargar\_consolidado() usa el fallback datos/tarea\_join/investigadores\_consolidado.xlsx (presente). No requiere datos/raw ni duckdb. Con poetry install el script corre de punta a punta en headless.

\paragraph{¿Produce lo esperado?}
SÍ, coincidencia total: existen artifacts/sprint2\_genero\_ocde/fig01\_barras\_pct\_femenino.png, fig02\_heatmap\_pct\_femenino.png, fig03\_lineas\_evolucion.png, fig04\_brecha\_genero.png y evidencias/genero\_pct\_femenino\_por\_gran\_area.csv, genero\_tabla\_pivot\_gran\_area.csv, genero\_brecha\_por\_gran\_area.csv, genero\_pct\_femenino\_por\_area.csv. Las 8 salidas declaradas en el docstring y en el README están presentes y con los nombres exactos, lo que indica ejecución exitosa previa.

\begin{tcolorbox}[riesgo]
\textbf{Bugs / problemas detectados:}
\begin{itemize}
  \item scripts/sprint2\_genero\_ocde.py:132-135 --- b0.reindex(areas) con areas derivada solo de b1: si el conjunto de gran áreas difiere entre convocatorias extremas, las barras de b0 quedan en NaN (no se dibujan) sin advertencia, y las áreas exclusivas de a0 desaparecen del gráfico.
\begin{quote}\footnotesize\ttfamily
\noindent\emph{Evidencia en scripts/sprint2\_genero\_ocde.py (líneas 132-135):}
\end{quote}
\begin{lstlisting}
    ax.bar([i - width / 2 for i in x], b0.reindex(areas), width,
           label=str(a0), color="steelblue", alpha=0.8)
    ax.bar([i + width / 2 for i in x], b1.reindex(areas), width,
           label=str(a1), color="tomato", alpha=0.8)
\end{lstlisting}

  \item scripts/sprint2\_genero\_ocde.py:81 --- sns.heatmap con vmin=0, vmax=60 fijos: cualquier gran área con \% femenino >60 (esperable en ciencias de la salud o educación) satura el rojo/verde de RdYlGn y pierde discriminación visual.
\begin{quote}\footnotesize\ttfamily
\noindent\emph{Evidencia en scripts/sprint2\_genero\_ocde.py (línea 81):}
\end{quote}
\begin{lstlisting}
        vmax=60,
\end{lstlisting}

  \item scripts/sprint2\_genero\_ocde.py:108 --- ax.set\_ylim(0, 70) en fig03: series que superen 70\% (posible en áreas pequeñas) quedan cortadas en el borde superior sin indicación.
\begin{quote}\footnotesize\ttfamily
\noindent\emph{Evidencia en scripts/sprint2\_genero\_ocde.py (línea 108):}
\end{quote}
\begin{lstlisting}
    ax.set_ylim(0, 70)
\end{lstlisting}

\end{itemize}
\end{tcolorbox}

\paragraph{Recomendación de auditoría}
\textbf{mantener --- análisis correcto y reproducible con salidas verificadas; las debilidades son de visualización (escalas fijas) y de robustez del alineamiento de áreas, corregibles con cambios menores.}

\medskip
\hrule

\subsubsection{scripts/sprint2\_matrices\_transicion.py}
\textbf{Rol:} orquestacion · \textbf{Líneas:} 125

\paragraph{Descripción funcional}
Calcula y visualiza las matrices de probabilidad de transición entre categorías (Junior/Asociado/Senior/Emérito) para los 5 periodos consecutivos de las 6 convocatorias (2013-2021), en dos variantes: observada (solo investigadores retenidos) y extendida (incluye el estado 'Desaparece'). Usa src/analisis/longitudinal (construir\_panel, comparar\_periodo, matriz\_transicion, matrices\_todos\_periodos). Salidas: 2 figuras de heatmaps en artifacts/sprint2\_transiciones/ y 20 CSVs (conteos y probabilidades $\times$ ext/obs $\times$ 5 periodos) en evidencias/.

\paragraph{Análisis de la lógica}
main() construye el panel (ID\_PERSONA\_PR, año, categoria, orden), imprime las matrices por periodo en consola con un bucle comparar\_periodo+matriz\_transicion, y luego recalcula todo con matrices\_todos\_periodos para las figuras y la exportación. fig\_heatmaps() dibuja un heatmap por periodo en una sola fila de subplots (1$\times$n) con sns.heatmap(annot=True, fmt='.2f', vmin=0, vmax=1). Los nombres de archivo derivan del periodo reemplazando el guion largo '--' por '\_' (p. ej. matriz\_probabilidades\_ext\_2015\_2017.csv).

\paragraph{Fortalezas}
\begin{itemize}
  \item Presenta explícitamente las dos variantes metodológicas (observada vs extendida con 'Desaparece'), lo que permite ver el sesgo de selección por deserción en lugar de ocultarlo.
  \item La lógica estadística está delegada en src/analisis/longitudinal (crosstab, renormalización por fila), manteniendo el script como orquestador puro.
  \item El manejo del guion largo '--' en nombres de archivo evita problemas de codificación de rutas en Windows.
  \item Salidas verificadas: las 2 figuras y los 20 CSVs de la nomenclatura actual existen en el clon.
\end{itemize}

\paragraph{Debilidades}
\begin{itemize}
  \item Cómputo duplicado: el bucle de impresión (líneas 83-91) recalcula las matrices que matrices\_todos\_periodos vuelve a calcular en la línea 94-95; con 77.237 registros y 5 periodos es ineficiente y propenso a divergencias si cambia la lógica.
  \item La matriz 'observada' renormaliza por fila excluyendo a los desaparecidos: sus probabilidades no suman a la población real y pueden sobreestimar la movilidad; el script no advierte de esta limitación interpretativa.
  \item Evidencias con nomenclatura heredada: en evidencias/ coexisten archivos de una versión anterior (matriz\_transicion\_extendida\_*.csv, matriz\_probabilidades\_extendida\_*.csv y variantes sin prefijo ext/obs como matriz\_transicion\_2017\_2019.csv) que este script ya no produce, duplicando y confundiendo la evidencia.
  \item fig\_heatmaps asume n>=1 periodo: con un solo par de años plt.subplots(1, n) devolvería un Axes escalar y zip() fallaría (robustez marginal).
  \item Sin validación de que el panel contenga exactamente 6 convocatorias / 5 periodos ni de tamaños de celda mínimos.
\end{itemize}

\paragraph{Mejoras propuestas}
\begin{itemize}
  \item Calcular matrices una sola vez (fuera del bucle de impresión) y reutilizarlas para print, figuras y exportación.
  \item Añadir advertencia explícita cuando la matriz observada descarte >X\% de la fila por deserción, y priorizar la lectura de la matriz extendida en el informe.
  \item Eliminar o archivar los CSVs legacy (extendida\_/sin prefijo) para dejar una sola convención de nombres.
  \item Proteger fig\_heatmaps para n==1 (ax = axes si escalar) y añadir un assert del número de periodos esperados.
  \item Escribir un test que verifique que las filas de la matriz de probabilidades sumen 1.0 (obs) y 1.0 incluyendo Desaparece (ext).
\end{itemize}

\paragraph{Ejecutabilidad}
Ejecutable en el clon solo tras instalar dependencias (pandas, matplotlib, seaborn; .venv del clon vacío, requirements.txt incompleto). El input existe: datos/tarea\_join/investigadores\_consolidado.xlsx vía cargar\_consolidado() (fallback del CSV raw inexistente). No requiere duckdb ni datos/raw. Con poetry install, corre de punta a punta.

\paragraph{¿Produce lo esperado?}
SÍ: existen artifacts/sprint2\_transiciones/fig01\_heatmaps\_ext.png y fig02\_heatmaps\_obs.png, y en evidencias/ los 20 CSVs con la nomenclatura actual (matriz\_transicion\_ext\_/matriz\_probabilidades\_ext\_/matriz\_transicion\_obs\_/matriz\_probabilidades\_obs\_ $\times$ 2013\_2014, 2014\_2015, 2015\_2017, 2017\_2019, 2019\_2021), coincidiendo con lo declarado (evidencias/matriz\_transicion\_*.csv, matriz\_probabilidades\_*.csv). Adicionalmente hay 8 archivos de nomenclatura antigua (extendida\_/sin prefijo) que evidencian ejecuciones previas con una versión distinta del script.

\begin{tcolorbox}[riesgo]
\textbf{Bugs / problemas detectados:}
\begin{itemize}
  \item scripts/sprint2\_matrices\_transicion.py:83-95 --- las matrices se calculan dos veces (bucle de impresión y matrices\_todos\_periodos): doble costo de cómputo sobre 77.237 registros y riesgo de que las cifras impresas difieran de las exportadas si la función cambiara entre ambas llamadas.
\begin{quote}\footnotesize\ttfamily
\noindent\emph{Evidencia en scripts/sprint2\_matrices\_transicion.py (líneas 83-95):}
\end{quote}
\begin{lstlisting}
    for a0, a1 in zip(anios, anios[1:]):
        comp = comparar_periodo(panel, a0, a1)
        _, probs_obs = matriz_transicion(comp, incluir_desaparece=False)
        _, probs_ext = matriz_transicion(comp, incluir_desaparece=True)
        print(f"\n{'='*55}")
        print(f"Periodo {a0}–{a1} — Observada (solo retenidos):")
[...]
\end{lstlisting}

  \item scripts/sprint2\_matrices\_transicion.py:49 --- plt.subplots(1, n) con n=1 devuelve un solo Axes (no iterable); el zip(axes, periodos) fallaría con TypeError si el panel tuviera un único periodo.
\begin{quote}\footnotesize\ttfamily
\noindent\emph{Evidencia en scripts/sprint2\_matrices\_transicion.py (línea 49):}
\end{quote}
\begin{lstlisting}
    fig, axes = plt.subplots(1, n, figsize=(5 * n, 4))
\end{lstlisting}

  \item Evidencias/ legacy: los archivos matriz\_transicion\_extendida\_2017\_2019.csv, matriz\_probabilidades\_extendida\_2019\_2021.csv y matriz\_transicion\_2017\_2019.csv (sin prefijo) ya no son generados por este script; quedan como evidencia huérfana que puede confundir la trazabilidad.
\end{itemize}
\end{tcolorbox}

\paragraph{Recomendación de auditoría}
\textbf{mantener --- es la entrega central del análisis longitudinal y produce evidencia verificable en el clon; requiere eliminar el cómputo duplicado, advertir el sesgo de la matriz observada y limpiar la nomenclatura legacy.}

\medskip
\hrule

\subsubsection{scripts/sprint2\_panel\_longitudinal.py}
\textbf{Rol:} orquestacion · \textbf{Líneas:} 229

\paragraph{Descripción funcional}
Orquesta el análisis longitudinal del Sprint 2 (Issue \#11): carga el consolidado (cargar\_consolidado + transformar), construye el panel de investigadores por convocatoria con src/analisis/longitudinal.py, genera 7 figuras (evolución total, por género, por gran área, categorías, nuevos ingresos, tracking por periodo, tasas de retención) y exporta 4+ CSVs de evidencias.

\paragraph{Análisis de la lógica}
main() en 4 pasos: [1/4] carga y transforma; [2/4] genera figuras descriptivas (funciones fig\_* con matplotlib/seaborn, backend Agg); [3/4] construye el panel (construir\_panel), calcula tracking (tracking\_todos\_periodos) y retención (tasa\_retencion\_por\_periodo); [4/4] exporta resumen, ingresos, tasas y tracking por categoría a evidencias/. Patrón claro: una función por figura, salidas documentadas en el docstring.

\paragraph{Fortalezas}
\begin{itemize}
  \item Docstring con uso y lista explícita de salidas (artifacts/ y evidencias/).
  \item Delega la lógica en src/analisis/longitudinal.py (separación de responsabilidades correcta).
  \item Usa matplotlib.use('Agg') (headless, reproducible).
  \item Maneja los 5 periodos consecutivos de las 6 convocatorias de forma genérica (zip de años).
\end{itemize}

\paragraph{Debilidades}
\begin{itemize}
  \item Estilo global de seaborn/matplotlib aplicado al módulo (efecto secundario al importar).
  \item Sin tests ni manejo de errores (si cargar\_consolidado fallara, crash sin mensaje claro).
  \item Figuras con tamaños fijos (figsize) poco adaptables a distintos formatos de reporte.
  \item No guarda los DataFrames intermedios (panel) para reutilización.
\end{itemize}

\paragraph{Mejoras propuestas}
\begin{itemize}
  \item Extraer la configuración de estilo a una función o archivo de tema.
  \item Añadir un pequeño test de humo (panel con datos sintéticos).
  \item Permitir configurar figsize por parámetro y guardar el panel en evidencias/.
\end{itemize}

\paragraph{Ejecutabilidad}
Sí, verificada empíricamente en la auditoría 2026: ejecutado en un clon fresco (venv Python, pandas 3.0.5) con EXIT=0 en 108.4 s. Requiere el XLSX de datos/tarea\_join (presente). NOTA: al correr sobre el XLSX versionado (3 convocatorias) genera solo 2 periodos de tracking, frente a los 5 periodos de las evidencias versionadas --- consecuencia del consolidado incompleto, no del script.

\paragraph{¿Produce lo esperado?}
Sí: genera las 7 figuras en artifacts/sprint2\_longitudinal/ (versionadas en el clon) y los CSVs en evidencias/ (panel\_longitudinal\_resumen\_todos\_periodos.csv, ingresos\_por\_convocatoria, tasas\_retencion, panel\_longitudinal\_cat\_*.csv). Los artefactos existentes coinciden con lo declarado, pero con 6 convocatorias (los regenerados con el clon tienen 3).

\begin{tcolorbox}[riesgo]
\textbf{Bugs / problemas detectados:}
\begin{itemize}
  \item Ninguno crítico; dependencia implícita de que el consolidado tenga las 6 convocatorias para reproducir los 5 periodos del README (problema de datos, no de código).
\end{itemize}
\end{tcolorbox}

\paragraph{Recomendación de auditoría}
\textbf{mantener --- bien estructurado y verificable; solo requiere el consolidado completo de 6 convocatorias para reproducir los resultados del README.}

\medskip
\hrule

\subsubsection{scripts/sprint2\_territorial.py}
\textbf{Rol:} orquestacion · \textbf{Líneas:} 142

\paragraph{Descripción funcional}
Analiza la concentración territorial de los investigadores con el índice de Herfindahl-Hirschman (HHI) por departamento y por región para las 6 convocatorias (2013-2021), más las cuotas (\%) por departamento y el top-10 de departamentos. Usa src/analisis/territorial (hhi\_por\_convocatoria, top\_territorios, tabla\_cuotas) sobre transformar(cargar\_consolidado()). Salidas: 3 figuras PNG en artifacts/sprint2\_territorial/ y 3 CSVs en evidencias/.

\paragraph{Análisis de la lógica}
main() genera fig\_top\_departamentos (barras apiladas por convocatoria del top 10 con pct por año), calcula e imprime el HHI departamental y regional, genera fig\_hhi\_evolucion (2 subplots lado a lado con anotaciones de valor) y fig\_heatmap\_cuotas (heatmap \% por departamento$\times$convocatoria, top 15), que además retorna cuotas\_pivot para exportar. Exporta territorial\_hhi\_por\_convocatoria.csv, territorial\_hhi\_region\_por\_convocatoria.csv y territorial\_cuotas\_departamento.csv.

\paragraph{Fortalezas}
\begin{itemize}
  \item Usa correctamente el HHI normalizado a [0,1] sobre cuotas reales (cálculo delegado en src/analisis/territorial.hhi), con exclusión de valores NO REPORTADO.
  \item Triple perspectiva complementaria: top-10, evolución HHI (depto y región) y cuotas completas $\rightarrow$ buena triangulación del hallazgo de concentración.
  \item Las figuras anotan los valores exactos de HHI, facilitando la lectura sin tabla auxiliar.
  \item Salidas verificadas: las 3 figuras y los 3 CSVs existen en el clon en las rutas declaradas.
\end{itemize}

\paragraph{Debilidades}
\begin{itemize}
  \item fig\_hhi\_evolucion mezcla tipos en las coordenadas de anotación: el eje X se plotea con strings (astype(str)) pero annotate recibe str(int(row['ANO\_CONVO\_INT'])), dependiendo de la conversión categórica implícita de matplotlib; frágil ante cambios de versión y de difícil depuración.
  \item int(row['ANO\_CONVO\_INT']) lanzaría ValueError si hubiera NaN en el año (los filtros de hhi\_por\_convocatoria no garantizan limpieza total del año).
  \item El HHI por región depende de la columna NME\_REGION\_RES\_PR derivada en Transformacion; el script no verifica que exista, fallando con KeyError poco descriptivo si la transformación cambia.
  \item Sin validación del tamaño muestral por año (el HHI de un año con pocos registros no es comparable) ni intervalo/estabilidad del índice.
\end{itemize}

\paragraph{Mejoras propuestas}
\begin{itemize}
  \item Anotar usando las mismas posiciones categóricas del plot (p. ej. range(len(data)) o plt.xticks con los strings) en lugar de pasar el string como coordenada.
  \item Proteger con .dropna() sobre ANO\_CONVO\_INT antes del gráfico y validar la presencia de NME\_REGION\_RES\_PR con un mensaje claro.
  \item Exportar también n\_total y n\_territorios junto al HHI (ya los retorna hhi\_por\_convocatoria) para contextualizar la evolución.
  \item Añadir test de humo: HHI Bogotá+Antioquia $\approx$ 51\% de concentración reportado en el README.
\end{itemize}

\paragraph{Ejecutabilidad}
Ejecutable en el clon solo tras instalar dependencias (pandas, matplotlib, seaborn; .venv del clon vacío, requirements.txt incompleto). El input existe: datos/tarea\_join/investigadores\_consolidado.xlsx vía cargar\_consolidado(). No requiere duckdb. Con poetry install, corre de punta a punta; el único supuesto frágil es la columna NME\_REGION\_RES\_PR producida por Transformacion.

\paragraph{¿Produce lo esperado?}
SÍ, coincidencia total: existen artifacts/sprint2\_territorial/fig01\_top\_departamentos.png, fig02\_hhi\_evolucion.png, fig03\_heatmap\_cuotas.png y evidencias/territorial\_hhi\_por\_convocatoria.csv, territorial\_hhi\_region\_por\_convocatoria.csv, territorial\_cuotas\_departamento.csv, todos con los nombres exactos del docstring y del README, lo que confirma ejecución previa exitosa.

\begin{tcolorbox}[riesgo]
\textbf{Bugs / problemas detectados:}
\begin{itemize}
  \item scripts/sprint2\_territorial.py:82 --- ax.annotate con x = str(int(row['ANO\_CONVO\_INT'])): se mezcla una cadena como coordenada de datos en un eje categórico (ploteado con astype(str)); funciona solo por la conversión unitaria implícita de matplotlib y es un punto de fallo frágil.
\begin{quote}\footnotesize\ttfamily
\noindent\emph{Evidencia en scripts/sprint2\_territorial.py (línea 82):}
\end{quote}
\begin{lstlisting}
                (str(int(row["ANO_CONVO_INT"])), row["hhi"]),
\end{lstlisting}

  \item scripts/sprint2\_territorial.py:82 --- int(row['ANO\_CONVO\_INT']) lanza ValueError si el año es NaN; hhi\_por\_convocatoria filtra geo y año (dropna en sub), pero el DataFrame impreso/graficado aquí no garantiza ausencia de NaN en otras filas.
\begin{quote}\footnotesize\ttfamily
\noindent\emph{Evidencia en scripts/sprint2\_territorial.py (línea 82):}
\end{quote}
\begin{lstlisting}
                (str(int(row["ANO_CONVO_INT"])), row["hhi"]),
\end{lstlisting}

  \item scripts/sprint2\_territorial.py:123 --- accede a NME\_REGION\_RES\_PR sin verificar su existencia: si Transformacion renombra/elimina esa columna, el script cae con KeyError sin mensaje orientativo.
\begin{quote}\footnotesize\ttfamily
\noindent\emph{Evidencia en scripts/sprint2\_territorial.py (línea 123):}
\end{quote}
\begin{lstlisting}
    hhi_region = hhi_por_convocatoria(df, col_geo="NME_REGION_RES_PR")
\end{lstlisting}

\end{itemize}
\end{tcolorbox}

\paragraph{Recomendación de auditoría}
\textbf{mantener --- el análisis HHI es metodológicamente sólido y las salidas están verificadas; corregir las anotaciones del gráfico de evolución y blindar la dependencia de la columna de región.}

\medskip
\hrule

\subsubsection{scripts/sprint3\_grafo\_interactivo.py}
\textbf{Rol:} orquestacion · \textbf{Líneas:} 72

\paragraph{Descripción funcional}
Genera dos visualizaciones interactivas HTML (Pyvis) del grafo de co-filiación institucional: el grafo completo y un subgrafo filtrado con grado >= 2. Usa src/analisis/redes (construir\_pares, construir\_grafo, subgrafo\_grado\_minimo, generar\_html\_pyvis, metricas\_grafo) sobre transformar(cargar\_consolidado()). Salidas: hallazgos/sprint3\_grafo\_completo.html (docstring: 498 nodos) y hallazgos/sprint3\_grafo\_filtrado.html (146 nodos).

\paragraph{Análisis de la lógica}
main() construye los pares de doble afiliación (construir\_pares sin filtro de año), el grafo global (construir\_grafo con anio=None), el subgrafo de grado mínimo, imprime las métricas de ambos y genera los dos HTML con generar\_html\_pyvis, que en src/analisis/redes.py calcula betweenness (weight='weight'), escala tamaño por grado ponderado y color por betweenness. No exporta CSVs ni figuras estáticas; sus únicos artefactos son los dos HTML en hallazgos/.

\paragraph{Fortalezas}
\begin{itemize}
  \item Entrega un artefacto de alto valor de comunicación (grafo interactivo navegable) y lo versiona en hallazgos/ como entregable, no como figura estática.
  \item Delega toda la lógica de red en src/analisis/redes (normalización de instituciones, pesos, betweenness), manteniendo el script mínimo y legible.
  \item El filtrado por grado >= 2 elimina nodos hoja y produce un subgrafo interpretable; las métricas de ambos grafos se imprimen para contraste.
  \item Salidas verificadas: ambos HTML existen en el clon y están versionados en git.
\end{itemize}

\paragraph{Debilidades}
\begin{itemize}
  \item Los títulos de los HTML hardcodean '(2019)' (líneas 53 y 62) pero el grafo se construye con TODAS las convocatorias (construir\_pares(df) sin filtro de año y construir\_grafo con anio=None): la etiqueta contradice el dato real (co-filiación global 2013-2021) y es engañosa para quien abre el HTML.
  \item Los conteos '498 nodos' y '146 nodos' del docstring son literales de una corrida puntual; al re-ejecutar con datos actualizados el docstring queda desactualizado (drift documental).
  \item No genera evidencias tabulares (pares, métricas) propias: quien quiera verificar los 498/146 nodos debe re-ejecutar o abrir el HTML.
  \item Sin parámetro de año: no permite generar el grafo por convocatoria (p. ej. solo 2019) aunque la capa analitica lo soporta (construir\_grafo acepta anio).
\end{itemize}

\paragraph{Mejoras propuestas}
\begin{itemize}
  \item Eliminar el año hardcodeado de los títulos o calcularlo dinámicamente (p. ej. '2013-2021' derivado del rango real del panel) para no contradecir los datos.
  \item Derivar los conteos de nodos del docstring dinámicamente o quitar los números literales.
  \item Exportar además evidencias/redes\_grafo\_metricas.csv (m completo y m filtrado) para trazabilidad de los 498/146 nodos.
  \item Añadir un flag --anio para reproducir la etiqueta 2019 de forma honesta (filtrando el grafo por convocatoria).
\end{itemize}

\paragraph{Ejecutabilidad}
Ejecutable en el clon solo tras instalar dependencias: pyproject declara networkx \textasciicircum{}3.2 y pyvis \textasciicircum{}0.3, pero el .venv del clon está vacío (solo pip/setuptools) y requirements.txt no los lista. El input existe: datos/tarea\_join/investigadores\_consolidado.xlsx vía cargar\_consolidado(). No requiere duckdb. Con poetry install, corre de punta a punta y sobrescribe los HTML existentes.

\paragraph{¿Produce lo esperado?}
SÍ: hallazgos/sprint3\_grafo\_completo.html y hallazgos/sprint3\_grafo\_filtrado.html existen en el clon (rutas exactas del docstring y del README, que documenta 'hallazgos/ --- HTMLs interactivos (grafos Pyvis)'), lo que confirma ejecución previa. Los conteos 498/146 nodos no son verificables por análisis estático sin abrir los HTML, pero la existencia y el versionado de ambos archivos es coherente con una corrida exitosa. El script no declara ni produce figuras PNG ni CSVs, y así es.

\begin{tcolorbox}[riesgo]
\textbf{Bugs / problemas detectados:}
\begin{itemize}
  \item scripts/sprint3\_grafo\_interactivo.py:53 y 62 --- títulos hardcodeados 'Co-filiacion institucional --- Grafo completo (2019)' y '(2019)': el grafo se construye con todas las convocatorias (construir\_pares(df) línea 39 y construir\_grafo(pares) línea 40 con anio=None en src/analisis/redes.py:83-90), por lo que la etiqueta 2019 es incorrecta y contradice el alcance real del grafo.
\begin{quote}\footnotesize\ttfamily
\noindent\emph{Evidencia en scripts/sprint3\_grafo\_interactivo.py (línea 53):}
\end{quote}
\begin{lstlisting}
        titulo="Co-filiacion institucional — Grafo completo (2019)",
\end{lstlisting}

\begin{quote}\footnotesize\ttfamily
\noindent\emph{Evidencia en src/analisis/redes.py (líneas 83-90):}
\end{quote}
\begin{lstlisting}
def construir_grafo(pares: pd.DataFrame, anio: int = None) -> nx.Graph:
    """
    Construye grafo de co-filiacion a partir de los pares.

    Si anio != None filtra por convocatoria. Peso de arista = n investigadores compartidos.
    """
[...]
\end{lstlisting}

  \item scripts/sprint3\_grafo\_interactivo.py:10-11 --- el docstring fija '498 nodos' y '146 nodos' como constantes; si el dataset cambia, el docstring y los HTML pueden divergir sin ningún aviso.
\begin{quote}\footnotesize\ttfamily
\noindent\emph{Evidencia en scripts/sprint3\_grafo\_interactivo.py (líneas 10-11):}
\end{quote}
\begin{lstlisting}
    hallazgos/sprint3_grafo_completo.html      (498 nodos)
    hallazgos/sprint3_grafo_filtrado.html      (146 nodos, grado >= 2)
\end{lstlisting}

\end{itemize}
\end{tcolorbox}

\paragraph{Recomendación de auditoría}
\textbf{mantener --- el entregable interactivo es valioso y está versionado; corregir el año hardcodeado en los títulos (el bug más visible) y desacoplar los conteos del docstring para evitar drift documental.}

\medskip
\hrule

\subsubsection{scripts/sprint3\_redes.py}
\textbf{Rol:} orquestacion · \textbf{Líneas:} 134

\paragraph{Descripción funcional}
Análisis de redes de co-filiación institucional: grafo donde nodo = institución y arista = investigadores con doble afiliación, a nivel global (2013-2021) y por convocatoria. Usa src/analisis/redes (construir\_pares, construir\_grafo, metricas\_grafo, tabla\_nodos, tabla\_aristas, metricas\_por\_convocatoria) sobre transformar(cargar\_consolidado()). Salidas: 3 figuras PNG en artifacts/sprint3\_redes/ (distribución de grado, evolución de aristas/nodos, top-20 instituciones por grado ponderado) y 4 CSVs en evidencias/ (pares, nodos, aristas, métricas por convocatoria).

\paragraph{Análisis de la lógica}
main() extrae los pares de doble afiliación (solo INST\_FILIA con '|'), construye el grafo global con peso = investigadores compartidos, imprime métricas (nodos, aristas, densidad, componentes, hub), genera 3 figuras con funciones dedicadas y exporta 4 CSVs. metricas\_por\_convocatoria reutiliza construir\_grafo(pares, anio) para la evolución anual. La normalización de instituciones (colapsar sedes/seccionales) ocurre en src/analisis/redes.normalizar\_institucion.

\paragraph{Fortalezas}
\begin{itemize}
  \item Cobertura completa del análisis de red: grafo global, evolución por convocatoria, top instituciones y exportación tabular de nodos/aristas para reproducibilidad.
  \item El uso de grado ponderado (weight) en nodos y aristas refleja correctamente el conteo de investigadores compartidos.
  \item La normalización de instituciones (sedes, paréntesis, mayúsculas) está delegada y documentada en la capa analitica, evitando duplicados espurios.
  \item Salidas verificadas: las 3 figuras y los 4 CSVs existen en el clon en las rutas declaradas.
\end{itemize}

\paragraph{Debilidades}
\begin{itemize}
  \item Línea 102: 'instituciones mencionadas' = nunique(inst\_a) + nunique(inst\_b) cuenta dos veces toda institución que aparece en ambas columnas; la cifra impresa está sobreestimada (debería ser la unión).
  \item fig\_distribucion\_grado dibuja una barra por nodo sin etiquetas de eje X legibles: con cientos de nodos la figura es poco informativa; un histograma log-log del grado sería más analítico.
  \item No reporta métricas de comunidad (modularidad, componentes gigantes) pese a imprimir n\_componentes; el análisis de 'poder institucional' queda en grado/betweenness solamente.
  \item Depende de la columna INST\_FILIA con separador '|': si Transformacion cambia el formato (p. ej. ';'), construir\_pares devuelve un DataFrame vacío sin advertencia.
\end{itemize}

\paragraph{Mejoras propuestas}
\begin{itemize}
  \item Corregir el conteo de instituciones con pd.unique(concat(inst\_a, inst\_b)) o una unión de conjuntos.
  \item Sustituir las barras por nodo por un histograma de distribución de grado (escala log-log) y anotar el grado medio.
  \item Añadir detección de comunidad (nx.community.louvain\_communities) y tamaño del componente gigante al resumen.
  \item Validar en construir\_pares que INST\_FILIA contenga '|' y advertir si el número de pares es 0 o anómalo.
\end{itemize}

\paragraph{Ejecutabilidad}
Ejecutable en el clon solo tras instalar dependencias (pandas, matplotlib, seaborn, networkx; .venv del clon vacío, requirements.txt no los lista pero pyproject sí). El input existe: datos/tarea\_join/investigadores\_consolidado.xlsx vía cargar\_consolidado(). No requiere duckdb. Con poetry install, corre de punta a punta.

\paragraph{¿Produce lo esperado?}
SÍ, coincidencia total: existen artifacts/sprint3\_redes/fig01\_distribucion\_grado.png, fig02\_evolucion\_aristas.png, fig03\_top\_instituciones.png y evidencias/redes\_pares\_cofiliacion.csv, redes\_nodos\_global.csv, redes\_aristas\_global.csv, redes\_metricas\_por\_convocatoria.csv, con los nombres exactos del docstring y del README; confirma ejecución previa exitosa.

\begin{tcolorbox}[riesgo]
\textbf{Bugs / problemas detectados:}
\begin{itemize}
  \item scripts/sprint3\_redes.py:102 --- pares['inst\_a'].nunique() + pares['inst\_b'].nunique() sobreestima 'instituciones mencionadas' al sumar cardinalidades de columnas que comparten elementos; la cifra impresa en consola no es el número real de instituciones.
\begin{quote}\footnotesize\ttfamily
\noindent\emph{Evidencia en scripts/sprint3\_redes.py (línea 102):}
\end{quote}
\begin{lstlisting}
    print(f"      {pares['inst_a'].nunique() + pares['inst_b'].nunique()} instituciones mencionadas")
\end{lstlisting}

  \item scripts/sprint3\_redes.py:48-57 --- fig\_distribucion\_grado genera una barra por nodo (len(grados) barras) sin eje X útil: con cientos de instituciones la figura se vuelve ilegible y aporta poco frente a un histograma de la distribución de grado.
\begin{quote}\footnotesize\ttfamily
\noindent\emph{Evidencia en scripts/sprint3\_redes.py (líneas 48-57):}
\end{quote}
\begin{lstlisting}
def fig_distribucion_grado(G) -> None:
    grados = sorted([d for _, d in G.degree()], reverse=True)
    fig, ax = plt.subplots(figsize=(10, 4))
    ax.bar(range(len(grados)), grados, color="steelblue", width=1.0)
    ax.set_xlabel("Nodo (ordenado por grado)")
    ax.set_ylabel("Grado")
[...]
\end{lstlisting}

  \item scripts/sprint3\_redes.py:100 --- construir\_pares(df) asume implícitamente el separador '|' en INST\_FILIA (src/analisis/redes.py:61-64); si la transformación cambia el formato, el grafo queda vacío sin ningún warning.
\begin{quote}\footnotesize\ttfamily
\noindent\emph{Evidencia en scripts/sprint3\_redes.py (línea 100):}
\end{quote}
\begin{lstlisting}
    pares = construir_pares(df)
\end{lstlisting}

\begin{quote}\footnotesize\ttfamily
\noindent\emph{Evidencia en src/analisis/redes.py (líneas 61-64):}
\end{quote}
\begin{lstlisting}
    multi = df[df["INST_FILIA"].str.contains("|", na=False, regex=False)].copy()
    multi[["inst_a", "inst_b"]] = (
        multi["INST_FILIA"]
        .str.split("|", n=1, expand=True)
\end{lstlisting}

\end{itemize}
\end{tcolorbox}

\paragraph{Recomendación de auditoría}
\textbf{mantener --- es el análisis de redes estándar del observatorio con salidas verificadas y evidencia completa; corregir el doble conteo impreso y mejorar la figura de distribución de grado.}

\medskip
\hrule

\subsubsection{scripts/sprint4\_diversidad.py}
\textbf{Rol:} orquestacion · \textbf{Líneas:} 172

\paragraph{Descripción funcional}
Analiza las variables de diversidad (víctima de conflicto, grupo étnico, discapacidad) comparando a los investigadores reconocidos en la convocatoria 2021 con proporciones poblacionales de referencia DANE 2018 / RUV 2021. Usa src/analisis/diversidad (cobertura\_por\_convocatoria, distribucion\_categoria, comparar\_dane, interseccional\_genero\_etnia, categoria\_por\_minoria) sobre transformar(cargar\_consolidado()). Salidas: 3 figuras PNG en artifacts/sprint4\_diversidad/ y 7 CSVs en evidencias/ (uno más que los 6 del docstring).

\paragraph{Análisis de la lógica}
main() calcula la cobertura temporal de las 3 variables (solo se capturan desde 2021, hallazgo crítico documentado en el docstring), filtra df\_2021 = df[ANO\_CONVO\_INT == 2021], imprime las distribuciones excluyendo NO DISPONIBLE/NO REGISTRA, ejecuta comparar\_dane (subrepresentación de minorías vs DANE/RUV con razón pct\_dane/pct\_minc) y genera 3 figuras + exporta 7 CSVs. fig\_categorias\_por\_minoria apila la distribución de categorías académicas por grupo étnico minoritario.

\paragraph{Fortalezas}
\begin{itemize}
  \item El hallazgo metodológico central está bien documentado: las variables de diversidad solo existen desde la convocatoria 894 (2021), por lo que todo el análisis se acota honestamente a 2021.
  \item La comparación con DANE 2018 / RUV 2021 con razón de subrepresentación es metodológicamente apropiada y da un resultado de política pública accionable.
  \item Exporta evidencia completa: cobertura, distribuciones por variable, comparación DANE, cruce interseccional género$\times$etnia y categorías por etnia.
  \item Salidas verificadas: las 3 figuras y los 7 CSVs existen en el clon.
\end{itemize}

\paragraph{Debilidades}
\begin{itemize}
  \item Drift docstring-código: el docstring declara 6 CSVs pero el código escribe 7 (agrega diversidad\_distribucion\_victimas\_2021.csv en la línea 158, no listado).
  \item Bug metodológico heredado de src/analisis/diversidad.comparar\_dane (líneas 78-80): el comentario afirma 'excluir NINGUN GRUPO ETNICO del denominador para ver minorias', pero el código solo excluye NO DISPONIBLE; quienes respondieron 'NINGUN GRUPO ETNICO' permanecen en el denominador, diluyendo los porcentajes de minorías y la razón de subrepresentación.
  \item Categorías académicas hardcodeadas con tildes ('INVESTIGADOR SÉNIOR', 'INVESTIGADOR EMÉRITO', línea 105): si la normalización de NME\_CLASIFICACION\_PR usa otro casing/acentuación, la figura puede salir vacía (el filtro cats\_disponibles lo mitiga pero no avisa).
  \item El filtro 'NING' (línea 103) elimina cualquier valor que contenga esas tres letras (p. ej. un grupo 'NINGUNA' mal codificado) sin validación del vocabulario real.
  \item No hay intervalos de confianza ni tests de significancia para las diferencias vs DANE; la 'subrepresentación' se reporta sin incertidumbre.
\end{itemize}

\paragraph{Mejoras propuestas}
\begin{itemize}
  \item Actualizar el docstring para incluir diversidad\_distribucion\_victimas\_2021.csv (o eliminar esa exportación si no es entregable).
  \item Corregir comparar\_dane para excluir 'NINGUN GRUPO ETNICO' del denominador (como dice su comentario) y añadir un test que verifique denominador y suma de pct.
  \item Derivar las categorías de las columnas reales del DataFrame (p. ej. de un catálogo) en lugar de hardcodear strings con acentos.
  \item Añadir IC (binomial) a los porcentajes vs DANE y una nota sobre el tamaño muestral de 2021.
  \item Reportar n de 'NINGUN GRUPO ETNICO' y NO DISPONIBLE para que el lector vea la cobertura real del denominador.
\end{itemize}

\paragraph{Ejecutabilidad}
Ejecutable en el clon solo tras instalar dependencias (pandas, matplotlib, seaborn; .venv del clon vacío, requirements.txt incompleto). El input existe: datos/tarea\_join/investigadores\_consolidado.xlsx vía cargar\_consolidado(). Requiere que Transformacion genere ANO\_CONVO\_INT == 2021 para la convocatoria 894 (no verificado en el script). No requiere duckdb ni datos/raw. Con poetry install, corre de punta a punta.

\paragraph{¿Produce lo esperado?}
SÍ: existen artifacts/sprint4\_diversidad/fig01\_cobertura\_temporal.png, fig02\_comparacion\_dane.png, fig03\_categorias\_por\_minoria.png y los 7 CSVs de evidencias (diversidad\_cobertura\_por\_convocatoria.csv, diversidad\_distribucion\_etnia\_2021.csv, diversidad\_distribucion\_discapacidad\_2021.csv, diversidad\_distribucion\_victimas\_2021.csv, diversidad\_comparacion\_dane.csv, diversidad\_interseccional\_genero\_etnia.csv, diversidad\_categorias\_por\_etnia.csv). Coincide con el código; difiere del docstring en 1 CSV (victimas\_2021 no declarado), lo que confirma que se ejecutó una versión del script y hubo drift documental.

\begin{tcolorbox}[riesgo]
\textbf{Bugs / problemas detectados:}
\begin{itemize}
  \item src/analisis/diversidad.py:78-80 --- comentario vs código: 'excluir NINGUN GRUPO ETNICO del denominador' pero el código filtra solo 'NO DISPONIBLE' (etnia = df\_2021[TXT\_GRUPO\_ETNICO != 'NO DISPONIBLE']; n\_etnia\_resp = len(etnia) incluye NINGUN GRUPO ETNICO): los pct\_minciencias y las razones de subrepresentación de minorías quedan diluidos por el denominador inflado. Bug metodológico que afecta la figura fig02 y evidencias/diversidad\_comparacion\_dane.csv.
\begin{quote}\footnotesize\ttfamily
\noindent\emph{Evidencia en src/analisis/diversidad.py (líneas 78-80):}
\end{quote}
\begin{lstlisting}
    # Etnia: excluir "NINGUN GRUPO ETNICO" del denominador para ver minorias
    etnia = df_2021[df_2021["TXT_GRUPO_ETNICO"] != "NO DISPONIBLE"]
    n_etnia_resp = len(etnia)
\end{lstlisting}

  \item scripts/sprint4\_diversidad.py:105 --- lista hardcodeada con tildes ['INVESTIGADOR SÉNIOR', 'INVESTIGADOR EMÉRITO']: depende del casing/acentuación exacta que produzca Transformacion en NME\_CLASIFICACION\_PR; si difiere, la figura 03 muestra categorías faltantes en silencio.
\begin{quote}\footnotesize\ttfamily
\noindent\emph{Evidencia en scripts/sprint4\_diversidad.py (línea 105):}
\end{quote}
\begin{lstlisting}
    cats = ["INVESTIGADOR JUNIOR", "INVESTIGADOR ASOCIADO", "INVESTIGADOR SÉNIOR", "INVESTIGADOR EMÉRITO"]
\end{lstlisting}

  \item scripts/sprint4\_diversidad.py:103 --- cat\_etnia[...].str.contains('NING', case=False, na=False) es un filtro de subcadena: eliminaría también valores como 'NINGUNA...' si el vocabulario de la variable cambiara.
\begin{quote}\footnotesize\ttfamily
\noindent\emph{Evidencia en scripts/sprint4\_diversidad.py (línea 103):}
\end{quote}
\begin{lstlisting}
    cat_etnia = cat_etnia[~cat_etnia["TXT_GRUPO_ETNICO"].str.contains("NING", case=False, na=False)]
\end{lstlisting}

  \item scripts/sprint4\_diversidad.py:158 --- exporta diversidad\_distribucion\_victimas\_2021.csv sin declararlo en el docstring (líneas 12-22), generando evidencia no documentada.
\begin{quote}\footnotesize\ttfamily
\noindent\emph{Evidencia en scripts/sprint4\_diversidad.py (línea 158):}
\end{quote}
\begin{lstlisting}
    dist_vict.to_csv(EVIDENCIAS / "diversidad_distribucion_victimas_2021.csv", index=False)
\end{lstlisting}

\end{itemize}
\end{tcolorbox}

\paragraph{Recomendación de auditoría}
\textbf{mejorar --- el análisis es el más sensible metodológicamente (comparación con DANE/RUV) y tiene salidas verificadas, pero el denominador incorrecto en comparar\_dane compromete la cifra central de subrepresentación; corregir ese bug antes de publicar y alinear docstring-código.}

\medskip
\hrule

\subsubsection{scripts/sprint5\_duckdb.py}
\textbf{Rol:} orquestacion · \textbf{Líneas:} 193

\paragraph{Descripción funcional}
Extiende el modelo dimensional en observatorio.duckdb (creado por scripts/sprint2\_duckdb.py) con fact\_produccion, dim\_producto, dim\_grupo y la vista vw\_investigador\_x\_produccion. Usa src/ingesta.cargar\_produccion y src/analisis/produccion.normalizar\_produccion. Entradas: datos/raw/produccion\_grupos.csv y la DB previa. Salida: DB ampliada + reporte evidencias/duckdb\_resumen\_sprint5.txt con 5 consultas de validación.

\paragraph{Análisis de la lógica}
Carga y normaliza producción (IDs a Int64, extrae ANO\_CONVO\_INT), registra el DataFrame como tabla temporal 'prod\_raw' y crea fact\_produccion (una fila por id\_persona\_pd $\times$ id\_convocatoria $\times$ id\_producto\_pd), dim\_producto y dim\_grupo con SELECT DISTINCT. Crea la vista vw\_investigador\_x\_produccion con LEFT JOIN fact\_clasificacion fc ON f.id\_persona\_pd = fc.id\_persona\_pr AND f.id\_convocatoria = fc.id\_convocatoria, definiendo es\_reconocido = (fc.id\_persona\_pr IS NOT NULL). Luego ejecuta 5 consultas: productos por convocatoria; cobertura por año con pct\_reconocido; productividad por categoría = COUNT(*) / COUNT(DISTINCT id\_persona\_pd) unida a dim\_categoria (cat.categoria\_normalizada / orden\_normalizado); top 10 grupos por n\_productos; y tipologías dominantes con ventana SUM(COUNT(*)) OVER () para el porcentaje. El reporte se escribe concatenando los to\_string() de cada DataFrame.

\paragraph{Fortalezas}
\begin{itemize}
  \item Modelado dimensional limpio: separa hechos (fact\_produccion) de dimensiones (dim\_producto, dim\_grupo) y expone una vista lista para análisis cruzados, reutilizable por otros scripts o la app.
  \item Idempotente: DROP TABLE/VIEW IF EXISTS antes de cada CREATE, permitiendo re-ejecución sin corrupción acumulada.
  \item Validación integrada al pipeline: imprime conteos de filas por tabla y ejecuta consultas de negocio que verifican el resultado en el mismo run.
  \item Los artefactos declarados existen en el clon y el reporte reproduce exactamente las métricas esperadas (ver produce\_esperado).
\end{itemize}

\paragraph{Debilidades}
\begin{itemize}
  \item La vista asume que el id\_convocatoria del dataset de producción es el mismo dominio que el de fact\_clasificacion; si no coinciden (p. ej. tipos str vs int o convocatorias de grupos vs de clasificación), el LEFT JOIN degrada silenciosamente a es\_reconocido = False sin ningún aviso.
  \item No reporta la cobertura del join (cuántos productos quedan sin clasificación asociada), clave para interpretar el pct\_reconocido.
  \item La definición de 'reconocido' aquí es 'reconocido en esa misma convocatoria', distinta de la usada en sprint5\_produccion.py (cualquier convocatoria), lo que produce cifras de cobertura no comparables entre scripts del mismo sprint.
  \item No hay validación de tipos entre id\_persona\_pd (Int64) y id\_persona\_pr de fact\_clasificacion (depende de sprint2); un desajuste de tipo en DuckDB rompe el join en runtime, no en validación.
  \item La modificación de la DB no es transaccional: si falla a mitad del script, la DB queda con un subconjunto de objetos nuevos.
\end{itemize}

\paragraph{Mejoras propuestas}
\begin{itemize}
  \item Agregar validación post-join: COUNT productos sin match en fact\_clasificacion y advertir si supera un umbral.
  \item Añadir CAST/verificación explícita de tipos de id\_persona e id\_convocatoria antes de construir la vista.
  \item Escribir un test de humo que verifique n\_investigadores $\approx$ 77.237 y que prom\_productos por categoría esté en el rango esperado (Junior \textasciitilde{}30, Asociado \textasciitilde{}65, Senior \textasciitilde{}121).
  \item Parametrizar DB\_PATH y permitir --solo-validacion para no reconstruir tablas.
  \item Documentar el supuesto de que fact\_produccion.id\_convocatoria es el id de la convocatoria de clasificación del grupo.
\end{itemize}

\paragraph{Ejecutabilidad}
NO ejecutable tal cual con el clon: falta datos/processed/observatorio.duckdb (solo hay .gitkeep; requiere ejecutar scripts/sprint2\_duckdb.py antes, que a su vez necesita el raw de investigadores) y datos/raw/produccion\_grupos.csv (gitignored, \textasciitilde{}1.2 GB). Requiere la dependencia duckdb instalada. El resto (pathlib, sys) es estándar. Con el dataset y la DB presentes, el script corre de punta a punta.

\paragraph{¿Produce lo esperado?}
SÍ: evidencias/duckdb\_resumen\_sprint5.txt existe y contiene las 5 consultas. La consulta productividad\_por\_categoria reproduce Junior 29.99, Asociado 65.36, Senior 121.49, Emérito 33.75 (redondeando: 30/65/121), que es exactamente la métrica que los investigadores buscaban. La cobertura por año va de 51.88\% (2013) a 63.34\% (2021). No hay discrepancias entre lo declarado y lo generado.

\begin{tcolorbox}[riesgo]
\textbf{Bugs / problemas detectados:}
\begin{itemize}
  \item scripts/sprint5\_duckdb.py:155-157 --- el LEFT JOIN a fact\_clasificacion no valida que f.id\_convocatoria = fc.id\_convocatoria tenga coincidencias; si los dominios de convocatoria difieren (producción vs clasificación), toda la columna es\_reconocido queda en False y la cobertura reportada (51-63\%) sería espuria sin ningún error visible.
\begin{quote}\footnotesize\ttfamily
\noindent\emph{Evidencia en scripts/sprint5\_duckdb.py (líneas 155-157):}
\end{quote}
\begin{lstlisting}
        LEFT JOIN fact_clasificacion fc
               ON f.id_persona_pd = fc.id_persona_pr
              AND f.id_convocatoria = fc.id_convocatoria
\end{lstlisting}

  \item scripts/sprint5\_duckdb.py:59 --- COUNT(*)::DOUBLE es sintaxis válida en DuckDB, pero el SUM(CASE WHEN es\_reconocido THEN 1 ELSE 0 END) de la consulta cobertura\_reconocimiento devuelve tipo DOUBLE (visible como '123683.0' en el reporte) frente a COUNT(*) entero; inconsistencia de tipos menor en el output.
\begin{quote}\footnotesize\ttfamily
\noindent\emph{Evidencia en scripts/sprint5\_duckdb.py (línea 59):}
\end{quote}
\begin{lstlisting}
               ROUND(COUNT(*)::DOUBLE / COUNT(DISTINCT v.id_persona_pd), 2) AS prom_productos
\end{lstlisting}

  \item scripts/sprint5\_duckdb.py:162 --- valida COUNT(*) de fact\_produccion/dim\_producto/dim\_grupo pero nunca valida la vista vw\_investigador\_x\_produccion (el objeto más usado después).
\begin{quote}\footnotesize\ttfamily
\noindent\emph{Evidencia en scripts/sprint5\_duckdb.py (línea 162):}
\end{quote}
\begin{lstlisting}
        n = con.execute(f"SELECT COUNT(*) FROM {tabla}").fetchone()[0]
\end{lstlisting}

  \item scripts/sprint5\_duckdb.py:96-136 --- DROP TABLE IF EXISTS sin transacción; una excepción entre DROPs deja la DB parcialmente actualizada (riesgo de estado inconsistente).
\begin{quote}\footnotesize\ttfamily
\noindent\emph{Evidencia en scripts/sprint5\_duckdb.py (líneas 96-136):}
\end{quote}
\begin{lstlisting}
    con.execute("DROP TABLE IF EXISTS fact_produccion")
    con.execute("""
        CREATE TABLE fact_produccion AS
        SELECT
            id_persona_pd,
            id_convocatoria,
[...]
\end{lstlisting}

\end{itemize}
\end{tcolorbox}

\paragraph{Recomendación de auditoría}
\textbf{mantener --- es el núcleo del modelo DuckDB del sprint 5, bien estructurado y con resultados verificados; requiere solo las mejoras de validación de join y tipos para blindar la cobertura reportada.}

\medskip
\hrule

\subsubsection{scripts/sprint5\_produccion.py}
\textbf{Rol:} orquestacion · \textbf{Líneas:} 251

\paragraph{Descripción funcional}
Análisis integral de producción de grupos: cobertura de reconocimiento (autores únicos y por convocatoria), productividad por categoría, brecha de género $\times$ gran área OCDE, concentración territorial, mix de tipologías y reconocidos sin producción. Usa src/ingesta (cargar\_consolidado, cargar\_produccion), src/Transformacion.transformar y 10 funciones de src/analisis/produccion. Entradas: investigadores consolidados + producción Socrata; salidas: 6 figuras PNG y 8 evidencias (7 CSVs + 1 JSON).

\paragraph{Análisis de la lógica}
Flujo 1$\rightarrow$8: (1) carga y transforma; (2) cobertura\_autores\_unicos devuelve dict de solape de IDs; (3) cobertura\_reconocidos agrupa por ANO\_CONVO\_INT marcando es\_reconocido = ID\_PERSONA\_PD ∈ set(ID\_PERSONA\_PR del padrón EN CUALQUIER convocatoria); (4) productividad\_por\_categoria agrupa producción por (persona, convocatoria), hace left merge con el padrón y agrega por (año, NME\_CLASIFICACION\_PR) media/mediana/pct con $\geq$1 producto; (5) brecha\_productividad\_genero pivotea promedios F/M por área OCDE; (6) productividad\_territorial calcula \% productos vs \% investigadores y ratio; (7) mix\_tipologias y tipologias\_por\_categoria (porcentajes por año/categoría) y reconocidos\_sin\_produccion (intersección de sets por convocatoria); (8) genera 6 figuras matplotlib/seaborn y exporta CSVs y JSON.

\paragraph{Fortalezas}
\begin{itemize}
  \item Orquestación clara y progresiva [1/8]\dots{}[8/8] con impresión de resultados intermedios (auditable en consola).
  \item Separación limpia entre lógica de análisis (src/analisis/produccion.py) y presentación (figuras), facilitando testing y reuso.
  \item Salidas exhaustivas y documentadas en el docstring; las 6 figuras y las 8 evidencias existen en el clon.
  \item Incluye indicador de calidad de captura (reconocidos sin producción) y métricas robustas (mediana además de media).
\end{itemize}

\paragraph{Debilidades}
\begin{itemize}
  \item Definición de 'reconocido' inconsistente con sprint5\_duckdb.py: cobertura\_reconocidos marca como reconocido a quien lo fue en cualquier convocatoria (src/analisis/produccion.py:63-65), mientras la vista DuckDB exige la misma convocatoria; las coberturas resultantes (63.34\% DuckDB vs. CSV) no son comparables.
  \item Carga \textasciitilde{}3.16M filas de producción en pandas y hace merges/grupos completos en memoria: alto consumo de RAM y sin streaming ni particionado.
  \item normalizar\_produccion convierte IDs a numérico Int64 (src/analisis/produccion.py:38-40): cualquier id con ceros a la izquierda se pierde, riesgo silencioso de cruces falsos si el padrón conserva el formato string.
  \item El pct de la figura 1 se dibuja a total+5000 unidades fijas, escala que no se adapta a convocatorias pequeñas.
  \item No hay validación de que el join producción$\leftrightarrow$padrón haya encontrado coincidencias (una caída de cobertura se reporta como hallazgo, no como error).
\end{itemize}

\paragraph{Mejoras propuestas}
\begin{itemize}
  \item Unificar la definición de 'reconocido' entre scripts (recomendado: misma convocatoria) y documentarla en un único lugar.
  \item Validar tipos y solape de IDs antes de los merges (advertir si el match < umbral).
  \item Reducir memoria: leer solo columnas necesarias de producción o procesar por convocatoria.
  \item Añadir tests sobre produccion\_productividad\_por\_categoria.csv (p. ej. productos\_promedio Senior > Junior en todas las convocatorias).
  \item Hacer la posición de la etiqueta de \% proporcional al total (ax.text con offset relativo).
\end{itemize}

\paragraph{Ejecutabilidad}
NO ejecutable tal cual con el clon: cargar\_produccion() lanza FileNotFoundError porque datos/raw/produccion\_grupos.csv está gitignored y ausente. cargar\_consolidado() sí funciona (cae a datos/tarea\_join/investigadores\_consolidado.xlsx). Dependencias: pandas, numpy, matplotlib, seaborn (todas estándar del proyecto). Con el CSV de producción en datos/raw, el script corre completo.

\paragraph{¿Produce lo esperado?}
SÍ: las 6 figuras de artifacts/sprint5\_produccion/ y las 8 evidencias (produccion\_cobertura\_por\_convocatoria.csv, produccion\_productividad\_por\_categoria.csv, produccion\_brecha\_genero.csv, produccion\_concentracion\_territorial.csv, produccion\_mix\_tipologias.csv, produccion\_tipologias\_por\_categoria.csv, produccion\_reconocidos\_sin\_produccion.csv, produccion\_resumen\_cobertura.json) existen. El CSV de productividad muestra Junior 9.9-17.9 y Senior 39.0-50.1 productos promedio según convocatoria, coherente con el hallazgo 'Senior produce más'. Nota: la cifra 30/65/121 mencionada en el encargo corresponde a la métrica acumulada de sprint5\_duckdb.py, no a este CSV (que es por convocatoria).

\begin{tcolorbox}[riesgo]
\textbf{Bugs / problemas detectados:}
\begin{itemize}
  \item src/analisis/produccion.py:63-65 (invocado desde sprint5\_produccion.py:203) --- cobertura\_reconocidos marca es\_reconocido sin condicionar la convocatoria: un producto de 2013 firmado por alguien reconocido solo en 2021 cuenta como 'firmado por reconocido', sobreestimando la cobertura frente a la definición por convocatoria del DuckDB (documentado en el docstring de la función como decisión, pero no advertido en el script orquestador).
\begin{quote}\footnotesize\ttfamily
\noindent\emph{Evidencia en src/analisis/produccion.py (líneas 63-65):}
\end{quote}
\begin{lstlisting}
    ids_reconocidos = set(inv["ID_PERSONA_PR"].dropna().astype("Int64"))
    p = prod.copy()
    p["es_reconocido"] = p["ID_PERSONA_PD"].isin(ids_reconocidos)
\end{lstlisting}

\begin{quote}\footnotesize\ttfamily
\noindent\emph{Evidencia en sprint5\_produccion.py (línea 203):}
\end{quote}
\begin{lstlisting}
    cov = cobertura_reconocidos(prod, inv)
\end{lstlisting}

  \item scripts/sprint5\_produccion.py:75-78 --- la etiqueta de porcentaje se posiciona en total+5000 y el bucle itera zip(n\_reconocidos, n\_productos, pct\_reconocidos): si n\_productos < 5000 la etiqueta queda fuera del rango visible del eje (cosmético).
\begin{quote}\footnotesize\ttfamily
\noindent\emph{Evidencia en scripts/sprint5\_produccion.py (líneas 75-78):}
\end{quote}
\begin{lstlisting}
    for i, (rec, total, pct) in enumerate(
        zip(cob_df["n_reconocidos"], cob_df["n_productos"], cob_df["pct_reconocidos"])
    ):
        ax.text(i, total + 5000, f"{pct:.1f}%", ha="center", fontsize=9)
\end{lstlisting}

  \item scripts/sprint5\_produccion.py:243-244 --- import json dentro de main() (estilo menor; debería ir al tope con los demás imports).
\begin{quote}\footnotesize\ttfamily
\noindent\emph{Evidencia en scripts/sprint5\_produccion.py (líneas 243-244):}
\end{quote}
\begin{lstlisting}
    with open(EVIDENCIAS / "produccion_resumen_cobertura.json", "w", encoding="utf-8") as f:
        json.dump(cov_unicos, f, ensure_ascii=False, indent=2)
\end{lstlisting}

  \item src/analisis/produccion.py:38-40 --- pd.to\_numeric(errors='coerce') convierte IDs no numéricos a NA silenciosamente; si ID\_PERSONA\_PD tiene formatos mixtos, filas completas se pierden del análisis sin reporte.
\begin{quote}\footnotesize\ttfamily
\noindent\emph{Evidencia en src/analisis/produccion.py (líneas 38-40):}
\end{quote}
\begin{lstlisting}
    for c in ("ID_PERSONA_PD", "ID_CONVOCATORIA"):
        if c in df.columns:
            df[c] = pd.to_numeric(df[c], errors="coerce").astype("Int64")
\end{lstlisting}

\end{itemize}
\end{tcolorbox}

\paragraph{Recomendación de auditoría}
\textbf{mantener --- es el script central del sprint 5, con salidas verificadas y documentadas; las mejoras prioritarias son unificar la definición de 'reconocido' entre scripts y validar la calidad del join antes de interpretar coberturas.}

\medskip
\hrule

\subsubsection{scripts/sprint6\_geografia\_institucional.py}
\textbf{Rol:} orquestacion · \textbf{Líneas:} 302

\paragraph{Descripción funcional}
Análisis geográfico institucional: cruza el departamento de residencia del investigador con el departamento de la sede de su institución principal (vía el mapping de la tabla maestra IES), respondiendo la pregunta del director sobre dónde se forman/afilian los investigadores de departamentos como Vichada (¿en Meta o Bogotá?). Usa src/ingesta.cargar\_consolidado, src/Transformacion.transformar y evidencias/mapping\_inst\_filia\_to\_ies.csv. Salidas: 3 figuras y 3 CSVs.

\paragraph{Análisis de la lógica}
asignar\_institucion\_principal(): toma la primera institución de INST\_FILIA (split por '|'), la normaliza (upper/strip) y hace lookup en el mapping por inst\_filia\_str $\rightarrow$ (nombre\_canonico, sigla, departamento\_sede, naturaleza); normaliza departamentos con NORM\_DPTO + eliminación de acentos (unicodedata); genera DPTO\_INSTITUCION y DPTO\_RESIDENCIA canónicos. tabla\_flujos(): nunique de ID\_PERSONA\_PR por par (residencia, institución), excluyendo NO\_MAPEADA/SIN\_FILIA. resumen\_por\_residencia(): por departamento, total/local/fuera, \% local y destino externo principal. Figuras: heatmap top 15$\times$12 de flujos, \% local por departamento (top 25) y destino de departamentos pequeños (VICHADA, GUAINIA, etc.).

\paragraph{Fortalezas}
\begin{itemize}
  \item Responde exactamente la pregunta del director con un indicador accionable (\% local vs. destino externo principal) y foco explícito en departamentos pequeños.
  \item Normalización consistente de departamentos (NORM\_DPTO + quita acentos) que unifica el formato del padrón ('BOGOTÁ, D. C.') con el de la tabla maestra ('Bogotá D.C.').
  \item Exporta el flujo completo (pares residencia$\rightarrow$institución con conteos) y no solo resúmenes, permitiendo reproducir el análisis.
  \item Artefactos verificados: 3 figuras y 3 CSVs existen; el resumen muestra resultados coherentes (BOGOTA 91.8\% local, ANTIOQUIA 89.2\%, destinos externos hacia BOGOTA).
\end{itemize}

\paragraph{Debilidades}
\begin{itemize}
  \item El lookup depende de coincidencia exacta entre INST\_FILIA (tras upper/strip, SIN quitar acentos) y las claves de mapping (mayúsculas sin acentos): un padrón con tildes ('UNIVERSIDAD TECNOLÓGICA') caería a NO\_MAPEADA silenciosamente; el script solo reporta el \% global de mapeados, no por string.
  \item Tomar la primera institución cuando INST\_FILIA trae varias (separadas por |) es una simplificación no documentada como limitación.
  \item EXTRANJERO/EXTERIOR/NO IDENTIFICADO no se excluyen de los flujos y pueden ocupar lugares en los top (heatmap/top destinos).
  \item Los tops (15$\times$12, top 25, umbral) están hardcodeados sin parámetros CLI.
\end{itemize}

\paragraph{Mejoras propuestas}
\begin{itemize}
  \item Quitar acentos (unicodedata) a INST\_FILIA antes del lookup y reportar el \% de match por grupo de error (NO\_MAPEADA con top-10 strings).
  \item Excluir EXTRANJERO/EXTERIOR/NO IDENTIFICADO de los flujos o tratarlos como categoría aparte en las figuras.
  \item Documentar la regla de 'primera institución' y evaluar asignar por frecuencia/antigüedad en vez de orden textual.
  \item Parametrizar top\_n y umbrales por CLI.
\end{itemize}

\paragraph{Ejecutabilidad}
SÍ, ejecutable tal cual con el clon: cargar\_consolidado() resuelve vía datos/tarea\_join/investigadores\_consolidado.xlsx, el mapping existe (evidencias/mapping\_inst\_filia\_to\_ies.csv) y solo requiere pandas, numpy, matplotlib y seaborn.

\paragraph{¿Produce lo esperado?}
SÍ: artifacts/sprint6\_geografia/\{fig\_heatmap\_residencia\_institucion,fig\_pct\_local\_por\_dpto,fig\_top\_destinos\_pequenos\}.png y evidencias/\{geografia\_flujo\_completo,geografia\_resumen\_por\_dpto\_residencia,geografia\_top\_flujos\}.csv existen. El resumen confirma el hallazgo buscado (Bogotá 91.8\% local, departamentos pequeños con destinos externos), y geografia\_flujo\_completo.csv alimenta el Sankey territorial.

\begin{tcolorbox}[riesgo]
\textbf{Bugs / problemas detectados:}
\begin{itemize}
  \item scripts/sprint6\_geografia\_institucional.py:81-88 --- INST\_FILIA solo se pasa a upper/strip, sin quitar acentos, mientras las claves del mapping (generadas desde MAPEO y el dump) son mayúsculas sin acentos; cualquier tilde en el dato fuente rompe el match y degrada a NO\_MAPEADA sin aviso específico.
\begin{quote}\footnotesize\ttfamily
\noindent\emph{Evidencia en scripts/sprint6\_geografia\_institucional.py (líneas 81-88):}
\end{quote}
\begin{lstlisting}
    primera = (df["INST_FILIA"]
               .fillna("")
               .astype(str)
               .str.split("|")
               .str[0]
               .str.strip()
[...]
\end{lstlisting}

  \item scripts/sprint6\_geografia\_institucional.py:88 --- el reemplazo cubre '', 'NAN' y 'NO REPORTADO', pero si el padrón usa otro marcador de ausencia (p. ej. 'NO DISPONIBLE' o 'SIN INFORMACION') quedaría como institución ficticia y se contaría en los flujos.
\begin{quote}\footnotesize\ttfamily
\noindent\emph{Evidencia en scripts/sprint6\_geografia\_institucional.py (línea 88):}
\end{quote}
\begin{lstlisting}
    df["INST_PRIMERA"] = primera.replace({"": np.nan, "NAN": np.nan, "NO REPORTADO": np.nan})
\end{lstlisting}

  \item scripts/sprint6\_geografia\_institucional.py:118-119 --- se excluyen NO\_MAPEADA/SIN\_FILIA pero no 'EXTRANJERO'/'EXTERIOR'/'NO IDENTIFICADO': los flujos internacionales/desconocidos pueden dominar los top y desdibujar el mensaje territorial.
\begin{quote}\footnotesize\ttfamily
\noindent\emph{Evidencia en scripts/sprint6\_geografia\_institucional.py (líneas 118-119):}
\end{quote}
\begin{lstlisting}
    sub = df[~df["DPTO_INSTITUCION"].isin(["NO_MAPEADA", "SIN_FILIA"])]
    sub = sub[sub["DPTO_RESIDENCIA"].notna()]
\end{lstlisting}

  \item scripts/sprint6\_geografia\_institucional.py:141-145 --- fuera.iloc[0] depende de que el grupo mantenga el orden descendente global por n\_investigadores; correcto por construcción (flujo ya ordenado), pero frágil si tabla\_flujos se reordenara.
\begin{quote}\footnotesize\ttfamily
\noindent\emph{Evidencia en scripts/sprint6\_geografia\_institucional.py (líneas 141-145):}
\end{quote}
\begin{lstlisting}
        fuera = grupo[grupo["DPTO_INSTITUCION"] != dpto]
        if len(fuera):
            top = fuera.iloc[0]
            top_destino = top["DPTO_INSTITUCION"]
            top_n = int(top["n_investigadores"])
\end{lstlisting}

\end{itemize}
\end{tcolorbox}

\paragraph{Recomendación de auditoría}
\textbf{mantener --- es el análisis estrella del sprint 6 (pregunta directa del director, evidencia verificada y reutilizable); requiere endurecer la normalización de strings y el manejo de categorías no territoriales para que la cobertura de mapeo no degrade silenciosamente.}

\medskip
\hrule

\subsubsection{scripts/sprint6\_ocde\_composicion.py}
\textbf{Rol:} orquestacion · \textbf{Líneas:} 183

\paragraph{Descripción funcional}
Responde dos preguntas del director sobre la producción nacional: qué rama del conocimiento produce más y cómo se compone por tipo de medición (nuevo conocimiento, formación, apropiación), segmentando por gran área OCDE. Cruza producción (Socrata 33dq-ab5a) con el padrón por (ID\_PERSONA\_PD = ID\_PERSONA\_PR, ID\_CONVOCATORIA) mediante inner join. Usa src/ingesta, src/Transformacion.transformar y src/analisis/produccion.normalizar\_produccion. Salidas: 3 figuras y 3 CSVs.

\paragraph{Análisis de la lógica}
cruzar(): merge inner de producción con (ID\_PERSONA\_PR$\rightarrow$ID\_PERSONA\_PD, ID\_CONVOCATORIA, NME\_GRAN\_AREA\_PR) y filtra 'NO REGISTRA'/'NO REPORTADO'. Volumen = groupby área con size(); productividad = n\_productos (filas producto-por-convocatoria) / n\_investigadores (nunique de ID\_PERSONA\_PR del padrón en esa área, todas las convocatorias); composición = groupby (área, tipo medición) con pct dentro del área. Tres figuras: barras horizontales de volumen, barras apiladas de composición con PALETA\_TIPO, y barras de productividad.

\paragraph{Fortalezas}
\begin{itemize}
  \item Responde directamente la pregunta del director con una lectura clara (volumen por área, composición por tipo de medición).
  \item Paleta de tipos de medición explícita y consistente con la presentación.
  \item Filtra áreas no informadas y excluye del denominador de investigadores a áreas sin producción (evita ratios 0).
  \item Artefactos verificados: 3 figuras y 3 CSVs (ocde\_volumen, ocde\_productividad, ocde\_composicion) existen en el clon.
\end{itemize}

\paragraph{Debilidades}
\begin{itemize}
  \item El denominador de productividad mezcla unidades: n\_productos cuenta filas producto $\times$ convocatoria (una persona puede sumar productos de 5 convocatorias), mientras n\_investigadores es nunique de personas en TODO el padrón $\rightarrow$ infla la productividad por área hasta \textasciitilde{}5$\times$ frente a un promedio por persona-convocatoria.
  \item El inner join reduce la 'producción nacional' a productos firmados por reconocidos: el título '¿Qué rama produce más en Colombia?' es engañoso si se interpreta como toda la producción del país.
  \item No reporta cuántos productos quedan fuera del cruce (autores no reconocidos o convocatorias sin match), pese a imprimir solo len(df) del resultado.
  \item La productividad se calcula sin controlar categoría ni año: un área con muchos Senior en 2021 (convocatoria más productiva) se beneficia del efecto composición.
\end{itemize}

\paragraph{Mejoras propuestas}
\begin{itemize}
  \item Calcular productividad como productos por (persona, convocatoria, área) y luego promediar por persona-área, o dividir por persona-convocatoria únicos.
  \item Reportar la pérdida del inner join (n productos sin match / por qué).
  \item Renombrar la métrica a 'producción de investigadores reconocidos por área' y añadir caveat en el título.
  \item Controlar por año y categoría (p. ej. productividad estandarizada) antes de afirmar que un área produce más.
\end{itemize}

\paragraph{Ejecutabilidad}
NO ejecutable tal cual con el clon: cargar\_produccion() falla por ausencia de datos/raw/produccion\_grupos.csv (gitignored). Requiere pandas, matplotlib y seaborn. Con el dataset de producción presente, el resto del pipeline (padrón vía datos/tarea\_join) está disponible.

\paragraph{¿Produce lo esperado?}
SÍ: artifacts/sprint6\_ocde/\{fig\_volumen\_por\_area,fig\_composicion\_tipo\_por\_area,fig\_productividad\_por\_area\}.png y evidencias/\{ocde\_volumen\_por\_area,ocde\_productividad\_por\_area,ocde\_composicion\_por\_area\}.csv existen y son coherentes con la estructura del script.

\begin{tcolorbox}[riesgo]
\textbf{Bugs / problemas detectados:}
\begin{itemize}
  \item scripts/sprint6\_ocde\_composicion.py:151-154 --- n\_productos/n\_investigadores mezcla granularidades: numerador filas (persona$\times$convocatoria$\times$producto) y denominador personas únicas en todo el periodo; si una persona aparece en varias convocatorias, su producción completa se divide por 1, inflando sistemáticamente la productividad de todas las áreas.
\begin{quote}\footnotesize\ttfamily
\noindent\emph{Evidencia en scripts/sprint6\_ocde\_composicion.py (líneas 151-154):}
\end{quote}
\begin{lstlisting}
    prod_area = volumen.merge(inv_por_area, on="gran_area")
    prod_area["productos_por_investigador"] = (
        prod_area["n_productos"] / prod_area["n_investigadores"]
    ).round(2)
\end{lstlisting}

  \item scripts/sprint6\_ocde\_composicion.py:62 --- how='inner' descarta en silencio los productos sin autor reconocido o con ID\_CONVOCATORIA no coincidente; el print de [2/5] solo muestra len(df) del cruce y no la pérdida frente a len(prod).
\begin{quote}\footnotesize\ttfamily
\noindent\emph{Evidencia en scripts/sprint6\_ocde\_composicion.py (línea 62):}
\end{quote}
\begin{lstlisting}
    df = prod.merge(inv_slim, on=["ID_PERSONA_PD", "ID_CONVOCATORIA"], how="inner")
\end{lstlisting}

  \item scripts/sprint6\_ocde\_composicion.py:64 --- el filtro solo cubre 'NO REGISTRA'/'NO REPORTADO'; variantes como 'SIN INFORMACIÓN' o NaN en NME\_GRAN\_AREA\_PR pasarían y crearían una categoría basura en volumen/composición.
\begin{quote}\footnotesize\ttfamily
\noindent\emph{Evidencia en scripts/sprint6\_ocde\_composicion.py (línea 64):}
\end{quote}
\begin{lstlisting}
    df = df[~df["NME_GRAN_AREA_PR"].isin(["NO REGISTRA", "NO REPORTADO"])]
\end{lstlisting}

\end{itemize}
\end{tcolorbox}

\paragraph{Recomendación de auditoría}
\textbf{mejorar --- el análisis responde la pregunta del director y los entregables existen, pero la métrica de productividad por área es metodológicamente inconsistente (unidades mezcladas) y debe corregirse antes de usarse como hallazgo.}

\medskip
\hrule

\subsubsection{scripts/sprint6\_sankey\_categoria.py}
\textbf{Rol:} orquestacion · \textbf{Líneas:} 184

\paragraph{Descripción funcional}
Genera diagramas de Sankey de transición de categoría (Junior/Asociado/Senior/Emérito + 'Desaparece') para cada par de convocatorias consecutivas 2013$\rightarrow$2021 (5 pares). Usa src/ingesta.cargar\_consolidado, src/Transformacion.transformar y src/analisis/longitudinal (construir\_panel, comparar\_periodo, matriz\_transicion, ORDEN\_CATEGORIAS). Salidas: 5 HTML interactivos + 5 PNG estáticos y evidencias/sankey\_transiciones.csv con todos los flujos.

\paragraph{Análisis de la lógica}
Construye el panel longitudinal por convocatoria, obtiene los años ordenados (2013, 2014, 2015, 2017, 2019, 2021) y sus pares consecutivos. Por par: comparar\_periodo(panel, a, b) $\rightarrow$ matriz\_transicion(comp, incluir\_desaparece=True) devuelve (conteos, probs); reordena filas por ORDEN\_CATEGORIAS; nodos = categoría+año (izquierda t, derecha t+1) con colores de COLOR\_CAT; flujos = celdas > 0 con color rgba del origen; figura go.Sankey con arrangement='snap' y hovertemplate. Exporta HTML (plotly CDN) y PNG vía write\_image, y consolida todos los flujos en sankey\_transiciones.csv con clasificación Sube/Baja/Mantiene/Desaparece (\_etiqueta\_transicion) y resumen por periodo.

\paragraph{Fortalezas}
\begin{itemize}
  \item Respeta un orden canónico de categorías (ORDEN\_CATEGORIAS) para filas y columnas, evitando ordenamientos arbitrarios en el diagrama.
  \item Los flujos se exportan en formato largo (periodo, origen, destino, n), lo que permite reproducir el análisis sin el Sankey.
  \item Estado 'Desaparece' explícito: mide retención/abandono real entre convocatorias.
  \item Artefactos verificados: los 5 pares de HTML/PNG y sankey\_transiciones.csv (76 líneas) existen en el clon.
\end{itemize}

\paragraph{Debilidades}
\begin{itemize}
  \item fig.write\_image exige kaleido, dependencia no declarada en el docstring ni en requirements (si falta, el script muere en la primera iteración).
  \item Los HTML usan include\_plotlyjs='cdn': sin conexión a internet no renderizan.
  \item Cualquier categoría del panel que no esté en ORDEN\_CATEGORIAS se excluye silenciosamente de fila\_orden (línea 54) y desaparece del Sankey sin warning, con la consiguiente pérdida de investigadores del total.
  \item El título suma total\_t incluyendo los que 'Desaparecen' (correcto como total en t, pero puede confundir si se lee como flujo completo).
  \item Parámetros (top, umbral, altura) hardcodeados sin interfaz CLI.
\end{itemize}

\paragraph{Mejoras propuestas}
\begin{itemize}
  \item Declarar kaleido en requirements/README o manejar su ausencia con un try/except que solo escriba HTML.
  \item Validar que todas las categorías del panel estén en ORDEN\_CATEGORIAS y advertir/fallar si no.
  \item Escribir los PNG con un backend alternativo (p. ej. captura de HTML con playwright) o documentar el requisito kaleido.
  \item Exportar también las probabilidades de transición (matriz) como evidencia complementaria.
\end{itemize}

\paragraph{Ejecutabilidad}
SÍ en cuanto a datos desde el clon: cargar\_consolidado() resuelve vía datos/tarea\_join/investigadores\_consolidado.xlsx. Requiere plotly Y kaleido instalados; sin kaleido, write\_image (línea 146) lanza error y aborta. Con ambas librerías el script corre completo.

\paragraph{¿Produce lo esperado?}
SÍ: artifacts/sprint6\_sankey/sankey\_2013\_2014.\{html,png\} \dots{} sankey\_2019\_2021.\{html,png\} y evidencias/sankey\_transiciones.csv existen. El CSV muestra flujos coherentes (2013-2014: Junior$\rightarrow$Junior 1967, Junior$\rightarrow$Desaparece 2823; Senior$\rightarrow$Senior 591, etc.) y 5 periodos (2013-2014, 2014-2015, 2015-2017, 2017-2019, 2019-2021).

\begin{tcolorbox}[riesgo]
\textbf{Bugs / problemas detectados:}
\begin{itemize}
  \item scripts/sprint6\_sankey\_categoria.py:52-55 --- fila\_orden filtra las categorías presentes en ORDEN\_CATEGORIAS, pero si construir\_panel produce una categoría no mapeada (p. ej. variante con acento 'INVESTIGADOR SÉNIOR' vs 'Senior'), esa fila se elimina del Sankey y del conteo sin ningún mensaje; el total del título quedaría subestimado.
\begin{quote}\footnotesize\ttfamily
\noindent\emph{Evidencia en scripts/sprint6\_sankey\_categoria.py (líneas 52-55):}
\end{quote}
\begin{lstlisting}
    conteos, _ = matriz_transicion(comp, incluir_desaparece=True)
    # Reordenar filas siguiendo el orden canónico
    fila_orden = [c for c in ORDEN_CATEGORIAS if c in conteos.index]
    conteos = conteos.loc[fila_orden]
\end{lstlisting}

  \item scripts/sprint6\_sankey\_categoria.py:105 --- total\_t = conteos.sum().sum() incluye la columna 'Desaparece'; el texto del título 'investigadores reconocidos en t' es correcto, pero si fila\_orden ya excluyó categorías, total\_t no coincide con el padrón de t.
\begin{quote}\footnotesize\ttfamily
\noindent\emph{Evidencia en scripts/sprint6\_sankey\_categoria.py (línea 105):}
\end{quote}
\begin{lstlisting}
    total_t = conteos.sum().sum()
\end{lstlisting}

  \item scripts/sprint6\_sankey\_categoria.py:146-147 --- write\_image sin comprobación de kaleido: dependencia implícita que rompe la ejecución en entornos mínimos.
\begin{quote}\footnotesize\ttfamily
\noindent\emph{Evidencia en scripts/sprint6\_sankey\_categoria.py (líneas 146-147):}
\end{quote}
\begin{lstlisting}
        fig.write_image(ARTIFACTS / f"{nombre}.png",
                        width=1100, height=520, scale=2)
\end{lstlisting}

\end{itemize}
\end{tcolorbox}

\paragraph{Recomendación de auditoría}
\textbf{mantener --- produce el entregable de transiciones solicitado con evidencia verificada; la prioridad es declarar kaleido y validar el mapeo de categorías para no perder investigadores silenciosamente.}

\medskip
\hrule

\subsubsection{scripts/sprint6\_sankey\_territorial.py}
\textbf{Rol:} orquestacion · \textbf{Líneas:} 135

\paragraph{Descripción funcional}
Materializa el Sankey territorial residencia del investigador $\rightarrow$ departamento de su institución, usando como entrada directa evidencias/geografia\_flujo\_completo.csv (generado por sprint6\_geografia\_institucional.py), sin recargar el padrón. Respuesta visual a la pregunta del director sobre los flujos de departamentos pequeños hacia Meta/Bogotá. Salidas: sankey\_territorial.html (interactivo) y sankey\_territorial.png (estático).

\paragraph{Análisis de la lógica}
construir\_sankey(): selecciona top 12 departamentos de residencia por volumen total, top 10 de institución y flujos con n\_investigadores $\geq$ 15; construye nodos 'X (res.)' e 'Y (inst.)' con colores de PALETA por residencia (nodos de institución en gris); flujos con alpha 0.40 si el flujo es local (residencia == institución) y 0.65 si es transversal; figura go.Sankey con arrangement='snap' y hovertemplate con formato de miles. Exporta HTML (plotly CDN) y PNG vía write\_image a 2x.

\paragraph{Fortalezas}
\begin{itemize}
  \item Diseño de pipeline limpio: consume el CSV ya generado por el script de geografía (una sola fuente de verdad, sin duplicar la lógica de cruce).
  \item Parámetros de saturación explícitos (top\_res, top\_inst, umbral\_min) que evitan un diagrama ilegible con miles de flujos.
  \item Distinción visual local vs. transversal (alpha) y paleta estable.
  \item Artefactos verificados: sankey\_territorial.html y sankey\_territorial.png existen en el clon.
\end{itemize}

\paragraph{Debilidades}
\begin{itemize}
  \item write\_image requiere kaleido (dependencia no declarada): el PNG falla en entornos sin esa librería aunque el HTML se escriba bien.
  \item HTML con include\_plotlyjs='cdn': requiere internet para renderizar.
  \item Los tops y el umbral están fijos en main() sin interfaz CLI para explorar variantes.
  \item No exporta el subconjunto de flujos usados en el diagrama (CSV), limitando la reproducibilidad del gráfico exacto.
  \item Depende de que geografia\_flujo\_completo.csv use los strings normalizados; si ese script cambia de normalización, este se rompe sin mensaje claro.
\end{itemize}

\paragraph{Mejoras propuestas}
\begin{itemize}
  \item Declarar kaleido o envolver write\_image en try/except para degradar a solo HTML con aviso.
  \item Exportar el CSV de los flujos filtrados usados en la figura.
  \item Parametrizar top\_res/top\_inst/umbral\_min por argparse para iterar rápido.
  \item Añadir validación de columnas esperadas (DPTO\_RESIDENCIA, DPTO\_INSTITUCION, n\_investigadores) al leer el CSV.
\end{itemize}

\paragraph{Ejecutabilidad}
SÍ, ejecutable tal cual con el clon: la entrada evidencias/geografia\_flujo\_completo.csv existe; requiere plotly y kaleido instalados (sin kaleido falla solo la escritura del PNG). No depende de datos/raw ni del padrón.

\paragraph{¿Produce lo esperado?}
SÍ: artifacts/sprint6\_sankey/sankey\_territorial.html y sankey\_territorial.png existen, completando el set de 6 sankeys del sprint 6 (5 de categoría + 1 territorial). La lógica produce exactamente el flujo solicitado (residencia $\rightarrow$ institución con umbral de ruido).

\begin{tcolorbox}[riesgo]
\textbf{Bugs / problemas detectados:}
\begin{itemize}
  \item scripts/sprint6\_sankey\_territorial.py:125-129 --- write\_image sin comprobación de kaleido: dependencia implícita; en un entorno sin kaleido el script aborta tras haber generado el HTML (comportamiento no degradable).
\begin{quote}\footnotesize\ttfamily
\noindent\emph{Evidencia en scripts/sprint6\_sankey\_territorial.py (líneas 125-129):}
\end{quote}
\begin{lstlisting}
    fig.write_html(ARTIFACTS / "sankey_territorial.html",
                    include_plotlyjs="cdn",
                    full_html=True)
    fig.write_image(ARTIFACTS / "sankey_territorial.png",
                     width=1200, height=620, scale=2)
\end{lstlisting}

  \item scripts/sprint6\_sankey\_territorial.py:69-70 --- top\_residencias.index(row['DPTO\_RESIDENCIA']) asume que toda fila de sub está en la lista top; correcto por el filtro isin previo (línea 56-58), pero si top\_residencias tuviera duplicados (no los tiene, viene de nlargest), el index() devolvería la primera aparición.
\begin{quote}\footnotesize\ttfamily
\noindent\emph{Evidencia en scripts/sprint6\_sankey\_territorial.py (líneas 69-70):}
\end{quote}
\begin{lstlisting}
        s_idx = top_residencias.index(row["DPTO_RESIDENCIA"])
        t_idx = len(top_residencias) + top_instituciones.index(row["DPTO_INSTITUCION"])
\end{lstlisting}

  \item scripts/sprint6\_sankey\_territorial.py:56-58 --- el filtro por umbral\_min $\geq$ 15 es absoluto, no proporcional al total del departamento: un flujo pequeño de VICHADA (dpto de pocos investigadores) puede quedar excluido justo donde el director quiere verlo; convendría umbral relativo por origen.
\begin{quote}\footnotesize\ttfamily
\noindent\emph{Evidencia en scripts/sprint6\_sankey\_territorial.py (líneas 56-58):}
\end{quote}
\begin{lstlisting}
    sub = flujo[flujo["DPTO_RESIDENCIA"].isin(top_residencias) &
                flujo["DPTO_INSTITUCION"].isin(top_instituciones) &
                (flujo["n_investigadores"] >= umbral_min)].copy()
\end{lstlisting}

\end{itemize}
\end{tcolorbox}

\paragraph{Recomendación de auditoría}
\textbf{mantener --- script simple, enfocado y con artefactos verificados que completan el entregable del director; mejoras menores (kaleido declarado, CSV de flujos filtrados, umbral relativo para departamentos pequeños) aumentan su robustez y utilidad.}

\medskip
\hrule

\subsubsection{scripts/sprint6\_tabla\_maestra\_ies.py}
\textbf{Rol:} orquestacion · \textbf{Líneas:} 635

\paragraph{Descripción funcional}
Construye el borrador de la tabla maestra canónica de IES a partir de evidencias/instituciones\_unicas\_raw.csv (dump de strings únicos de INST\_FILIA con frecuencias), aplicando un diccionario MAPEO hardcodeado de \textasciitilde{}250 claves string $\rightarrow$ (nombre\_canonico, sigla, departamento\_sede, naturaleza). Solo usa pandas, sin tocar el padrón ni producción. Salidas: tabla\_maestra\_ies.csv, mapping\_inst\_filia\_to\_ies.csv e ies\_pendientes\_revision.csv.

\paragraph{Análisis de la lógica}
aplicar\_mapeo(row): si inst\_filia\_str ∈ MAPEO devuelve la tupla canónica con estado 'asignado'; si no, NA con estado 'pendiente\_revision'. main(): lee el dump, aplica el mapeo fila a fila (apply axis=1), calcula cobertura sobre apariciones (asignadas/total), agrupa las asignadas por (nombre\_canonico, sigla, naturaleza) con agregados (departamentos únicos, apariciones, variantes) ordenadas por apariciones desc, y exporta los 3 archivos, imprimiendo cobertura y top 20.

\paragraph{Fortalezas}
\begin{itemize}
  \item Borrador transparente y auditable: el mapeo completo es legible, con comentarios por pasadas y objetivo de cobertura (>90\%).
  \item Reporta cobertura real sobre apariciones del padrón y separa explícitamente las entradas pendientes de revisión.
  \item Cero dependencias externas más allá de pandas y ejecutable con el clon (el input instituciones\_unicas\_raw.csv existe, 2763 filas).
  \item Decisiones de negocio explícitas: una IES por sede/razón social con departamento\_sede por ubicación física.
\end{itemize}

\paragraph{Debilidades}
\begin{itemize}
  \item Mapeo 100\% manual y hardcodeado (\textasciitilde{}250 de 2763 strings): frágil ante variantes nuevas (acentos, guiones, 'SAS', 'LTDA') y costoso de mantener; es una solución de 'borrador', no un componente de datos.
  \item La columna 'sedes\_registradas' en realidad cuenta departamentos distintos (nunique de departamento\_sede), no sedes físicas: nombre engañoso.
  \item Coincidencia exacta por mayúsculas-sin-acentos: si el dump cambia de normalización (p. ej. conservara tildes), todo el mapeo se rompe en silencio y la cobertura cae sin error.
  \item Variantes duplicadas sin mapear (p. ej. 'UNIVERSIDAD NACIONAL DE COLOMBIA' sin sede) asumen Bogotá por defecto sin validación.
  \item apply(axis=1) con función Python es lento, aunque tolerable para \textasciitilde{}2.8k filas.
\end{itemize}

\paragraph{Mejoras propuestas}
\begin{itemize}
  \item Externalizar el MAPEO a un CSV/JSON versionado (fuente de datos, no código) con validación de unicidad de claves.
  \item Añadir normalización difusa o reglas (contiene 'UNIVERSIDAD NACIONAL DE COLOMBIA' $\rightarrow$ UNAL) para cubrir variantes no listadas, reportando la confianza.
  \item Renombrar 'sedes\_registradas' a 'departamentos\_sede' o contar sedes reales.
  \item Añadir un test que falle si la cobertura baja de un umbral (p. ej. 85\%) ante cambios del dump.
\end{itemize}

\paragraph{Ejecutabilidad}
SÍ, ejecutable tal cual con el clon: solo necesita pandas y evidencias/instituciones\_unicas\_raw.csv (existente, columnas inst\_filia\_str, n\_apariciones, pct, pct\_acum --- coinciden con lo que lee el script). No depende de datos/raw ni de observatorio.duckdb.

\paragraph{¿Produce lo esperado?}
SÍ: evidencias/tabla\_maestra\_ies.csv, evidencias/mapping\_inst\_filia\_to\_ies.csv e evidencias/ies\_pendientes\_revision.csv existen. El mapping tiene las columnas esperadas (inst\_filia\_str, nombre\_canonico, sigla, departamento\_sede, naturaleza, estado) y es el insumo que consumen sprint6\_geografia\_institucional.py y la app.

\begin{tcolorbox}[riesgo]
\textbf{Bugs / problemas detectados:}
\begin{itemize}
  \item scripts/sprint6\_tabla\_maestra\_ies.py:607 --- 'sedes\_registradas' = nunique de departamento\_sede: dos sedes en el mismo departamento (p. ej. dos campus de la UNAL en Bogotá) cuentan como 1, mientras dos sedes en departamentos distintos cuentan como 2; la métrica no es número de sedes.
\begin{quote}\footnotesize\ttfamily
\noindent\emph{Evidencia en scripts/sprint6\_tabla\_maestra\_ies.py (línea 607):}
\end{quote}
\begin{lstlisting}
            sedes_registradas=("departamento_sede", "nunique"),
\end{lstlisting}

  \item scripts/sprint6\_tabla\_maestra\_ies.py:38-562 --- claves hardcodeadas sin validación de duplicados ni de normalización: si instituciones\_unicas\_raw.csv se regenera con otro case/acentos, la cobertura cae sin ninguna señal de error.
\begin{quote}\footnotesize\ttfamily
\noindent\emph{Evidencia en scripts/sprint6\_tabla\_maestra\_ies.py (líneas 38-562):}
\end{quote}
\begin{lstlisting}
MAPEO = {
    # --- Top 20 por frecuencia -------------------------------------------
    "UNIVERSIDAD DE ANTIOQUIA":
        ("Universidad de Antioquia", "UdeA", "Antioquia", "publica"),
    "UNIVERSIDAD NACIONAL DE COLOMBIA SEDE BOGOTA (UNIVERSIDAD NACIONAL DE COLOMBIA)":
        ("Universidad Nacional de Colombia", "UNAL", "Bogotá D.C.", "publica"),
[...]
\end{lstlisting}

  \item scripts/sprint6\_tabla\_maestra\_ies.py:592 --- pd.concat([raw, asignaciones], axis=1) asume que aplicar\_mapeo devuelve una Series por fila alineada por índice; correcto aquí, pero frágil si se reordenara raw antes.
\begin{quote}\footnotesize\ttfamily
\noindent\emph{Evidencia en scripts/sprint6\_tabla\_maestra\_ies.py (línea 592):}
\end{quote}
\begin{lstlisting}
    asignaciones = raw.apply(aplicar_mapeo, axis=1)
\end{lstlisting}

  \item scripts/sprint6\_tabla\_maestra\_ies.py:620-623 --- ies\_pendientes\_revision.csv incluye todas las entradas sin mapeo, pero el script no imprime las 10 pendientes más frecuentes para facilitar la revisión manual (solo cobertura agregada).
\begin{quote}\footnotesize\ttfamily
\noindent\emph{Evidencia en scripts/sprint6\_tabla\_maestra\_ies.py (líneas 620-623):}
\end{quote}
\begin{lstlisting}
    pendientes = mapping[mapping["estado"] == "pendiente_revision"].copy()
    pendientes.to_csv(EVIDENCIAS / "ies_pendientes_revision.csv",
                       index=False, encoding="utf-8")
    cobertura_pendientes = pendientes["n_apariciones"].sum() / total_apariciones * 100
\end{lstlisting}

\end{itemize}
\end{tcolorbox}

\paragraph{Recomendación de auditoría}
\textbf{refactorizar --- como borrador cumple y sus productos existen, pero un catálogo canónico de IES no debería vivir en código: externalizar el mapeo a datos versionados, renombrar métricas y añadir cobertura mínima garantizada antes de que otros scripts (geografía, app) dependan de él en producción.}

\medskip
\hrule

\subsubsection{scripts/sprint6\_validacion\_calidad.py}
\textbf{Rol:} orquestacion · \textbf{Líneas:} 121

\paragraph{Descripción funcional}
Valida y documenta el tratamiento de valores atípicos de la variable EDAD\_ANOS\_PR del padrón: detecta los casos > 100 (máximo observado 956, error de digitación), compara tres estrategias (original, eliminación, imputación por mediana de grupo) y genera evidencia auditable. Usa src/ingesta.cargar\_consolidado, src/Transformacion.transformar y src/analisis/calidad (detectar\_atipicos\_edad, resumen\_tratamiento, filtrar, imputar\_mediana, UMBRAL\_EDAD\_MAX). Salidas: 2 CSVs, 1 JSON y 1 figura.

\paragraph{Análisis de la lógica}
Flujo 1$\rightarrow$4: (1) transformar(cargar\_consolidado()); (2) detectar\_atipicos\_edad devuelve los registros con EDAD\_ANOS\_PR > 100 ordenados desc; (3) resumen\_tratamiento calcula n\_total/n\_atipicos/n\_validos y justificación, y se construye un DataFrame comparativo edad\_stats con media/mediana/desv/max para original, filtrar() (eliminación) e imputar\_mediana(); (4) histogramas antes/después con umbral trazado como axvline, y exportación de atipicos a CSV, comparación a CSV y resumen a JSON.

\paragraph{Fortalezas}
\begin{itemize}
  \item Metodología explícita y auditable: umbral documentado (100), justificación cuantitativa (10/77.237 = 0.013\%) y decisión adoptada (eliminación) con alternativa de imputación comparada.
  \item La evidencia confirma los números esperados: n=77.237 registros, máx original 956.25 $\rightarrow$ 100 tras filtro (evidencias/calidad\_comparacion.csv).
  \item Script pequeño, legible y sin dependencias exóticas (pandas, matplotlib, seaborn).
  \item Reutiliza funciones de negocio de src/analisis/calidad.py con docstrings de decisión metodológica.
\end{itemize}

\paragraph{Debilidades}
\begin{itemize}
  \item n\_registros usa len(df) (cuenta filas, incluidas las de EDAD NaN) mientras media/mediana/std ignoran NaN: unidades de conteo mixtas en la misma tabla comparativa.
  \item La mediana se deja sin redondear mientras media y desv\_estandar se redondean a 2 decimales (formato inconsistente).
  \item El script solo valida edad; otros problemas documentados en calidad.py (departamento NO DISPONIBLE 5.2\%, género NO REPORTADO 2.7\%, convocatoria 833 con año 2019) se citan pero no se verifican ni exportan.
  \item La figura compara 'original' vs 'filtrado' pero no muestra el efecto de la imputación, pese a que se calcula en el DataFrame comparativo.
\end{itemize}

\paragraph{Mejoras propuestas}
\begin{itemize}
  \item Calcular n\_registros como count no-NA de EDAD\_ANOS\_PR para consistencia con las demás métricas.
  \item Redondear mediana igual que media y documentar el criterio.
  \item Ampliar la validación a duplicados, nulos por columna y distribución de tipos, exportando un perfil de calidad completo.
  \item Fijar seed/versión del dataset en el JSON de resumen para trazabilidad.
\end{itemize}

\paragraph{Ejecutabilidad}
SÍ, ejecutable tal cual con el clon: cargar\_consolidado() resuelve desde datos/tarea\_join/investigadores\_consolidado.xlsx (no requiere raw Socrata); depende solo de pandas, matplotlib y seaborn. El único requisito no verificado es que transformar() produzca EDAD\_ANOS\_PR, lo cual confirman los artefactos existentes.

\paragraph{¿Produce lo esperado?}
SÍ: evidencias/calidad\_atipicos\_edad.csv, evidencias/calidad\_comparacion.csv, evidencias/calidad\_resumen.json y artifacts/sprint6\_calidad/fig\_edad\_antes\_despues.png existen. El CSV de comparación confirma 77.237 registros, media 44.84 $\rightarrow$ 44.73, desv 13.23 $\rightarrow$ 10.04 y máx 956.25 $\rightarrow$ 100.0, coherente con la decisión de eliminación.

\begin{tcolorbox}[riesgo]
\textbf{Bugs / problemas detectados:}
\begin{itemize}
  \item scripts/sprint6\_validacion\_calidad.py:62-84 --- edad\_stats mezcla unidades: len(df) cuenta filas (incluye NaN en EDAD) mientras mean/median/std de pandas ignoran NaN; si hubiera filas con edad faltante, n\_registros y las demás columnas no serían comparables.
\begin{quote}\footnotesize\ttfamily
\noindent\emph{Evidencia en scripts/sprint6\_validacion\_calidad.py (líneas 62-84):}
\end{quote}
\begin{lstlisting}
    edad_stats = pd.DataFrame({
        "metodo": ["original", "eliminacion (>100)", "imputacion mediana por grupo"],
        "n_registros": [len(df), len(df_filtrado), len(df_imputado)],
        "media": [
            round(df["EDAD_ANOS_PR"].mean(), 2),
            round(df_filtrado["EDAD_ANOS_PR"].mean(), 2),
[...]
\end{lstlisting}

  \item scripts/sprint6\_validacion\_calidad.py:70-73 --- df['EDAD\_ANOS\_PR'].median() sin round frente a round(..., 2) en media (formato inconsistente; en el CSV aparece 43.29 por coincidencia del valor).
\begin{quote}\footnotesize\ttfamily
\noindent\emph{Evidencia en scripts/sprint6\_validacion\_calidad.py (líneas 70-73):}
\end{quote}
\begin{lstlisting}
        "mediana": [
            df["EDAD_ANOS_PR"].median(),
            df_filtrado["EDAD_ANOS_PR"].median(),
            df_imputado["EDAD_ANOS_PR"].median(),
\end{lstlisting}

  \item scripts/sprint6\_validacion\_calidad.py:50-51 --- detectar\_atipicos\_edad solo imprime, no informa el conteo de NaN de EDAD\_ANOS\_PR, que podría alterar la interpretación del umbral.
\begin{quote}\footnotesize\ttfamily
\noindent\emph{Evidencia en scripts/sprint6\_validacion\_calidad.py (líneas 50-51):}
\end{quote}
\begin{lstlisting}
    atipicos = detectar_atipicos_edad(df)
    print(atipicos.to_string(index=False))
\end{lstlisting}

\end{itemize}
\end{tcolorbox}

\paragraph{Recomendación de auditoría}
\textbf{mantener --- cumple su propósito (decisión metodológica auditable) con evidencia verificada; requiere solo ajustes menores de consistencia de conteos y redondeo.}

\medskip
\hrule

\subsection{Visualización y dashboard}
\noindent Archivos: 2.

\subsubsection{app/streamlit\_app.py}
\textbf{Rol:} visualizacion · \textbf{Líneas:} 17

\paragraph{Descripción funcional}
Placeholder de 17 lineas que solo configura la pagina (titulo 'Observatorio Ustadistica', icono de grafico, layout wide) y muestra un titulo con el aviso 'Dashboard en construccion. Consulta el README para el plan de desarrollo.' No carga datos, no importa src/ y no contiene ningun analisis. Es la entrada que ejecuta el Dockerfile (CMD streamlit run app/streamlit\_app.py).

\paragraph{Análisis de la lógica}
Estructura minima: import streamlit; st.set\_page\_config (page\_title 'Observatorio Ustadistica', page\_icon \textbackslash\{\}U0001f4ca, layout wide); st.title; st.info. Sin st.cache\_data, sin filtros, sin tabs, sin conexion con src/. El docstring documenta el uso con `poetry run streamlit run app/streamlit\_app.py`.

\paragraph{Fortalezas}
\begin{itemize}
  \item Sintaxis Streamlit valida: no crashea al arrancar
  \item Sirve como esqueleto para desarrollo futuro (README del proyecto menciona 'plan de desarrollo')
\end{itemize}

\paragraph{Debilidades}
\begin{itemize}
  \item Es un placeholder que contradice al dashboard real (streamlit\_app.py, 983 lineas, 7 tabs)
  \item El README documenta el deploy del dashboard real (main file streamlit\_app.py en Streamlit Cloud, linea 192) pero el Dockerfile despliega este placeholder: dos 'dashboards' distintos segun la via de deploy
  \item La imagen Docker no copia el dashboard real (COPY solo src/, app/, datos/) ni hallazgos/, asi que en el contenedor el usuario final ve solo el cartel de 'en construccion'
  \item Duplica la identidad del proyecto (titulo distinto al del dashboard real: 'Observatorio Ustadistica' vs 'Observatorio MinCiencias - Informe critico')
\end{itemize}

\paragraph{Mejoras propuestas}
\begin{itemize}
  \item Eliminar el placeholder y apuntar el Dockerfile al dashboard real (copiar streamlit\_app.py de la raiz, hallazgos/ y datos/), o
  \item Mover el dashboard real a app/streamlit\_app.py (haciendo app/ la entrada canonica) y actualizar README y Streamlit Cloud en consecuencia, o
  \item Si se mantiene como esqueleto, excluirlo del Dockerfile y documentar que no es el dashboard de produccion
\end{itemize}

\paragraph{Ejecutabilidad}
Sí, se ejecuta tal cual (solo requiere streamlit instalado), pero solo muestra el aviso de 'Dashboard en construccion'. No lee ningun dato ni modulo del proyecto.

\paragraph{¿Produce lo esperado?}
NO produce el dashboard de 7 capitulos documentado: renderiza un cartel estatico. En el despliegue Docker, el contenedor (python:3.12-slim + poetry + datos/) muestra este placeholder en vez del dashboard real, contradiciendo el README y el proposito del observatorio.

\begin{tcolorbox}[riesgo]
\textbf{Bugs / problemas detectados:}
\begin{itemize}
  \item Bug de deploy (no de sintaxis): el Dockerfile ejecuta `streamlit run app/streamlit\_app.py` (linea 22) -> el contenedor despliega el placeholder y no el dashboard real; ademas streamlit\_app.py de la raiz ni siquiera se copia a la imagen (COPY src/, app/, datos/ solamente)
  \item Contradiccion de documentacion: README linea 175 (`poetry run streamlit run streamlit\_app.py`) y linea 192 (Streamlit Cloud main file streamlit\_app.py) apuntan al dashboard real, mientras que el docstring de este archivo y el Dockerfile apuntan al placeholder
  \item La imagen Docker tampoco copia hallazgos/sprint3\_grafo\_filtrado.html: aunque se cambiara el CMD al dashboard real, la tab Redes fallaria en el contenedor
\end{itemize}
\end{tcolorbox}

\paragraph{Recomendación de auditoría}
\textbf{eliminar: el placeholder no aporta valor de produccion y rompe la consistencia del deploy (Docker muestra algo distinto a lo documentado y a Streamlit Cloud). La accion correcta es hacer que el Dockerfile ejecute el dashboard real (streamlit\_app.py de la raiz o movido a app/) y copiar hallazgos/; o, si se prefiere app/ como canonica, mover el dashboard real ahi y borrar este archivo.}

\medskip
\hrule

\subsubsection{streamlit\_app.py}
\textbf{Rol:} visualizacion · \textbf{Líneas:} 983

\paragraph{Descripción funcional}
Dashboard Streamlit principal del Observatorio MinCiencias con 7 pestañas tipo capítulo (Calidad, Trayectoria, Territorio, Diversidad, Redes, Productividad, Datos crudos). Lee el consolidado vía src/ingesta.cargar\_consolidado (XLSX de respaldo en datos/tarea\_join), lo transforma con src/Transformacion.transformar y lo cruza con \textasciitilde{}30 funciones de src/analisis/*. Cada sección sigue el patrón pregunta -> evidencia -> hallazgo -> caveat, con filtros globales en la barra lateral.

\paragraph{Análisis de la lógica}
Estructura: constantes (COORDENADAS\_DEPTO, POBLACION\_DANE\_2018, CAT\_CORTA) -> cargar\_datos() con @st.cache\_data (aplica transformar y crea edad\_valida) -> sidebar\_filtros() (año, género, gran área, depto top-15, categoría, checkbox excluir calidad) -> cabecera() con 6 métricas -> 7 tabs que llaman a una funcion seccion\_* cada una. Caches: cargar\_datos, panel\_completo(len(df), df) y cargar\_produccion\_normalizada. Conexion con src/: sys.path.insert(ROOT/src) e imports de ingesta, Transformacion y analisis.\{longitudinal, territorial, genero, diversidad, redes, produccion\}. La seccion 5 embebe hallazgos/sprint3\_grafo\_filtrado.html (existe) y la 6 depende de datos/raw/produccion\_grupos.csv (NO existe en el repo).

\paragraph{Fortalezas}
\begin{itemize}
  \item Arquitectura limpia por secciones con una funcion por capitulo; facil de auditar
  \item Narrativa critica consistente (pregunta/hallazgo/caveat) con caveats metodologicos muy buenos (p. ej. subrepresentacion vs poblacion con educacion superior, Eméritos vitalicios)
  \item Manejo defensivo: try/except que degrada con st.info en vez de crashear (HHI, genero, DANE, longitudinal)
  \item Reutiliza las funciones reales de src/analisis; no duplica logica de analisis en la UI
  \item Filtros globales coherentes entre secciones y metricas de cabecera bien elegidas
  \item Imports y rutas verificados: todos los modulos src/analisis importados existen con las firmas usadas
\end{itemize}

\paragraph{Debilidades}
\begin{itemize}
  \item La seccion 6 (Productividad, capitulo 6 del informe) queda vacia en esta maquina: falta datos/raw/produccion\_grupos.csv (\textasciitilde{}1,2 GB); solo muestra un warning
  \item El grafico HHI de la seccion 3 esta roto por un bug de shape y degrada silenciosamente a 'HHI no disponible'
  \item Numeros hardcodeados en texto y logica: '30.086 investigadores unicos', fallback 16796 para la anomalia de convocatoria, '6 convocatorias (2013-2021)' en la caption
  \item Caché de panel\_completo con clave debil (len(df)): filtros con el mismo tamaño colisionan y pueden devolver paneles obsoletos
  \item Docstring dice '6 secciones tipo capitulos' pero hay 7 tabs (desincronizacion interna)
  \item Seccion 6 no cubierta por try/except: si el CSV de produccion existiera vacio, cov.iloc[0]/iloc[-1] lanzarian IndexError sin captura
  \item Metricas de cabecera usan df\_full mientras las secciones usan df filtrado: el texto 'global' mezcla poblaciones
\end{itemize}

\paragraph{Mejoras propuestas}
\begin{itemize}
  \item Corregir el bug del HHI: eliminar el .reset\_index() o renombrar usando las columnas reales (ANO\_CONVO\_INT, hhi)
  \item Documentar el prerequisito del dataset de produccion en la UI (ya hay warning) o agregar fallback a evidencias/produccion\_*.csv
  \item Calcular en runtime los numeros hardcodeados (unicos, convocatorias) en vez de fijarlos
  \item Cachear panel\_completo con un hash real del filtro (tupla de selecciones) en vez de len(df)
  \item Envolver la seccion 6 completa en try/except para evitar IndexError con produccion vacia
  \item Sincronizar docstring (6 -> 7 secciones) y caption
\end{itemize}

\paragraph{Ejecutabilidad}
Parcial. `streamlit run streamlit\_app.py` desde la raiz funciona con streamlit, pandas, numpy, plotly y openpyxl (requirements.txt los declara). cargar\_consolidado cae al XLSX datos/tarea\_join/investigadores\_consolidado.xlsx (existe; datos/raw/investigadores\_consolidado.csv no existe). Con ese insumo, las secciones 1, 2, 3 (sin HHI), 4, 5 y 7 deberian renderizar. La seccion 6 (productividad) no se ejecuta: requiere descargar datos/raw/produccion\_grupos.csv con `python -m src.ingesta.produccion`; la app lo maneja con st.warning y return. En entorno poetry (Docker) openpyxl figura en poetry.lock (grupo main) aunque pyproject.toml no lo declara: riesgo de lock desincronizado.

\paragraph{¿Produce lo esperado?}
Sí, el dashboard implementa los 7 capitulos del informe y coincide con el README ('streamlit\_app.py \# Dashboard - 7 tabs'). Sin embargo, el capitulo 6 (productividad) no produce nada sin el dataset de produccion, y el HHI (parte del capitulo 3) nunca dibuja por el bug de shape. El placeholder app/streamlit\_app.py no contradice a este archivo en local ni en Streamlit Cloud (el README apunta a este), pero sí en el despliegue Docker.

\begin{tcolorbox}[riesgo]
\textbf{Bugs / problemas detectados:}
\begin{itemize}
  \item Seccion 3: hhi\_por\_convocatoria(df, col\_geo=col).reset\_index() devuelve 5 columnas (index + ANO\_CONVO\_INT + hhi + n\_total + n\_territorios) y luego hhi\_df.columns = ['Año', 'HHI'] lanza ValueError 'Length mismatch'; el except lo captura y el grafico HHI nunca se muestra (lineas 434-435)
  \item Seccion 6: depende de datos/raw/produccion\_grupos.csv que no esta en el repo; toda la seccion (4 metricas + 6 graficos) se degrada a un warning
  \item Seccion 1: fallback hardcodeado `else 16796` y umbral magico ID\_CONVOCATORIA == 20 sin validacion -> numero potencialmente incorrecto en silencio
  \item panel\_completo(len(df), df) usa len(df) como hash de cache -> colisiones entre filtros del mismo tamano (linea 347)
  \item Si el CSV de produccion existiera pero vacio, cov.iloc[0]/cov.iloc[-1] (lineas 761-762) lanzarian IndexError no capturado y crashearia la app
  \item Caption y docstring desincronizados: '6 secciones' en el header (linea 7) vs 7 tabs reales
\end{itemize}
\end{tcolorbox}

\paragraph{Recomendación de auditoría}
\textbf{mejorar: es la pieza central del proyecto, bien estructurada y con 5 de 7 secciones funcionales con el XLSX actual. Corregir el bug del HHI, robustecer la cache, documentar/proveer el dataset de produccion y eliminar los hardcodes antes de considerar un refactor.}

\medskip
\hrule

\subsection{Documentación ejecutable}
\noindent Archivos: 1.

\subsubsection{docs/manual.ipynb}
\textbf{Rol:} documentacion · \textbf{Líneas:} 805

\paragraph{Descripción funcional}
Manual reproducible de 49 celdas (23 markdown + 26 codigo) generado por scripts/generar\_manual.py (no ejecutado: todas las celdas tienen execution\_count null). 14 secciones que reproducen los hallazgos del observatorio: calidad (atipicos de edad), trayectoria y matrices de transicion, Sankeys de categoria, concentracion territorial (HHI), tabla maestra de IES, geografia institucional (residencia vs institucion), brecha de genero, diversidad vs DANE/RUV, productividad por categoria, brecha de genero en productividad, composicion OCDE y recomendaciones de politica publica.

\paragraph{Análisis de la lógica}
Celda 1: ROOT = pathlib.Path.cwd() con fallback a ROOT.parent si (ROOT/src) no existe (funciona desde la raiz del repo o desde docs/); inserta src/ en sys.path. Celda 2: inv = transformar(cargar\_consolidado()) y prod = normalizar\_produccion(cargar\_produccion()). Secciones 2-3: funciones de src/analisis/calidad.py y longitudinal.py (construir\_panel, comparar\_periodo, matriz\_transicion, matrices\_todos\_periodos). Secciones 4, 7, 8 y 13: subprocess.run(['python3', 'scripts/sprint6\_*.py'], cwd=ROOT). Secciones 6-8 y 13: pd.read\_csv de evidencias/*.csv. IFrame apunta a artifacts/sprint6\_sankey/*.html. Todo el codigo del manual reutiliza los modulos reales de src/ (no duplica logica).

\paragraph{Fortalezas}
\begin{itemize}
  \item Intencion didactica y reproducible clara: cada hallazgo con su codigo minimo y su caveat
  \item Reutiliza el codigo real (src/ingesta, src/Transformacion, src/analisis/*): el manual no inventa logica paralela
  \item Los insumos que referencia existen en el repo: evidencias/tabla\_maestra\_ies.csv, mapping\_inst\_filia\_to\_ies.csv, geografia\_*.csv, ocde\_*.csv, scripts/sprint6\_*.py y artifacts/sprint6\_sankey/*.html
  \item Caveats de alta calidad (Eméritos vitalicios, comparacion contra poblacion con educacion superior, captura de doble afiliacion solo en 2021)
  \item Resolucion de ROOT con fallback (mejor que la del EDA)
\end{itemize}

\paragraph{Debilidades}
\begin{itemize}
  \item Dos celdas con bugs de API que crashean: m['conteos'] (matrices\_todos\_periodos devuelve tuplas) y columna= en hhi\_por\_convocatoria (la firma es col\_geo=)
  \item subprocess hardcodea 'python3', no portable a Windows (el repo vive en Windows)
  \item Depende de datos/raw/produccion\_grupos.csv (\textasciitilde{}1,2 GB) que NO esta en el repo: la celda 2 de carga crashea y las secciones 11-13 no tienen prod
  \item Inconsistencias numericas entre el texto del manual y el codigo real: '22 registros con edad >100 (0,03\%)' vs src/analisis/calidad.py '10 registros (0,013\%)'; '211 IES canonicas / 91\%' vs celda '258 strings mapeados'
  \item Nunca ejecutado de punta a punta (execution\_count null en las 26 celdas de codigo): el manual no tiene evidencia de reproducibilidad
  \item El generador scripts/generar\_manual.py contiene los mismos bugs (los hereda el notebook)
\end{itemize}

\paragraph{Mejoras propuestas}
\begin{itemize}
  \item Corregir la celda de matrices: m es una tupla (conteos, probs) -> iterar con for periodo, (conteos, probs) in matrices.items()
  \item Corregir la celda HHI: hhi\_por\_convocatoria(inv\_clean, col\_geo='NME\_DEPARTAMENTO\_RES\_PR')
  \item Usar 'python' o sys.executable en subprocess para portabilidad a Windows
  \item Anadir celda condicional que avise si falta el dataset de produccion y permita continuar con las secciones que no lo necesitan
  \item Armonizar cifras con src/analisis/calidad.py y scripts/sprint6\_tabla\_maestra\_ies.py, y regenerar con generar\_manual.py
  \item Ejecutar el notebook completo y commitear los outputs como evidencia de reproducibilidad
\end{itemize}

\paragraph{Ejecutabilidad}
Parcial y no reproducible de punta a punta tal cual. Con `poetry run jupyter notebook docs/manual.ipynb` (README) la carga de inv funciona (XLSX existe; openpyxl en lock/requirements), pero: (1) cargar\_produccion() lanza FileNotFoundError (datos/raw/produccion\_grupos.csv ausente) -> la celda 2 crashea y las secciones 11-13 carecen de prod; (2) la celda de matrices lanza TypeError (m['conteos'] sobre una tupla); (3) la celda HHI lanza TypeError (kwarg columna= inexistente); (4) los 4 subprocess con 'python3' fallan en Windows. Las secciones 6-10 (tabla maestra, geografia, sankeys, genero, diversidad) leerian archivos que sí existen en el repo y funcionarian si se superan las anteriores.

\paragraph{¿Produce lo esperado?}
El manual documenta los 12 hallazgos y sus rutas de evidencias/artifacts existen, pero NO es ejecutable tal cual: 3 celdas fallan por API incorrecta y 1 por dependencia de datos ausente. No contradice al dashboard (usa las mismas funciones de src/), pero la cifra '22 atipicos de edad' contradice a src/analisis/calidad.py ('10'), indicando que el texto del manual no se actualizo con el codigo final. El placeholder app/streamlit\_app.py es irrelevante para este notebook.

\begin{tcolorbox}[riesgo]
\textbf{Bugs / problemas detectados:}
\begin{itemize}
  \item Celda seccion 3: `print(m['conteos'])` - matrices\_todos\_periodos devuelve \{periodo: (conteos, probs)\} (tupla, segun docstring y codigo de src/analisis/longitudinal.py linea 164-179); m es una tupla -> TypeError: tuple indices must be integers
  \item Celda seccion 5: `hhi\_por\_convocatoria(inv\_clean, columna='NME\_DEPARTAMENTO\_RES\_PR')` - la firma es (df, col\_geo, col\_year) -> TypeError: unexpected keyword argument 'columna'
  \item Celda 2: `cargar\_produccion()` -> FileNotFoundError porque datos/raw/produccion\_grupos.csv no existe en el repo (hay que descargarlo, \textasciitilde{}1,2 GB)
  \item 4 celdas usan subprocess.run(['python3', ...]): en Windows no existe python3.exe -> FileNotFoundError (portabilidad rota)
  \item Texto vs codigo: '22 registros con edad superior a 100 (0,03\%)' vs src/analisis/calidad.py 'diez registros (0,013\%)'; '211 IES canonicas que mapean el 91\%' vs celda que reporta '258 strings mapeados' (scripts/sprint6\_tabla\_maestra\_ies.py)
\end{itemize}
\end{tcolorbox}

\paragraph{Recomendación de auditoría}
\textbf{refactorizar: el manual es valioso como documentacion de los hallazgos pero no es reproducible tal cual. Corregir los 2 bugs de API (m['conteos'], col\_geo), hacer portable el subprocess, documentar el prerequisito del dataset de produccion (o leer muestras de evidencias/) y armonizar las cifras con src/. Los mismos cambios deben aplicarse en scripts/generar\_manual.py para que el notebook se pueda regenerar.}

\medskip
\hrule

\subsection{Análisis exploratorio de datos}
\noindent Archivos: 1.

\subsubsection{notebooks/01\_eda.ipynb}
\textbf{Rol:} eda · \textbf{Líneas:} 240

\paragraph{Descripción funcional}
EDA exploratorio de 16 celdas (6 markdown + 10 codigo) sobre el dataset consolidado de investigadores reconocidos (77.237 registros, 30 variables, 6 convocatorias 2013-2021). Carga con src/ingesta.cargar\_consolidado, transforma con src/Transformacion.transformar y produce: tipos y faltantes, estadisticas de 3 variables numericas, frecuencias de 5 categoricas y 4 visualizaciones (genero, gran area OCDE, edad por convocatoria, top 15 departamentos). Es el punto de entrada analitico del proyecto, anterior a los sprint 2-6.

\paragraph{Análisis de la lógica}
Celda 1: ROOT = pathlib.Path().resolve().parent (comentario: 'raiz del repo (notebooks/../)', asume cwd = notebooks/); sys.path.insert(0, ROOT/src); imports de pandas/numpy/matplotlib/seaborn y `from ingesta import cargar\_consolidado`, `from Transformacion import transformar`. Celda 2: df\_raw = cargar\_consolidado(); df = transformar(df\_raw); print shape y head. Celdas siguientes analizan df directamente: dtypes, faltantes (isnull().mean()), heatmap de NaN, describe de EDAD\_ANOS\_PR/NRO\_ORDEN\_FORM\_PR/ORDEN\_CLAS\_PR, histogramas, value\_counts de 5 categoricas, pie de genero, barh de areas, boxplot de edad por convocatoria, barplot de top departamentos. Sin cache ni subprocesos; input unico: el consolidado.

\paragraph{Fortalezas}
\begin{itemize}
  \item Sencillo y autocontenido: funciona sin el dataset de produccion (solo necesita el consolidado)
  \item Usa el pipeline real de ingesta+transformacion, misma fuente que el dashboard (consistencia de datos)
  \item Estructura EDA estandar y legible: carga, calidad, numericas, categoricas, visualizaciones
  \item Todas las columnas que referencia (EDAD\_ANOS\_PR, NRO\_ORDEN\_FORM\_PR, ORDEN\_CLAS\_PR, NME\_REGION\_RES\_PR, etc.) existen en el header del CSV consolidado (30 columnas, verificado)
\end{itemize}

\paragraph{Debilidades}
\begin{itemize}
  \item Resolucion de ROOT fragil: depende de que el cwd del kernel sea notebooks/; si JupyterLab arranca el kernel desde la raiz del repo, ROOT apunta al padre del repo y `from ingesta import ...` lanza ModuleNotFoundError (a diferencia de manual.ipynb, no tiene fallback a ROOT.parent)
  \item El analisis de faltantes se hace sobre df YA transformado: transformar imputa mediana en numericas y 'NO REPORTADO' en categoricas (tratar\_faltantes), asi que el heatmap de NaN y los \% de faltantes quedan vacios o minimos (solo sobreviven COD\_DANE\_* y ANO\_CONVO\_FECHA); la seccion 'Calidad de datos' del EDA es poco informativa
  \item Nomenclatura inconsistente: el titulo dice 'Grupo 7' (la carpeta del repo es Grupo 8; Transformacion.py tambien dice Grupo 7)
  \item Usa palette= en sns.boxplot/sns.barplot (deprecacion en seaborn >= 0.14; cosmético, no rompe)
  \item No documenta como se lanza (cwd requerido) ni fija versiones de seaborn
\end{itemize}

\paragraph{Mejoras propuestas}
\begin{itemize}
  \item Adoptar el patron de ROOT del manual.ipynb: Path.cwd() con fallback a ROOT.parent (o derivar de la ubicacion del notebook) para funcionar desde cualquier cwd
  \item Correr el analisis de calidad sobre df\_raw (antes de transformar) o documentar que los faltantes ya fueron imputados/etiquetados
  \item Unificar la nomenclatura de grupo (Grupo 7 vs Grupo 8) en titulo y docstrings
  \item Quitar palette= o fijar seaborn<0.14 en pyproject
  \item Anadir celda que confirme shape == (77237, 30) como asercion de insumo
\end{itemize}

\paragraph{Ejecutabilidad}
Sí, con una condicion de entorno: se ejecuta tal cual si el kernel arranca con cwd = notebooks/ (p. ej. `jupyter notebook notebooks/01\_eda.ipynb` desde la raiz, que en el cliente clasico pone el cwd en notebooks/). Entrada: datos/tarea\_join/investigadores\_consolidado.xlsx via cargar\_consolidado (existe; el CSV de datos/raw no existe, se cae al XLSX). Requiere pandas, numpy, matplotlib, seaborn y openpyxl. No requiere el dataset de produccion. Si el cwd del kernel es la raiz del repo, el import de ingesta falla (ModuleNotFoundError).

\paragraph{¿Produce lo esperado?}
Cumple su rol de EDA: carga y describe el consolidado y genera las 4 visualizaciones sin errores bajo la condicion de cwd correcta. Nota de verificacion: el CSV paralelo datos/tarea\_join/investigadores\_consolidado.csv tiene \textasciitilde{}50.891 filas (no 77.237): conviene confirmar que el XLSX (fuente real de cargar\_consolidado) tenga las 77.237 filas y que no haya divergencia entre formatos. La seccion de faltantes es casi vacia por la imputacion previa de transformar. No contradice al dashboard (misma fuente de datos).

\begin{tcolorbox}[riesgo]
\textbf{Bugs / problemas detectados:}
\begin{itemize}
  \item ROOT = pathlib.Path().resolve().parent (celda 1): si el cwd del kernel no es notebooks/, ROOT no apunta a la raiz del repo y `from ingesta import cargar\_consolidado` lanza ModuleNotFoundError; no hay fallback (a diferencia de docs/manual.ipynb que sí lo tiene)
  \item Analisis de calidad sobre df transformado: transformar (tratar\_faltantes) imputa mediana en numericas y 'NO REPORTADO' en categoricas, por lo que df.isnull() queda casi sin NaN; el mapa de calor y los \% de faltantes de la seccion 2 no muestran la calidad real del dato crudo
  \item Titulo 'Grupo 7' vs carpeta 'Grupo 8' (y Transformacion.py tambien dice Grupo 7): inconsistencia de nomenclatura del proyecto
  \item sns.boxplot/palette y sns.barplot/palette: deprecacion de palette en seaborn>=0.14 (aviso, no falla)
\end{itemize}
\end{tcolorbox}

\paragraph{Recomendación de auditoría}
\textbf{mejorar (mantener): el EDA es valido y util como punto de entrada; endurecer la resolucion de rutas con el patron del manual.ipynb, mover el analisis de calidad a df\_raw (o aclarar que los faltantes ya fueron tratados), unificar la nomenclatura de grupo y fijar/eliminar palette=.}

\medskip
\hrule

\subsection{Pruebas automatizadas}
\noindent Archivos: 2.

\subsubsection{tests/\_\_init\_\_.py}
\textbf{Rol:} testing · \textbf{Líneas:} 0

\paragraph{Descripción funcional}
Marcador de paquete Python para tests/ (vacío, 0 líneas).

\paragraph{Análisis de la lógica}
Sin lógica.

\paragraph{Fortalezas}
\begin{itemize}
  \item Correcto por convención (permite descubrir tests con pytest).
\end{itemize}

\paragraph{Ejecutabilidad}
N/A.

\paragraph{¿Produce lo esperado?}
N/A.

\paragraph{Recomendación de auditoría}
\textbf{mantener --- estándar de Python.}

\medskip
\hrule

\subsubsection{tests/test\_ingesta.py}
\textbf{Rol:} testing · \textbf{Líneas:} 8

\paragraph{Descripción funcional}
Suite de tests del módulo de ingesta. Contiene un único test (test\_catalogo\_exists) que verifica que datos/catalogo.yaml existe.

\paragraph{Análisis de la lógica}
test\_catalogo\_exists usa Path('datos/catalogo.yaml').exists() y assert. Configuración de testpaths en pyproject.toml apunta a tests/.

\paragraph{Fortalezas}
\begin{itemize}
  \item Estructura de tests creada y ejecutable con pytest.
  \item Configuración en pyproject.toml (testpaths) correcta.
\end{itemize}

\paragraph{Debilidades}
\begin{itemize}
  \item Cobertura insuficiente: 1 test trivial; no prueba descargar/guardar/transformar ni ninguna función de análisis.
  \item pandera declarado en dev-dependencies pero sin uso (no hay contratos de esquema).
  \item El test depende del cwd (Path relativo): falla si pytest se ejecuta desde otra carpeta.
\end{itemize}

\paragraph{Mejoras propuestas}
\begin{itemize}
  \item Añadir tests de humo para cargar\_consolidado (shape, columnas, unicidad de ID\_PERSONA\_PR).
  \item Añadir contratos de esquema con pandera (completitud, tipos, codificación).
  \item Usar rutas relativas al proyecto (pathlib con \_\_file\_\_).
\end{itemize}

\paragraph{Ejecutabilidad}
Sí (pytest); verificado que el único test pasa si el catálogo existe. No se ejecutó pytest en la auditoría por prioridad, pero no hay dependencias externas más allá de pathlib.

\paragraph{¿Produce lo esperado?}
Parcial: valida solo la existencia del catálogo, no el comportamiento de la ingesta --- no cumple el objetivo de 'probar' el pipeline.

\begin{tcolorbox}[riesgo]
\textbf{Bugs / problemas detectados:}
\begin{itemize}
  \item Path relativo frágil (depende del directorio de ejecución).
\end{itemize}
\end{tcolorbox}

\paragraph{Recomendación de auditoría}
\textbf{mejorar --- ampliar la cobertura a transformación y análisis; es la acción de mayor impacto en la deuda técnica de pruebas.}

\medskip
\hrule

\subsection{Código legacy (etapa académica inicial)}
\noindent Archivos: 13.

\subsubsection{notebooks/tarea\_anlisis\_bases\_de\_datos/2019\_codigo.py}
\textbf{Rol:} legacy · \textbf{Líneas:} 174

\paragraph{Descripción funcional}
EDA descriptivo de la convocatoria 2019 (Investigadores\_Reconocidos\_por\_convocatoria\_20250909.csv): calcula el \% de faltantes por variable, conteos y gráficos de barras para género, región/departamento de nacimiento y residencia, histogramas de variables numéricas y tablas cruzadas (género x región, gran área x género). Es un script exportado de Colab y corresponde a la tarea de análisis de bases de datos.

\paragraph{Análisis de la lógica}
pandas + matplotlib: lectura con pd.read\_csv, limpieza de nombres de columnas, df.isna().sum() con \% de faltantes ordenado, value\_counts y pd.crosstab(normalize='index') para proporciones por fila, detección de columnas clave por subcadena ('GENER', 'REGION', 'DEPARTAMENTO'), histogramas y heatmap con imshow. No usa librerías de inferencia estadística; es puramente descriptivo.

\paragraph{Fortalezas}
\begin{itemize}
  \item Reporta faltantes por variable con porcentaje ordenado, práctica correcta de EDA
  \item Usa crosstab con proporciones por fila (composición de género dentro de cada gran área), lectura adecuada para comparar grupos
  \item Detecta columnas de forma defensiva por subcadena en el nombre, tolerante a cambios de esquema
  \item Excluye columnas tipo ID/código de los histogramas (evita describir claves)
\end{itemize}

\paragraph{Debilidades}
\begin{itemize}
  \item Ruta hardcodeada de Colab '/content/Investigadores\_Reconocidos\_por\_convocatoria\_20250909.csv', inexistente en el clon y en cualquier equipo local
  \item Usa display() de IPython/Colab; como script Python estándar lanza NameError
  \item Bloque de cálculo de faltantes duplicado (falt y faltantes): código muerto
  \item El heatmap de proporciones no está acompañado de ninguna prueba de independencia (chi-cuadrado)
  \item Etiquetas de porcentaje con desplazamiento fijo 'v + 50' sobre conteos absolutos, frágil si cambia la escala
  \item Las proporciones se calculan sobre n total incluyendo NaN en el denominador; sin documentar
  \item No guarda figuras ni resultados; sin estructura de función o parámetros
\end{itemize}

\paragraph{Mejoras propuestas}
\begin{itemize}
  \item Parametrizar la ruta (Pathlib/argumentos) en vez de hardcodear /content
  \item Reemplazar display() por print() para que corra como script
  \item Eliminar el bloque duplicado de faltantes
  \item Añadir chi2\_contingency a las tablas cruzadas y reportar residuos
  \item Guardar figuras y tablas en una carpeta de salida
\end{itemize}

\paragraph{Ejecutabilidad}
NO. Falta el input '/content/Investigadores\_Reconocidos\_por\_convocatoria\_20250909.csv' (ruta de Colab, no existe en el clon) y además usa display() de IPython, que falla fuera de un entorno de notebook.

\paragraph{¿Produce lo esperado?}
Metodología descriptiva correcta para el objetivo (EDA por convocatoria): no hay errores estadísticos graves porque no realiza inferencia. Único matiz: las proporciones usan el total con NaN y no se prueba la asociación en los cruces, por lo que los heatmaps son solo exploratorios.

\begin{tcolorbox}[riesgo]
\textbf{Bugs / problemas detectados:}
\begin{itemize}
  \item display() no existe en scripts Python estándar (NameError)
  \item Cálculo de faltantes duplicado (falt y faltantes) sin uso del primero
  \item plt.text(pct + 0.5, ...) asume coordenadas de barras horizontales con holgura fija
\end{itemize}
\end{tcolorbox}

\paragraph{Recomendación de auditoría}
\textbf{Mantener como referencia histórica del EDA inicial de la convocatoria 2019. No migrar: el pipeline moderno (src/ingesta/minciencias.py, src/analisis/calidad.py) ya sistematiza faltantes y perfiles descriptivos; este script no aporta lógica reutilizable.}

\medskip
\hrule

\subsubsection{notebooks/tarea\_anlisis\_bases\_de\_datos/datos\_2021.ipynb}
\textbf{Rol:} legacy · \textbf{Líneas:} 2338

\paragraph{Descripción funcional}
Notebook Colab con el EDA de la convocatoria 2021 (Convocatoria 894): filtra Investigadores\_Consolidado por año (21.094 registros x 30 variables), describe, analiza faltantes y duplicados, y gráfica género (pie), edad (histograma + detección de outliers >100 años), región de residencia, nivel de formación, gran área, área, clasificación (violinplot edad x clasificación) y proporciones de víctimas del conflicto, discapacidad y grupo étnico. Cierra con interpretaciones sustantivas en markdown.

\paragraph{Análisis de la lógica}
pandas + matplotlib + seaborn: filtra con df['ANO\_CONVO'].astype(str).str.contains('2021'), describe(), isnull().sum(), duplicated(), value\_counts(), sns.histplot/countplot/violinplot, y crea df\_2021\_filtered (edad <= 100) para los gráficos posteriores. Todo descriptivo; sin pruebas de hipótesis.

\paragraph{Fortalezas}
\begin{itemize}
  \item Identifica y documenta el outlier de edad (máximo 956.25 años) y filtra edad <= 100 antes de las gráficas (decisión de limpieza explícita)
  \item Revisa duplicados y faltantes por columna con interpretación
  \item El violinplot edad x clasificación apoya una conclusión sustantiva correcta (Eméritos mayores, Junior más jóvenes)
  \item Buena práctica de narrativa: cada gráfico va acompañado de interpretación en markdown
\end{itemize}

\paragraph{Debilidades}
\begin{itemize}
  \item Ruta hardcodeada de Colab '/content/Investigadores\_Consolidado.xlsx', inexistente en el clon
  \item La celda del countplot de nivel de formación x género usa df\_2021\_gender\_filtered, variable que NO está definida en ninguna celda del notebook (NameError al re-ejecutar; probablemente se borró la celda que la creaba)
  \item describe() incluye IDs y códigos (media de ID\_PERSONA\_PR y COD\_DANE sin sentido estadístico)
  \item Pie chart para género con categorías de n=4-5 ('Intersexual', 'No disponible'): visualización poco informativa
  \item Salidas embebidas en base64 inflan el archivo a 2338 líneas, dificultando el diff
  \item Filtrado por str.contains('2021') sobre la fecha depende del formato de la columna (frágil si el Excel cambia el tipo)
  \item Sin seed, sin guardado de resultados y sin checklist de calidad al final
\end{itemize}

\paragraph{Mejoras propuestas}
\begin{itemize}
  \item Definir df\_2021\_gender\_filtered o eliminar la celda rota
  \item Excluir IDs/códigos del describe()
  \item Reemplazar el pie por barras horizontales
  \item Limpiar salidas (Output -> Clear) antes de commitear
  \item Parametrizar la lectura del Excel y extraer el año con dt.year en vez de str.contains
\end{itemize}

\paragraph{Ejecutabilidad}
NO. Falta '/content/Investigadores\_Consolidado.xlsx' (no existe en el clon) y la celda de nivel de formación x género lanza NameError porque df\_2021\_gender\_filtered no está definida en el notebook.

\paragraph{¿Produce lo esperado?}
Metodología descriptiva correcta para el objetivo (EDA por convocatoria). No hay errores de inferencia porque no se realizan pruebas; el filtro edad <= 100 es razonable pero el umbral no se justifica y no se hace análisis de sensibilidad (cuántos registros y cómo cambian los perfiles).

\begin{tcolorbox}[riesgo]
\textbf{Bugs / problemas detectados:}
\begin{itemize}
  \item df\_2021\_gender\_filtered usada sin definición (NameError en re-ejecución)
  \item Filtro por str.contains('2021') sobre una columna de fecha: frágil ante cambios de formato
  \item describe() sobre columnas ID/código produce estadísticos sin interpretación
\end{itemize}
\end{tcolorbox}

\paragraph{Recomendación de auditoría}
\textbf{Mantener como referencia del EDA de la convocatoria 2021 (documenta hallazgos de calidad como el outlier de edad). No migrar tal cual: los análisis equivalentes modernos están en src/analisis/ (produccion.py, genero.py, longitudinal.py) con datos ya depurados.}

\medskip
\hrule

\subsubsection{notebooks/tarea\_anlisis\_bases\_de\_datos/inv\_17\_codigo\_graficas.R}
\textbf{Rol:} legacy · \textbf{Líneas:} 614

\paragraph{Descripción funcional}
EDA descriptivo de la convocatoria 2017 (Investigadores\_Reconocidos\_por\_convocatoria\_20250831.xlsx): tablas de frecuencia y gráficos (barras, pie, histograma con densidad, boxplot) para área, gran área, género, país/región/departamento/municipio de nacimiento y de residencia, nivel de formación, clasificación, edad, víctimas del conflicto, grupo étnico, discapacidad e INST\_FILIA. Detecta el valor imposible de edad (952 años) y crea una versión filtrada (edad < 100). Corresponde a la tarea de análisis de bases de datos.

\paragraph{Análisis de la lógica}
readxl + ggplot2 + dplyr + forcats: count() con filtros por umbral de frecuencia para dividir categorías altas/bajas, geom\_bar/geom\_col con coord\_flip y reorder/fct\_infreq, pie charts con coord\_polar, histograma con aes(y = ..density..) + geom\_density + vline de la media. Sin pruebas formales; análisis univariado exploratorio.

\paragraph{Fortalezas}
\begin{itemize}
  \item Reconoce explícitamente las variables de identificación (IDs, códigos DANE) y evita graficarlas, decisión de análisis correcta
  \item Detecta el valor imposible de edad (952) y construye edad\_sin\_error (edad < 100), buen control de calidad de datos
  \item Ordena categorías con reorder/fct\_infreq y separa frecuencias altas y bajas para legibilidad
  \item Autoevalúa gráficos poco informativos ('grafica sin sentido', 'no sirve la grafica'): honestidad analítica y criterio
\end{itemize}

\paragraph{Debilidades}
\begin{itemize}
  \item Ruta hardcodeada 'Downloads/Investigadores\_Reconocidos\_por\_convocatoria\_20250831.xlsx', inexistente en el clon
  \item Usa plot\_missing() de DataExplorer sin cargar library(DataExplorer): error en ejecución
  \item Referencia el objeto Investigadores\_Consolidado (línea 21) sin definirlo en el archivo: error de ejecución
  \item Pie charts con decenas de categorías (país de nacimiento, región): mala práctica de visualización
  \item Umbrales de filtrado arbitrarios y comentarios desalineados con el código (p. ej., comenta 'menor a 140' pero filtra n<40)
  \item Número mágico hardcodeado: sum(NME\_MUNICIPIO\_RES\_PR1\$n)/13001
  \item Sintaxis deprecada geom\_text(aes(label = ..count..)) y aes(y = ..density..) en ggplot2 >= 3.4 (usa after\_stat)
  \item No guarda gráficos ni tablas; sin documentación de parámetros ni versión de los datos
\end{itemize}

\paragraph{Mejoras propuestas}
\begin{itemize}
  \item Cargar DataExplorer o reemplazar plot\_missing por un cálculo propio de faltantes
  \item Definir Investigadores\_Consolidado o eliminar la línea que lo referencia
  \item Sustituir pie charts por barras horizontales cuando hay muchas categorías
  \item Documentar y justificar los umbrales (n>250, n>500, n<40, etc.)
  \item Usar after\_stat(count)/after\_stat(density) y parametrizar la ruta y el denominador (13001)
\end{itemize}

\paragraph{Ejecutabilidad}
NO. Falta 'Downloads/Investigadores\_Reconocidos\_por\_convocatoria\_20250831.xlsx' (no existe en el clon) y además falla en plot\_missing() por falta de library(DataExplorer) y en la referencia a Investigadores\_Consolidado (objeto no definido).

\paragraph{¿Produce lo esperado?}
Descriptivo correcto para el objetivo (EDA por convocatoria): no hay errores de inferencia porque no realiza pruebas. El filtro edad < 100 es apropiado. Los pie charts con muchas categorías son un problema de visualización, no estadístico.

\begin{tcolorbox}[riesgo]
\textbf{Bugs / problemas detectados:}
\begin{itemize}
  \item plot\_missing() sin library(DataExplorer): función no encontrada
  \item Investigadores\_Consolidado referenciado sin ser definido (línea 21)
  \item Sintaxis ..count../..density.. deprecada en ggplot2 moderno (deprecation warning/error según versión)
  \item Denominador hardcodeado 13001 en el cálculo de proporción de municipios de residencia
\end{itemize}
\end{tcolorbox}

\paragraph{Recomendación de auditoría}
\textbf{Mantener como referencia histórica del EDA de la convocatoria 2017 (documenta hallazgos de calidad de datos). No migrar: src/analisis/calidad.py y src/ingesta/minciencias.py sistematizan la limpieza; este script no aporta lógica reutilizable sin reescribirlo.}

\medskip
\hrule

\subsubsection{notebooks/tarea\_dimensiones\_y\_hechos/codigo\_dimensiones\_C.R}
\textbf{Rol:} legacy · \textbf{Líneas:} 40

\paragraph{Descripción funcional}
Crea tablas de dimensiones del modelo estrella a partir de Investigadores\_Consolidado: dimensión investigador (solo ID\_PERSONA\_PR), dimensión municipio de nacimiento (COD\_DANE\_NAC\_PR + país/región/departamento/municipio) y dimensión convocatoria (ID\_CONVOCATORIA + nombre + fecha de convocatoria), exportadas a tres xlsx. Corresponde a la tarea de dimensiones y hechos.

\paragraph{Análisis de la lógica}
readxl + openxlsx + dplyr: unique() sobre la columna de ID, distinct(COD\_DANE\_NAC\_PR, .keep\_all = TRUE) para el municipio, distinct(ID\_CONVOCATORIA, .keep\_all = TRUE) para la convocatoria, y write.xlsx() por dimensión. Sin joins; solo extracción de claves y descriptores.

\paragraph{Fortalezas}
\begin{itemize}
  \item Separa claves naturales (IDs) de atributos descriptivos, enfoque correcto de modelado dimensional
  \item Usa distinct() sobre la clave para garantizar unicidad en las dimensiones
  \item Genera artefactos accionables (xlsx) consumibles por la tabla de hechos
\end{itemize}

\paragraph{Debilidades}
\begin{itemize}
  \item Ruta hardcodeada 'Downloads/Investigadores\_Consolidado.xlsx', inexistente en el clon
  \item Dimensión investigador sin atributos (solo ID\_PERSONA\_PR): tabla sin valor analítico y sin vínculo a institución/universidad
  \item distinct(COD\_DANE\_NAC\_PR, .keep\_all = TRUE) toma la PRIMERA fila por código: si el mismo código aparece con nombres distintos o con NA (investigadores en el 'Exterior'), la dimensión es no determinista
  \item No maneja NA: los códigos DANE faltantes (Exterior) entran como fila NA en la dimensión
  \item Nombres de salida con errores de tipeo (Dimension\_invetigadores.xlsx, Dimension\_municipos\_de\_naciminetos.xlsx) que NO coinciden con los nombres que consume la tabla de hechos (dimensiones\_investigadores.xlsx del notebook): dos implementaciones paralelas de dimensiones incompatibles entre sí
  \item Sin claves sustitutas para municipio y convocatoria; sin validación de cardinalidad
  \item No hay dimensión de formación, clasificación, área ni residencia aquí (la residencia se hace en el notebook, no en este script)
\end{itemize}

\paragraph{Mejoras propuestas}
\begin{itemize}
  \item Unificar con las dimensiones del notebook dimensiones.ipynb (una sola implementación)
  \item Añadir claves sustitutas y atributos a todas las dimensiones (incluida la de investigador)
  \item Decidir y documentar el tratamiento de 'Exterior'/NA antes de construir la dimensión
  \item Corregir nombres de archivo y validar la cardinalidad de la dimensión contra la tabla de hechos
\end{itemize}

\paragraph{Ejecutabilidad}
NO. Falta 'Downloads/Investigadores\_Consolidado.xlsx' en el clon.

\paragraph{¿Produce lo esperado?}
El diseño (claves naturales + descriptores) es conceptualmente correcto para un modelo estrella, pero incompleto: la dimensión investigador no tiene atributos, la de municipio depende del orden de filas (.keep\_all = TRUE) y no se resuelve el caso 'Exterior'. Riesgo de claves inconsistentes entre dimensiones de archivos distintos.

\begin{tcolorbox}[riesgo]
\textbf{Bugs / problemas detectados:}
\begin{itemize}
  \item Typos en los nombres de archivo de salida (invetigadores, municipos de naciminetos) que rompen la trazabilidad con la tabla de hechos
  \item distinct() con .keep\_all = TRUE: resultado dependiente del orden de filas en la base
  \item Filas con COD\_DANE\_NAC\_PR = NA ('Exterior') entran como clave NA en la dimensión
\end{itemize}
\end{tcolorbox}

\paragraph{Recomendación de auditoría}
\textbf{Mantener como referencia histórica del modelado dimensional académico. No migrar: el modelo moderno está en src/modelo/dimensional.py; además conviene consolidar con el notebook de dimensiones para eliminar la duplicación incompatible.}

\medskip
\hrule

\subsubsection{notebooks/tarea\_dimensiones\_y\_hechos/dimensiones.ipynb}
\textbf{Rol:} legacy · \textbf{Líneas:} 289

\paragraph{Descripción funcional}
Notebook Colab que construye tres dimensiones del modelo estrella a partir de Investigadores\_Consolidado: Dim\_Universidad (INST\_FILIA con clave sustituta 'Universidad'), Dim\_Formacion (ID\_NIV\_FORMACION\_PR, NME\_NIV\_FORM\_PR, NRO\_ORDEN\_FORM\_PR) y Dim\_Residencia (COD\_DANE\_RES\_PR + país/región/departamento/municipio), y las exporta a dimensiones\_investigadores.xlsx (3 hojas) con descarga vía files.download. Es la implementación de dimensiones que luego consume tabla\_hechos\_codigo.R.

\paragraph{Análisis de la lógica}
pandas + openpyxl + google.colab.files: dropna().drop\_duplicates().reset\_index(drop=True) por dimensión, insert(0, 'Universidad', range(1, len+1)) para la clave sustituta, pd.ExcelWriter con una hoja por dimensión y files.download().

\paragraph{Fortalezas}
\begin{itemize}
  \item Agrega clave sustituta a Dim\_Universidad (buena práctica de modelado dimensional)
  \item Separa dimensiones por tema y las exporta en un único workbook con hojas nombradas
  \item Conserva la jerarquía de formación en NRO\_ORDEN\_FORM\_PR (ordinal), lo que permite ordenar después
  \item Salida estructurada y lista para el join de hechos
\end{itemize}

\paragraph{Debilidades}
\begin{itemize}
  \item Ruta 'Investigadores\_Consolidado (1).xlsx' (archivo subido a Colab), inexistente en el clon
  \item INST\_FILIA contiene LISTAS separadas por PIPE: drop\_duplicates() trata 'UNIV A|UNIV B' como una 'universidad' falsa (dimension inflada) y dropna() elimina a los investigadores sin filiación
  \item Dim\_Residencia con dropna() elimina TODAS las filas con COD\_DANE\_RES\_PR NA, es decir, todos los investigadores en el 'Exterior' quedan fuera de la dimensión: la tabla de hechos luego hace left\_join y pierde esas claves
  \item Sin clave sustituta para residencia y formación (solo claves naturales); la clave 'Universidad' no se usa en el join de hechos (une por INST\_FILIA)
  \item Modelo incompleto: no hay dimensión de nacimiento, de área de conocimiento ni de clasificación
  \item El código está comentado en su versión de upload (files.upload) y no es reproducible fuera de Colab
\end{itemize}

\paragraph{Mejoras propuestas}
\begin{itemize}
  \item Splitear INST\_FILIA por '|' (formato largo o primera institución documentado) antes de deduplicar
  \item Incluir 'Exterior' como categoría explícita (p. ej., código 0/99999) en lugar de dropna()
  \item Añadir claves sustitutas consistentes y dimensiones de nacimiento/área/clasificación
  \item Documentar la granularidad de la fila de hechos (persona x convocatoria) y validar cardinalidad de la dimensión
  \item Externalizar la lectura del Excel y la ruta de salida
\end{itemize}

\paragraph{Ejecutabilidad}
NO. Requiere subir el Excel en Colab (google.colab.files) y el archivo 'Investigadores\_Consolidado (1).xlsx' no está en el clon; files.download() tampoco existe fuera de Colab.

\paragraph{¿Produce lo esperado?}
Diseño de dimensiones razonable para un modelo estrella académico, pero con pérdida real de datos: el dropna() de Dim\_Residencia excluye al 'Exterior' y la granularidad de INST\_FILIA con pipes queda mal modelada. Estos defectos se propagan a la tabla de hechos (joins con NA).

\begin{tcolorbox}[riesgo]
\textbf{Bugs / problemas detectados:}
\begin{itemize}
  \item dropna() en Dim\_Residencia elimina las filas 'Exterior' (COD\_DANE NA), dejando hechos sin clave de residencia
  \item INST\_FILIA con pipes ('UNIV A|UNIV B') tratada como una sola institución en la dimensión
  \item La clave sustituta 'Universidad' no se corresponde con el join de la tabla de hechos (que une por INST\_FILIA)
\end{itemize}
\end{tcolorbox}

\paragraph{Recomendación de auditoría}
\textbf{Mantener como referencia del diseño dimensional académico. No migrar: src/modelo/dimensional.py es el reemplazo moderno; si se conserva la idea, corregir el manejo de pipes y de 'Exterior'.}

\medskip
\hrule

\subsubsection{notebooks/tarea\_dimensiones\_y\_hechos/tabla\_hechos\_codigo.R}
\textbf{Rol:} legacy · \textbf{Líneas:} 20

\paragraph{Descripción funcional}
Construye la tabla de hechos del modelo estrella: hace left\_join de Investigadores\_Consolidado con Dim\_Universidad (por INST\_FILIA) y con Dim\_genero (por NME\_GENERO\_PR), selecciona columnas por posición numérica y exporta Tabla\_hechos.xlsx. Corresponde a la tarea de dimensiones y hechos.

\paragraph{Análisis de la lógica}
readxl + openxlsx + dplyr: dos left\_join con las dimensiones, subconjunto por índices c(1,4,5,14,15,18,21,26,27:29,31,32) y write.xlsx(). La granularidad resultante es persona x convocatoria.

\paragraph{Fortalezas}
\begin{itemize}
  \item Integra las claves de dimensiones a la fila de hechos mediante joins, enfoque correcto de modelo estrella
  \item El join por INST\_FILIA y NME\_GENERO\_PR es conceptualmente válido si las dimensiones son únicas por esas columnas
\end{itemize}

\paragraph{Debilidades}
\begin{itemize}
  \item Investigadores\_Consolidado NO se define en este archivo: depende de un objeto global de una sesión R anterior (probablemente creado en codigo\_dimensiones\_C.R). No reproducible en una sesión limpia
  \item Selección de columnas por posición numérica (índices mágicos) sin documentar qué columnas son: frágil ante cualquier cambio de esquema
  \item El join por INST\_FILIA arrastra el problema de las dimensiones: investigadores sin filiación o con INST\_FILIA en PIPE quedan con clave de universidad NA o mal asignada
  \item No verifica cardinalidad: si la dimensión no es única por INST\_FILIA, el left\_join duplicaría filas silenciosamente
  \item Rutas hardcodeadas 'Downloads/dimensiones\_investigadores.xlsx' y 'Downloads/Dim\_genero (1).csv', inexistentes en el clon
  \item Nombre de salida sin carpeta explícita ('Tabla\_hechos.xlsx' en el working directory)
\end{itemize}

\paragraph{Mejoras propuestas}
\begin{itemize}
  \item Leer el Excel al inicio del script (autocontenido)
  \item Seleccionar columnas por nombre (dplyr::select) y documentar el diccionario de la tabla de hechos
  \item Validar que nrow() no cambie tras cada join y reportar claves NA
  \item Escribir con ruta explícita y verificable
\end{itemize}

\paragraph{Ejecutabilidad}
NO. Falta el objeto Investigadores\_Consolidado (no definido en el archivo) y los inputs 'Downloads/dimensiones\_investigadores.xlsx' y 'Downloads/Dim\_genero (1).csv' no existen en el clon.

\paragraph{¿Produce lo esperado?}
La lógica de hechos es correcta para un modelo estrella simple (fila persona x convocatoria con claves de dimensión), pero la selección por posición y la dependencia de objetos globales la hacen propensa a errores. Nota: la tabla de hechos no contiene medidas/indicadores, solo claves; aceptable como entregable académico pero de uso analítico limitado.

\begin{tcolorbox}[riesgo]
\textbf{Bugs / problemas detectados:}
\begin{itemize}
  \item Investigadores\_Consolidado no definido en el script (error de objeto no encontrado)
  \item Índices mágicos c(1,4,5,14,15,18,21,26,27:29,31,32) sin correspondencia verificable con el esquema
  \item Sin control de cardinalidad tras los joins (riesgo de duplicación de filas)
\end{itemize}
\end{tcolorbox}

\paragraph{Recomendación de auditoría}
\textbf{Mantener como referencia del entregable de hechos. No migrar: src/modelo/dimensional.py genera el modelo estrella moderno; si se quisiera conservar, reescribir con joins por nombre, lectura autocontenida y validaciones.}

\medskip
\hrule

\subsubsection{notebooks/tarea\_gran\_tabla/analisis2.ipynb}
\textbf{Rol:} legacy · \textbf{Líneas:} 294

\paragraph{Descripción funcional}
Análisis bivariado de la GRAN\_TABLA(SIN ID): matriz de correlación entre numéricas, countplots con porcentajes (género x convocatoria, clasificación x gran área, formación x clasificación), tablas de contingencia con pruebas chi-cuadrado y coeficiente de contingencia de Pearson (clasificación x región, formación x región, convocatoria x gran área), y Kruskal-Wallis + post-hoc de Dunn con Bonferroni para edad según gran área. NOTA: a pesar de la extensión .ipynb, el contenido es un script Python exportado de Colab (no es un notebook JSON), lo que sugiere una descarga mal nombrada.

\paragraph{Análisis de la lógica}
pandas + seaborn + scipy: pd.crosstab, chi2\_contingency, np.sqrt(chi2/(chi2+n)) para el coeficiente de contingencia, scipy.stats.kruskal y scikit\_posthocs.posthoc\_dunn(p\_adjust='bonferroni'). Incluye anotación manual de porcentajes sobre barras apiladas y !pip install dentro del script.

\paragraph{Fortalezas}
\begin{itemize}
  \item Usa chi-cuadrado con coeficiente de contingencia de Pearson: medida de asociación adecuada para variables nominales
  \item Elige Kruskal-Wallis (no paramétrico) para edad por gran área: robusto ante la no normalidad de la edad
  \item Aplica post-hoc de Dunn con corrección Bonferroni tras Kruskal-Wallis: secuencia inferencial correcta
  \item Ordena la tabla de contingencia por jerarquía de clasificación (ORDEN\_CLAS\_PR) para lectura sustantiva
\end{itemize}

\paragraph{Debilidades}
\begin{itemize}
  \item La extensión .ipynb es engañosa: el archivo es un script .py exportado de Colab; no abre como notebook
  \item Ruta hardcodeada '/content/GRAN\_TABLA(SIN ID).xlsx', inexistente en el clon
  \item !pip install scikit-posthocs dentro del código: magia de notebook que falla como script y rompe la reproducibilidad de dependencias
  \item Chi-cuadrado sin verificación de frecuencias esperadas: las tablas región x clasificación y región x formación tienen celdas pequeñas ('Exterior', 'No disponible') que invalidan la aproximación asintótica
  \item Correlación de Pearson sobre variables ordinales (NRO\_ORDEN\_FORM\_PR, ORDEN\_CLAS\_PR) tratadas como continuas; lo adecuado es Spearman
  \item El bloque de anotación de porcentajes usa aritmética frágil sobre ticks de matplotlib (código muerto 'categoria' y riesgo de división por cero)
  \item No guarda resultados ni figuras; sin semilla (no aplica) ni control de múltiples comparaciones más allá de Bonferroni
\end{itemize}

\paragraph{Mejoras propuestas}
\begin{itemize}
  \item Renombrar el archivo a .py y eliminar la magia !pip install (declarar dependencia en requirements)
  \item Parametrizar la ruta de lectura
  \item Reportar la mínima frecuencia esperada o agrupar categorías raras antes del chi-cuadrado (o usar test exacto)
  \item Usar Spearman para ordinales y simplificar las anotaciones de barras
  \item Guardar tablas de resultados en CSV
\end{itemize}

\paragraph{Ejecutabilidad}
NO. Falta '/content/GRAN\_TABLA(SIN ID).xlsx' y las magias de Colab (!pip install) fallan al ejecutarlo como script Python estándar.

\paragraph{¿Produce lo esperado?}
Metodología inferencial mayormente correcta: chi-cuadrado para asociación de nominales y Kruskal-Wallis + Dunn para edad por grupo. Riesgo: celdas con frecuencias esperadas bajas debilitan el chi-cuadrado, y con n \textasciitilde{} 50.000 casi cualquier asociación resulta significativa, por lo que el coeficiente de contingencia (bien calculado) es la métrica realmente informativa.

\begin{tcolorbox}[riesgo]
\textbf{Bugs / problemas detectados:}
\begin{itemize}
  \item Extensión .ipynb con contenido de script .py (no es JSON de notebook)
  \item !pip install dentro del archivo (falla fuera de notebook)
  \item Lógica de anotación de barras dependiente del espaciado de xticks (puede dividir por cero o etiquetar mal con hue)
  \item Pearson aplicado a variables ordinales
\end{itemize}
\end{tcolorbox}

\paragraph{Recomendación de auditoría}
\textbf{Mantener como referencia del análisis bivariado de la gran tabla. No migrar tal cual: los análisis equivalentes modernos están en src/analisis/ (genero.py, territorial.py, produccion.py); si se quiere conservar el chi-cuadrado, reimplementarlo en el pipeline con validación de supuestos y guardado de resultados.}

\medskip
\hrule

\subsubsection{notebooks/tarea\_gran\_tabla/analisis\_de\_clusters.R}
\textbf{Rol:} legacy · \textbf{Líneas:} 94

\paragraph{Descripción funcional}
Clustering k-prototypes sobre la GRAN\_TABLA(SIN ID) para datos mixtos: convierte las categóricas en factor/ordinal, elimina casos incompletos con cc(), aplica kproto(k = 3) con set.seed(123) y describe los clusters con tablas de proporción por variable (gran área, género, región de nacimiento/residencia, formación, clasificación, discapacidad) y medias de edad. Corresponde a la tarea de análisis multivariado de la gran tabla.

\paragraph{Análisis de la lógica}
klaR + clustMixType + factoextra + dplyr: mutate(across(c(1:10,16:23), factor) y ordered() para formación/clasificación), cc() para casos completos (retiene el 91.6\% de la base), kproto(GT\_NA, k = 3), filter por cluster y table()/mean() por grupo para perfilado.

\paragraph{Fortalezas}
\begin{itemize}
  \item Elección correcta del método: k-prototypes (clustMixType) es el algoritmo apropiado para datos mixtos categórico-numérico
  \item Fija semilla (set.seed(123)): el resultado del kproto es reproducible
  \item Reporta explícitamente la proporción de datos retenidos tras eliminar casos incompletos (91.6\%)
  \item Perfilado de clusters con proporciones por variable y media de edad: enfoque descriptivo adecuado para interpretar grupos
\end{itemize}

\paragraph{Debilidades}
\begin{itemize}
  \item k = 3 fijado sin validación: no hay curva de withinss/silhouette, ni comparación de varios k, ni análisis del lambda (importancia de variables) de kproto
  \item cc() (función de casos completos del paquete mice) elimina el 8.4\% de las filas sin evaluar sesgo: los registros faltantes pueden ser sistemáticos (p. ej., 'Exterior', sin filiación, sin código DANE)
  \item mice NO está cargado en este script (solo readxl, klaR, dplyr, factoextra, clustMixType): cc() solo existe si se ejecutó antes analisis\_multivariado\_base\_gran\_table.R en la misma sesión: no reproducible
  \item Incluye el outlier de edad (956 años) en el clusterizado sin filtrar (los análisis previos sí lo descartaban)
  \item Categorías 'No disponible'/'Exterior' dominantes distorsionan las distancias mixtas del kproto (el algoritmo estandariza categorías de forma uniforme)
  \item El perfilado es solo inspección visual de proporciones: sin chi-cuadrado/ANOVA por variable vs cluster ni tamaño de clusters reportado
  \item Ruta hardcodeada 'Downloads/GRAN\_TABLA(SIN ID).xlsx', inexistente en el clon
\end{itemize}

\paragraph{Mejoras propuestas}
\begin{itemize}
  \item Validar k con varios valores (p. ej., k = 2..8) usando métricas de clustMixType (lambda, withinss) o silhouette para datos mixtos
  \item Cargar mice (o definir casos completos con complete.cases()) y documentar el sesgo de la eliminación; evaluar imputación como alternativa
  \item Filtrar edad > 100 antes de clusterizar (coherencia con la limpieza previa)
  \item Reportar el lambda de cada variable (contribución al clustering) y validar perfiles con pruebas formales
  \item Documentar la granularidad y el número de filas por cluster
\end{itemize}

\paragraph{Ejecutabilidad}
NO. Falta 'Downloads/GRAN\_TABLA(SIN ID).xlsx' en el clon y, en una sesión limpia, cc() no existe porque mice no está cargado (depende de ejecutar antes el script multivariado en la misma sesión R).

\paragraph{¿Produce lo esperado?}
El método (k-prototypes) es el adecuado para datos mixtos, pero la selección de k sin validación y la eliminación de NA por casos completos sin evaluar sesgo son errores metodológicos: los clusters pueden ser artefactos del número de grupos impuesto y de la pérdida de 8.4\% de la muestra. Los perfiles por proporción son descriptivos y no validan la estructura encontrada.

\begin{tcolorbox}[riesgo]
\textbf{Bugs / problemas detectados:}
\begin{itemize}
  \item cc() (mice) usado sin cargar mice en el script: dependencia oculta de sesión previa
  \item k = 3 fijo sin justificación ni validación (no hay análisis de sensibilidad a k)
  \item Outlier de edad (956 años) incluido en el clusterizado
  \item Eliminación de casos incompletos (8.4\%) sin análisis de sesgo
\end{itemize}
\end{tcolorbox}

\paragraph{Recomendación de auditoría}
\textbf{Mantener como referencia académica de la tarea (muestra el uso correcto de k-prototypes para datos mixtos). No migrar: el pipeline moderno no incluye clustering; si se reintroduce, hacerlo con validación de k, manejo de NA documentado y datos depurados.}

\medskip
\hrule

\subsubsection{notebooks/tarea\_gran\_tabla/analisis\_multivariado\_base\_gran\_table.R}
\textbf{Rol:} legacy · \textbf{Líneas:} 405

\paragraph{Descripción funcional}
Análisis multivariado de la gran tabla: análisis de correspondencia (CA) para pares de categóricas (gran área x género, INST\_FILIA x género, nivel de formación x clasificación), análisis de correspondencia múltiple (MCA) para conjuntos de categóricas (5 variables; 3 variables de población) y FAMD (factor analysis of mixed data) con ANO\_CONVO, gran área, género, país de nacimiento y EDAD\_ANOS\_PR. Reporta eigenvalores, contribuciones y cos2, y construye biplots con ggrepel. Incluye análisis de faltantes (plot\_missing, md.pattern) y retira una fila totalmente vacía.

\paragraph{Análisis de la lógica}
FactoMineR (CA, MCA, FAMD) + factoextra + ggplot2 + ggrepel + mice + DataExplorer: table() -> CA(bc/ bd/ bf, graph = F), MCA(GT[, c(2,5,6,7,16)], graph = F) y MCA(GT[, c(20:22)], ncp = 16), FAMD(GT[-50891, c(2,5,6,7,15)], ncp = 20, graph = T); extrae b\$eig, \$row/\$col/\$var/\$quanti.var/\$quali.var coord/cos2/contrib y grafica varianza y biplots.

\paragraph{Fortalezas}
\begin{itemize}
  \item Elección correcta de método según tipo de datos: CA para tablas de contingencia 2 x k, MCA para múltiples categóricas, FAMD para datos mixtos
  \item Reporta contribuciones y cos2 (calidad de representación de filas/columnas), buenas prácticas de lectura de CA/MCA/FAMD
  \item Detecta y elimina la fila totalmente vacía (md.pattern / !is.finite) antes del FAMD
  \item Gráficos de varianza explicada y acumulada con umbral del 70\% para decidir número de dimensiones
\end{itemize}

\paragraph{Debilidades}
\begin{itemize}
  \item Ruta hardcodeada 'Downloads/GRAN\_TABLA(SIN ID).xlsx', inexistente en el clon
  \item CA sobre INST\_FILIA x género con miles de filas (instituciones y combinaciones PIPE sin normalizar): tabla dispersa donde las categorías con frecuencia 1 dominan el chi-cuadrado; resultado poco interpretable
  \item Remoción de la fila 50891 por índice hardcodeado (GT[-50891, ...]): frágil si cambia la base; debería ser un filtro por contenido
  \item Comentario del MCA (dice NME\_PAIS\_RES\_PR) no coincide con el código (columna 16): documentación desalineada
  \item Los CA no van acompañados de chisq.test: no se reporta significancia ni frecuencias esperadas
  \item MCA/FAMD con categorías 'No disponible'/'Exterior' y NAs sin tratamiento explícito documentado (depende del default na.method de FactoMineR)
  \item FAMD con ncp = 20 para solo 5 variables: sobredimensionado
  \item No guarda resultados (coordenadas/contribuciones) en archivos
\end{itemize}

\paragraph{Mejoras propuestas}
\begin{itemize}
  \item Excluir INST\_FILIA del CA o agregar a nivel de institución (eliminar pipes) antes de analizar
  \item Acompañar cada CA con chisq.test y verificación de frecuencias esperadas
  \item Eliminar la fila vacía con un filtro (p. ej., rowSums(is.na()) < ncol) en vez del índice 50891
  \item Alinear comentarios con columnas y documentar el tratamiento de NA
  \item Guardar eig/coord/cos2/contrib en xlsx para trazabilidad
\end{itemize}

\paragraph{Ejecutabilidad}
NO. Falta 'Downloads/GRAN\_TABLA(SIN ID).xlsx' en el clon.

\paragraph{¿Produce lo esperado?}
Metodología en general correcta: CA/MCA/FAMD están bien aplicados según el tipo de variables. Riesgos: el CA sobre INST\_FILIA con tabla dispersa infla la inercia de categorías raras (poca interpretabilidad); FAMD con ncp=20 es inofensivo pero sobredimensionado; la lectura de contribuciones sin pruebas de significancia puede llevar a sobreinterpretación.

\begin{tcolorbox}[riesgo]
\textbf{Bugs / problemas detectados:}
\begin{itemize}
  \item Índice hardcodeado 50891 para eliminar la fila vacía (rompe con cualquier cambio de la base)
  \item Comentario del MCA (NME\_PAIS\_RES\_PR) desalineado con la columna 16 del código
  \item CA sobre INST\_FILIA sin normalizar los valores PIPE (instituciones múltiples tratadas como categorías únicas)
  \item Sintaxis ggplot con aes(shape = '...') dentro de geom\_point que no varía la forma real (leyenda engañosa)
\end{itemize}
\end{tcolorbox}

\paragraph{Recomendación de auditoría}
\textbf{Mantener como referencia académica (buen ejemplo de aplicación de CA/MCA/FAMD). No migrar tal cual: el pipeline moderno no realiza análisis factorial; si se requiere en el futuro, reimplementarlo con datos limpios, validación de supuestos y guardado de resultados.}

\medskip
\hrule

\subsubsection{notebooks/tarea\_gran\_tabla/codigo\_analisis\_univariado.py}
\textbf{Rol:} legacy · \textbf{Líneas:} 123

\paragraph{Descripción funcional}
Análisis univariado de la gran tabla: revisa tipos de datos, describe() de las numéricas (NRO\_ORDEN\_FORM\_PR, ORDEN\_CLAS\_PR, EDAD\_ANOS\_PR) con \% de faltantes, histogramas y boxplots, value\_counts de las categóricas, gráficos de barras de las principales variables y un bloque robusto de parsing de ANO\_CONVO (texto dd/mm/aaaa o serial de Excel) para contar investigadores por año. Corresponde a la tarea de análisis univariado de la gran tabla.

\paragraph{Análisis de la lógica}
pandas + seaborn + matplotlib: describe().T con missing\_\% agregada, sns.histplot/boxplot, value\_counts(dropna=False), y parsing de fecha con doble estrategia: pd.to\_datetime(dayfirst=True) y, como respaldo, pd.to\_datetime(serial, unit='D', origin='1899-12-30'); barplot por año de convocatoria.

\paragraph{Fortalezas}
\begin{itemize}
  \item Parsing de fechas robusto con doble estrategia (texto y serial de Excel): buena práctica defensiva para datos heterogéneos
  \item Reporta \% de faltantes junto con los estadísticos descriptivos
  \item Separa explícitamente variables numéricas, temporales y categóricas antes de analizar
  \item Análisis claro y estructurado, fácil de seguir
\end{itemize}

\paragraph{Debilidades}
\begin{itemize}
  \item Ruta hardcodeada '/content/GRAN\_TABLA(SIN ID).xlsx' (Colab), inexistente en el clon
  \item Lee el Excel dos veces (todo el archivo y luego solo la columna ANO\_CONVO): ineficiente y evidencia de parche sobre el dato original
  \item Inconsistencia en el tratamiento de NA: las tablas usan value\_counts(dropna=False) pero los gráficos usan el default dropna=True
  \item Las columnas numéricas llegan como texto/object (heredado de la creación de la gran tabla con dtype='object' en codigo\_creacion\_gran\_tabla\_a\_partir\_del\_modelo\_estrella.py), obligando a conversiones ad hoc: síntoma de un problema de tipos aguas arriba
  \item Sin pruebas estadísticas ni guardado de resultados
\end{itemize}

\paragraph{Mejoras propuestas}
\begin{itemize}
  \item Leer ANO\_CONVO una sola vez y reutilizar el DataFrame
  \item Uniformizar el tratamiento de NA entre tablas y gráficos
  \item Corregir el dtype en el origen (creación de la gran tabla) en vez de parchear aquí
  \item Parametrizar la ruta y guardar el resumen descriptivo en CSV
\end{itemize}

\paragraph{Ejecutabilidad}
NO. Falta '/content/GRAN\_TABLA(SIN ID).xlsx' (no existe en el clon; es una ruta de Colab).

\paragraph{¿Produce lo esperado?}
Correcto como EDA univariado: no hay errores metodológicos porque no realiza inferencia. El parsing defensivo de fechas es una buena práctica. El único riesgo es interpretar estadísticos de variables ordinales (NRO\_ORDEN\_FORM\_PR) como continuas, pero el script no extrae conclusiones inferenciales.

\begin{tcolorbox}[riesgo]
\textbf{Bugs / problemas detectados:}
\begin{itemize}
  \item Doble lectura del Excel (desperdicio y riesgo de discrepancias entre lecturas)
  \item value\_counts con dropna=False en tablas vs dropna=True en gráficos (conteos inconsistentes)
  \item Dependencia de parche: supone que ANO\_CONVO viene mal tipada y la re-parsea
\end{itemize}
\end{tcolorbox}

\paragraph{Recomendación de auditoría}
\textbf{Mantener como referencia del EDA univariado de la gran tabla. No migrar: src/analisis/calidad.py y los scripts de sprint (sprint2\_duckdb.py, sprint2\_genero\_ocde.py) ya cubren descriptivos y perfiles con datos tipados correctamente.}

\medskip
\hrule

\subsubsection{notebooks/tarea\_gran\_tabla/codigo\_creacion\_gran\_tabla\_a\_partir\_del\_modelo\_estrella.py}
\textbf{Rol:} legacy · \textbf{Líneas:} 205

\paragraph{Descripción funcional}
Ensambla la gran tabla a partir del modelo estrella: lee todos los CSV/Excel de una carpeta, detecta la tabla de hechos por nombre (regex 'hecho|hechos|fact') o por máximo de filas, detecta claves de unión con heurísticas (columnas comunes únicas, llaves compuestas, columnas tipo ID) y hace merges left con guarda anti-1:N; exporta gran\_tabla.* y un reporte de uniones, y luego elimina columnas ID para producir GRAN\_TABLA(SIN ID). Corresponde a la tarea de creación de la gran tabla a partir del modelo estrella.

\paragraph{Análisis de la lógica}
pandas + re + pathlib: read\_any() por archivo/hoja, guess\_fact\_name(), candidate\_keys\_for\_dim() (nunique == len, combinaciones únicas, columnas id\_like), safe\_left\_join() que revierte el merge si aumenta el número de filas (1:N), merge(how='left') con rename de columnas con prefijo de dimensión, drop de IDs y export a CSV/Excel (xlsxwriter/openpyxl).

\paragraph{Fortalezas}
\begin{itemize}
  \item Heurísticas de detección de claves razonables y defensivas: verifica unicidad y revierte joins 1:N en lugar de duplicar filas silenciosamente
  \item Genera un reporte de uniones (reporte\_uniones.csv/xlsx) que permite auditar qué dimensiones se integraron y con qué claves
  \item Separa la versión completa (con IDs) de la versión analítica sin IDs (GRAN\_TABLA(SIN ID))
  \item Tolerante a formatos múltiples (CSV/Excel y múltiples hojas)
\end{itemize}

\paragraph{Debilidades}
\begin{itemize}
  \item Ruta hardcodeada '/content/A123' (placeholder del drive del autor), inexistente en el clon; y claves\_forzadas apunta a 'MiArchivo\_\_Dimension\_municipos\_de\_naciminetos\_R', nombre que nunca coincide con los archivos reales del repo: la única clave forzada no aplica y todo queda a merced de las heurísticas
  \item !pip install xlsxwriter y display() dentro del script: magias de Colab que fallan como script estándar
  \item Lee todo con dtype='object': las columnas numéricas salen como texto y contaminan la gran tabla (por eso los análisis posteriores re-parsan fechas y números, p. ej., codigo\_analisis\_univariado.py)
  \item La heurística de clave puede unir por una columna común equivocada (p. ej., NME\_CONVOCATORIA) sin verificación semántica de que sea clave en ambas tablas
  \item Renombra columnas con prefijos largos tipo 'Archivo\_\_Hoja\_\_col', generando nombres difíciles de usar
  \item No versiona inputs ni dimensiones: el resultado depende del contenido exacto de la carpeta
  \item Import de re duplicado (líneas 11 y 13) y variable pat\_fact definida dos veces
\end{itemize}

\paragraph{Mejoras propuestas}
\begin{itemize}
  \item Eliminar magias de Colab (pip/display) y parametrizar la carpeta de entrada
  \item Leer con tipos inferidos o tipar explícitamente las columnas numéricas
  \item Verificar semánticamente las claves (unicidad en hechos y en dimensión) y agregar validación cruzada
  \item Persistir el reporte de uniones como artefacto del pipeline y versionar los inputs
  \item Usar prefijos de columna cortos y estables
\end{itemize}

\paragraph{Ejecutabilidad}
NO. Falta la carpeta '/content/A123' con las tablas del modelo estrella (y los nombres forzados 'MiArchivo\_\_...' no existen en el repo); además !pip install y display() son incompatibles con un script estándar.

\paragraph{¿Produce lo esperado?}
La lógica de ensamblaje (star schema -> tabla denormalizada) es correcta en concepto y tiene buenas salvaguardas (reversión de 1:N, reporte de uniones). El problema grave es dtype='object', que degrada tipos numéricos/ordenados y obliga a parches aguas abajo, y la falta de versionado de inputs, que hace el resultado no reproducible con los datos actuales del repositorio.

\begin{tcolorbox}[riesgo]
\textbf{Bugs / problemas detectados:}
\begin{itemize}
  \item dtype='object' al leer Excel convierte todas las columnas en texto (pérdida de tipos numéricos)
  \item claves\_forzadas apunta a un nombre inexistente ('MiArchivo\_\_Dimension\_municipos\_de\_naciminetos\_R'): nunca se usa
  \item !pip install dentro de un archivo .py (falla fuera de notebook)
  \item La detección de la tabla de hechos por máximo de filas puede elegir la tabla equivocada si ninguna coincide con la regex
\end{itemize}
\end{tcolorbox}

\paragraph{Recomendación de auditoría}
\textbf{Mantener como referencia del proceso de ensamblaje (buena idea de reporte de uniones). No migrar tal cual: el modelo moderno (src/modelo/dimensional.py + scripts/sprint2\_duckdb.py) ya construye la gran tabla con tipos correctos y control de calidad; este script solo aporta la idea de auditar las uniones.}

\medskip
\hrule

\subsubsection{notebooks/tarea\_join/codigo\_join.py}
\textbf{Rol:} legacy · \textbf{Líneas:} 26

\paragraph{Descripción funcional}
Une verticalmente (concatenación) las bases por convocatoria 2017, 2019 y 2021 (Investigadores\_Reconocidos\_por\_convocatoria\_20250829*.csv) para crear Investigadores\_Consolidado.xlsx y .csv y descargarlos. Corresponde a la tarea de join (unión de las bases por convocatoria en una base consolidada).

\paragraph{Análisis de la lógica}
pandas: pd.concat([df\_2017, df\_2019, df\_2021], ignore\_index=True), to\_excel() y to\_csv(), y descarga con google.colab.files.download(). Operación de unión vertical (union all de registros persona x convocatoria).

\paragraph{Fortalezas}
\begin{itemize}
  \item Código simple y claro: hace exactamente lo que la tarea pide (unión vertical de las tres convocatorias)
  \item Genera la base en ambos formatos (Excel y CSV) para consumo posterior
\end{itemize}

\paragraph{Debilidades}
\begin{itemize}
  \item Rutas hardcodeadas de Colab '/content/Investigadores\_Reconocidos\_por\_convocatoria\_20250829*.csv', inexistentes en el clon
  \item from google.colab import files: solo funciona en Colab; como script local falla
  \item concat sin verificar que las tres bases tengan el mismo esquema (mismas columnas, mismo orden, mismos tipos): un cambio de esquema genera columnas NaN o tipos mezclados silenciosamente
  \item Sin deduplicación ni control de calidad: los duplicados de ID\_PERSONA\_PR se manejan después en arreglar\_base\_consultoria\_7.R (acoplamiento frágil entre scripts)
  \item Versiones de archivo inconsistentes entre tareas: aquí 20250829, pero 2019\_codigo.py usa 20250909 e inv\_17\_codigo\_graficas.R usa 20250831: los análisis de distintas tareas NO se hicieron sobre el mismo snapshot de datos
  \item Sin documentación de la granularidad resultante (fila = investigador x convocatoria)
\end{itemize}

\paragraph{Mejoras propuestas}
\begin{itemize}
  \item Verificar igualdad de columnas y tipos antes del concat (assert o comparación de listas)
  \item Deduplicar con criterio documentado en el mismo script o validar en uno solo
  \item Parametrizar rutas y fijar una única versión de archivos fuente para todo el proyecto
  \item Registrar un manifiesto de versiones de inputs
\end{itemize}

\paragraph{Ejecutabilidad}
NO. Faltan los tres CSV en '/content/' (no existen en el clon) y google.colab.files no existe fuera de Colab.

\paragraph{¿Produce lo esperado?}
El concat vertical es la operación correcta para unir convocatorias (no hay error metodológico en la unión en sí). El riesgo es la ausencia de control de esquema y de deduplicación, que genera inconsistencias aguas abajo (duplicados, tipos mezclados) que luego otros scripts deben parchear.

\begin{tcolorbox}[riesgo]
\textbf{Bugs / problemas detectados:}
\begin{itemize}
  \item Import de google.colab.files fuera de Colab (ImportError)
  \item concat sin validación de esquema: si una convocatoria cambia de columnas, se crean NaN silenciosos
  \item Versiones de archivo fuente distintas entre tareas (20250829 vs 20250909 vs 20250831): datos no comparables entre scripts
\end{itemize}
\end{tcolorbox}

\paragraph{Recomendación de auditoría}
\textbf{Mantener como referencia histórica de la tarea de join. No migrar: src/ingesta/minciencias.py hace la ingesta consolidada moderna con validación de esquema y versionado de datos.}

\medskip
\hrule

\subsubsection{notebooks/tarea\_tabla/arreglar\_base\_consultoria\_7.R}
\textbf{Rol:} legacy · \textbf{Líneas:} 125

\paragraph{Descripción funcional}
Depura Investigadores\_Consolidado: diagnostica duplicados de ID\_PERSONA\_PR, deduplica conservando una fila por persona con prioridad de convocatoria (21 > 20 > 19) mediante slice(1), y audita las diferencias de edad entre registros duplicados (dif\_edad > 4.1) revisando casos anómalos de forma manual con View. Corresponde a la tarea de arreglo de la base (preparación para análisis a nivel de persona).

\paragraph{Análisis de la lógica}
readxl + dplyr: duplicated() y sum(table(...) > 1) para diagnóstico, group\_by(ID\_PERSONA\_PR) + mutate(prioridad = case\_when(ID\_CONVOCATORIA == 21 \textasciitilde{} 3, ...)) + arrange(desc(prioridad)) + slice(1) + ungroup() para deduplicar, y mutate(dif\_edad = EDAD\_ANOS\_PR - dplyr::lag(EDAD\_ANOS\_PR)) dentro de group\_by para auditar la coherencia temporal de la edad entre convocatorias.

\paragraph{Fortalezas}
\begin{itemize}
  \item Enfoque de calidad de datos correcto: detecta duplicados, elige el registro con una jerarquía de convocatoria explícita y verificable, y valida la unicidad después de slice(1)
  \item Control de calidad interesante: calcula la diferencia de edad entre registros duplicados de la misma persona (debería ser \textasciitilde{}2-4 años entre convocatorias) y detecta inconsistencias (dif\_edad > 4.1)
  \item Decisión documentada en el código (comentarios con los casos revisados y sus diferencias)
  \item Reporta la proporción de filas retenidas tras la deduplicación
\end{itemize}

\paragraph{Debilidades}
\begin{itemize}
  \item Ruta hardcodeada 'Downloads/Investigadores\_Consolidado.xlsx', inexistente en el clon
  \item write.xlsx() sin cargar library(openxlsx) (solo readxl y dplyr están cargados) y sin extensión en el nombre ('Base\_sin\_duplicados'): error en ejecución y archivo sin formato
  \item La prioridad por convocatoria crea un snapshot heterogéneo: los investigadores que aparecen en varias convocatorias quedan representados con sus datos del año más reciente, mientras que los singletons conservan el único año que tienen; la edad, formación y clasificación de distintas personas provienen de años distintos (sesgo de selección de año)
  \item Los casos con dif\_edad > 4.1 se inspeccionan visualmente pero NO se decide ni se registra su tratamiento (no se excluyen, no se corrigen): auditoría incompleta
  \item slice(1) sin desempate: si la misma persona tiene dos registros con la misma prioridad (duplicado dentro de la misma convocatoria), se elige una fila arbitraria (no determinista)
  \item No valida el resultado contra el total esperado ni contra la base original más allá de nrow ratio
\end{itemize}

\paragraph{Mejoras propuestas}
\begin{itemize}
  \item Cargar openxlsx y escribir 'Base\_sin\_duplicados.xlsx' con ruta explícita
  \item Documentar y aplicar una política para los casos con dif\_edad > 4.1 (excluir, imputar o marcar)
  \item Añadir un orden determinista dentro de slice(1) (p. ej., arrange por año o por nivel de formación)
  \item Evaluar el sesgo del snapshot mixto: reportar cuántos registros provienen de cada convocatoria tras la deduplicación
  \item Comparar la edad imputada vs la declarada entre convocatorias como control adicional
\end{itemize}

\paragraph{Ejecutabilidad}
NO. Falta 'Downloads/Investigadores\_Consolidado.xlsx' en el clon y write.xlsx() falla por no tener cargado openxlsx (además el nombre de archivo carece de extensión).

\paragraph{¿Produce lo esperado?}
La deduplicación con jerarquía es una decisión válida para análisis a nivel persona, pero introduce un sesgo de selección de año (mezcla de snapshots temporales). El análisis de dif\_edad es un buen control de calidad incompleto: identifica anomalías pero no las resuelve ni documenta la decisión. La base resultante solo se valida por conteo, no por contenido.

\begin{tcolorbox}[riesgo]
\textbf{Bugs / problemas detectados:}
\begin{itemize}
  \item write.xlsx() sin library(openxlsx) y con nombre sin extensión (error de ejecución y archivo sin formato)
  \item slice(1) sin orden determinista ante empates de prioridad (resultado no reproducible)
  \item Caso de auditoría (dif\_edad > 4.1) revisado manualmente sin decisión ni registro del tratamiento
\end{itemize}
\end{tcolorbox}

\paragraph{Recomendación de auditoría}
\textbf{Mantener como referencia del proceso de limpieza y deduplicación. No migrar: la limpieza moderna está en src/ingesta/minciencias.py y src/analisis/calidad.py; si se necesita la base persona-nivel, reimplementar con política de año explícita, desempate determinista y registro de decisiones de calidad.}

\medskip
\hrule

\subsection{Otros archivos}
\noindent Archivos: 1.

\subsubsection{src/\_\_init\_\_.py}
\textbf{Rol:} ingesta · \textbf{Líneas:} 2

\paragraph{Descripción funcional}
Marcador de paquete Python para src/. Convierte la carpeta src/ en un paquete importable (los módulos se importan como 'from ingesta import ...' vía sys.path en los scripts).

\paragraph{Análisis de la lógica}
Sin lógica: archivo vacío/de marcador (2 líneas).

\paragraph{Fortalezas}
\begin{itemize}
  \item Correcto por convención: permite imports del paquete.
\end{itemize}

\paragraph{Debilidades}
\begin{itemize}
  \item Ninguna relevante; es boilerplate estándar.
\end{itemize}

\paragraph{Ejecutabilidad}
N/A (marcador de paquete).

\paragraph{¿Produce lo esperado?}
N/A.

\paragraph{Recomendación de auditoría}
\textbf{mantener --- estándar de Python.}

\medskip
\hrule

